% Copyright 2025-2026 by Jonatan Perren
%
% This file may be distributed and/or modified
%
% 1. under the LaTeX Project Public License and/or
% 2. under the GNU Free Documentation License.


\documentclass[a4paper,doc2]{ltxdoc}



\input{../pgfgo.cfg}
% Copyright 2019 by Till Tantau
%
% This file may be distributed and/or modified
%
% 1. under the LaTeX Project Public License and/or
% 2. under the GNU Free Documentation License.
%
% See the file doc/generic/pgf/licenses/LICENSE for more details.


% pgf version is defined in \pgfversion in file
% generic/pgf/utilities/pgfrcs.code.tex

\def\xcolorversion{2.00}

\usepackage[version=latest]{pgf}

\usepackage{xkeyval,calc,listings,tikz,fp}

\usepackage{imakeidx}
\makeindex
\usepackage{hyperref}
\hypersetup{%
    colorlinks=false, % use true to enable colors below:
    linkcolor=blue,%red,
    filecolor=blue,%magenta,
    urlcolor=blue,%cyan,
    citecolor=blue,
    pdfborder=0 0 0,
}

% We need lots of libraries...
\usetikzlibrary{
  3d,
  arrows,
  arrows.spaced,
  arrows.meta,
  bending,
  babel,
  calc,
  fit,
  patterns,
  patterns.meta,
  plotmarks,
  shapes.geometric,
  shapes.misc,
  shapes.symbols,
  shapes.arrows,
  shapes.callouts,
  shapes.multipart,
  shapes.gates.logic.US,
  shapes.gates.logic.IEC,
  circuits.logic.US,
  circuits.logic.IEC,
  circuits.logic.CDH,
  circuits.ee.IEC,
  datavisualization,
  datavisualization.polar,
  datavisualization.formats.functions,
  er,
  automata,
  backgrounds,
  chains,
  topaths,
  trees,
  petri,
  mindmap,
  matrix,
  calendar,
  folding,
  fadings,
  shadings,
  spy,
  through,
  turtle,
  positioning,
  scopes,
  decorations.fractals,
  decorations.shapes,
  decorations.text,
  decorations.pathmorphing,
  decorations.pathreplacing,
  decorations.footprints,
  decorations.markings,
  shadows,
  lindenmayersystems,
  intersections,
  fixedpointarithmetic,
  fpu,
  svg.path,
  external,
  graphs,
  graphs.standard,
  quotes,
  math,
  angles,
  views,
  animations,
  rdf,
  perspective,
}

\usepackage{ifluatex}
\newif\ifgdccodebasic
\newif\ifgdccodeogdf

\ifluatex

  \usetikzlibrary{graphdrawing}
  \usegdlibrary{trees,circular,layered,examples,force,phylogenetics,routing}

  % Test whether C code is available:
  \directlua{
    if pcall(require,'pgf_gd_examples_c_SimpleDemoC') then
      tex.print('\string\\gdccodebasictrue')
    end
  }

  \ifgdccodebasic
    \usegdlibrary{pgf_gd_examples_c_SimpleDemoC,pgf_gd_examples_c_SimpleDemoCPlusPlus}

    % Test whether OGDF code is available:
    \directlua{
      if pcall(require,'pgf_gd_ogdf_c_SimpleDemoOGDF') then
        tex.print('\string\\gdccodeogdftrue')
      end
    }

    \ifgdccodeogdf
      \usegdlibrary{pgf_gd_ogdf_c_SimpleDemoOGDF,ogdf}
    \fi

  \fi

\fi

\def\LuaTeX{Lua\TeX}%


\newif\ifpgfmanualexternalize
\pgfmanualexternalizefalse
\ifpgfmanualexternalize
  \tikzexternalize[
    mode=list only,export=true,% simply skips EVERY picture -> good for debugging the text.
  ]{pgfmanual}

  \tikzifexternalizing{%
    \pgfkeys{/pdflinks/codeexample links=false}%
  }{}%
\fi


\usepackage[a4paper,left=2.25cm,right=2.25cm,top=2.5cm,bottom=2.5cm,nohead]{geometry}
\usepackage{amsmath,amssymb}
\usepackage{xxcolor}
\usepackage{pifont}

\ifluatex
%  \usepackage[no-math]{fontspec}
%  \usepackage{luatextra}

  % TT: I have commented luatextra since it loads fontspec, which
  % currently breaks "$\mathrm{\Omega}$" (nothing is
  % shown). Also, fontspec loads so much stuff, I'm not really sure
  % any of it is needed...

  % \filedescription is defined in expl3, required by fontspec,
  % required by luatextra. Needs to be \relaxed since
  % pgfmanual-en-macros.tex defines an environment named filedescription
  \let\filedescription\relax
  \usepackage[utf8]{luainputenc}
  % dvisvgm does not support OpenType fonts so we have to bite the bullet and
  % use T1 where quotes are weird.  In LuaTeX \outputmode=0 means DVI.
  \ifnum\outputmode=0
    \usepackage[T1]{fontenc}
  \fi
\else
  \usepackage[utf8]{inputenc}
  \usepackage[T1]{fontenc}
\fi

% Promote `Missing character` reports to full errors
% require texlive 2021 or above
\tracinglostchars=3

\graphicspath{{../../img/}}
\input{pgfmanual-en-macros}

\usepackage{todonotes}
\newcommand\todosp[2][]{%        % Stefan Pinnow
    \todo[
        disable,
        color=blue!40,
        #1]{#2}
}

\makeatletter
\renewcommand*\l@section[2]{%
  \ifnum \c@tocdepth >\z@
    \addpenalty\@secpenalty
    \addvspace{1.0em \@plus\p@}%
    \setlength\@tempdima{2.5em}%
    \begingroup
      \parindent \z@ \rightskip \@pnumwidth
      \parfillskip -\@pnumwidth
      \leavevmode \bfseries
      \advance\leftskip\@tempdima
      \hskip -\leftskip
      #1\nobreak\hfil \nobreak\hb@xt@\@pnumwidth{\hss #2}\par
    \endgroup
  \fi}
\renewcommand*\l@subsection{\@dottedtocline{2}{2.5em}{3.3em}}
\renewcommand*\l@subsubsection{\@dottedtocline{3}{5.8em}{4.2em}}
\def\@pnumwidth{2.2em}
\makeatother

%\includeonly{pgfmanual-en-library-profiler}

% Global styles:
\tikzset{
  every plot/.style={prefix=plots/pgf-},
  shape example/.style={
    color=black!30,
    draw,
    fill=yellow!30,
    line width=.5cm,
    inner xsep=2.5cm,
    inner ysep=0.5cm}
}

\index{Options for graphics|see{Graphic options and styles}}
\index{Styles for graphics|see{Graphic options and styles}}
\index{Options for packages|see{Package options}}
\index{Handlers for keys|see{Key handlers}}
\index{File|see{Packages and files}}
\index{Layout|see{Page layout}}
\index{Node|see{Predefined node}}
\index{Data formats|see{Formats}}


% Animations:

\tikzset{make snapshot if necessary}

% Detect changed labels
% by David Carlisle https://tex.stackexchange.com/a/169245
\makeatletter
\def\@testdef#1#2#3{%
  \def\reserved@a{#3}%
  \expandafter\ifx\csname #1@#2\endcsname\reserved@a\else
    \typeout{^^JLabel `#2' changed:^^J\meaning\reserved@a^^J\expandafter\meaning\csname #1@#2\endcsname^^J}%
    \@tempswatrue
  \fi}
\makeatother

%% -----------------------------------------------------------------------------
%\includeonly{
%%    pgfmanual-en-introduction,
%%    pgfmanual-en-tutorial,
%%    pgfmanual-en-tutorial-nodes,
%%    pgfmanual-en-tutorial-Euclid,
%%    pgfmanual-en-tutorial-chains,
%%    pgfmanual-en-tutorial-map,
%%    pgfmanual-en-guidelines,
%%    pgfmanual-en-installation,
%%    pgfmanual-en-license,
%%    pgfmanual-en-drivers,
%%    pgfmanual-en-tikz-design,
%%    pgfmanual-en-tikz-scopes,
%%    pgfmanual-en-tikz-coordinates,
%%    pgfmanual-en-tikz-paths,
%%    pgfmanual-en-tikz-actions,
%%    pgfmanual-en-tikz-arrows,
%%    pgfmanual-en-tikz-shapes,
%%    pgfmanual-en-tikz-pics,
%%    pgfmanual-en-tikz-graphs,
%%    pgfmanual-en-tikz-matrices,
%%    pgfmanual-en-tikz-trees,
%%    pgfmanual-en-tikz-plots,
%%    pgfmanual-en-tikz-transparency,
%%    pgfmanual-en-tikz-decorations,
%%    pgfmanual-en-tikz-transformations,
%%    pgfmanual-en-tikz-animations,
%%    pgfmanual-en-gd-overview,
%%    pgfmanual-en-gd-usage-tikz,
%%    pgfmanual-en-gd-usage-pgf,
%%    pgfmanual-en-gd-trees,
%%    pgfmanual-en-gd-layered,
%%    pgfmanual-en-gd-force,
%%    pgfmanual-en-gd-circular,
%%    pgfmanual-en-gd-phylogenetics,
%%    pgfmanual-en-gd-edge-routing,
%%%    pgfmanual-en-gd-ogdf,
%%    pgfmanual-en-gd-algorithm-layer,
%%    pgfmanual-en-gd-algorithms-in-c,
%%    pgfmanual-en-gd-display-layer,
%%    pgfmanual-en-gd-binding-layer,
%%    pgfmanual-en-library-angles,
%%    pgfmanual-en-library-arrows,
%%    pgfmanual-en-library-automata,
%%    pgfmanual-en-library-babel,
%%    pgfmanual-en-library-backgrounds,
%%    pgfmanual-en-library-calc,
%%    pgfmanual-en-library-calendar,
%%    pgfmanual-en-library-chains,
%%    pgfmanual-en-library-circuits,
%%    pgfmanual-en-library-decorations,
%%    pgfmanual-en-library-er,
%%    pgfmanual-en-library-external,
%%    pgfmanual-en-library-fadings,
%%    pgfmanual-en-library-fit,
%%    pgfmanual-en-library-fixedpoint,
%%    pgfmanual-en-library-fpu,
%%    pgfmanual-en-library-lsystems,
%%    pgfmanual-en-library-math,
%%    pgfmanual-en-library-matrices,
%%    pgfmanual-en-library-mindmaps,
%%    pgfmanual-en-library-folding,
%%    pgfmanual-en-library-patterns,
%%    pgfmanual-en-library-petri,
%%    pgfmanual-en-library-plot-handlers,
%%    pgfmanual-en-library-plot-marks,
%%    pgfmanual-en-library-profiler,
%%    pgfmanual-en-library-rdf,
%%    pgfmanual-en-library-shadings,
%%    pgfmanual-en-library-shadows,
%%    pgfmanual-en-library-shapes,
%%    pgfmanual-en-library-spy,
%%    pgfmanual-en-library-svg-path,
%%    pgfmanual-en-library-edges,
%%    pgfmanual-en-library-through,
%%    pgfmanual-en-library-trees,
%%    pgfmanual-en-library-turtle,
%%    pgfmanual-en-library-views,
%%    pgfmanual-en-dv-introduction,
%%    pgfmanual-en-dv-main,
%%    pgfmanual-en-dv-formats,
%%    pgfmanual-en-dv-axes,
%%    pgfmanual-en-dv-visualizers,
%%    pgfmanual-en-dv-stylesheets,
%%    pgfmanual-en-dv-polar,
%%    pgfmanual-en-dv-backend,
%%    pgfmanual-en-pgfkeys,
%%    pgfmanual-en-pgffor,
%%    pgfmanual-en-pgfcalendar,
%%    pgfmanual-en-pages,
%%    pgfmanual-en-xxcolor,
%%    pgfmanual-en-module-parser,
%%    pgfmanual-en-math-design,
%%    pgfmanual-en-math-parsing,
%%    pgfmanual-en-math-commands,
%%    pgfmanual-en-math-algorithms,
%%    pgfmanual-en-math-numberprinting,
%%    pgfmanual-en-oo,
%%    pgfmanual-en-base-design,
%%    pgfmanual-en-base-scopes,
%%    pgfmanual-en-base-points,
%%    pgfmanual-en-base-paths,
%%    pgfmanual-en-base-decorations,
%%    pgfmanual-en-base-actions,
%%    pgfmanual-en-base-arrows,
%%    pgfmanual-en-base-nodes,
%%    pgfmanual-en-base-matrices,
%%    pgfmanual-en-base-transformations,
%%    pgfmanual-en-base-patterns,
%%    pgfmanual-en-base-images,
%%    pgfmanual-en-base-external,
%%    pgfmanual-en-base-plots,
%%    pgfmanual-en-base-layers,
%%    pgfmanual-en-base-shadings,
%%    pgfmanual-en-base-transparency,
%%    pgfmanual-en-base-animations,
%%    pgfmanual-en-base-internalregisters,
%%    pgfmanual-en-base-quick,
%%    pgfmanual-en-pgfsys-overview,
%%    pgfmanual-en-pgfsys-commands,
%%    pgfmanual-en-pgfsys-paths,
%%    pgfmanual-en-pgfsys-protocol,
%%    pgfmanual-en-pgfsys-animations,
%}
%% -----------------------------------------------------------------------------


%%% Local Variables:
%%% mode: latex
%%% TeX-master: "pgfmanual"
%%% End:


\usepackage{animate}
\usepackage{../../latex/pgf-go}
%LOS ARCHIVOS SGF DE LOS EJEMPLOS, POR SU NOMBRE (LOS TUTORIALES Y EL CAPITULO SGF LOS PRESENTAN).
\pgfgodefinesgf{game}{game.sgf}
\pgfgodefinesgf{lee}{../examples/lee.sgf}
\pgfgodefinesgf{variations}{variations.sgf}
\pgfgodefinesgf{problem}{../examples/problem.sgf}
\pgfgodefinesgf{josekis}{../examples/joseki.sgf}
%LOS KANJI DEL TEXTO (VER label/style).
\font\pgfgokanji = {file:ipaexm.ttf} at 9pt
%LA MISMA FUENTE COMO FAMILIA, PARA format = kanji (VER mark/label/format).
\DeclareFontFamily{TU}{ipaex}{}
\DeclareFontShape{TU}{ipaex}{m}{n}{<->file:ipaexm.ttf}{}
\DeclareFontShape{TU}{ipaex}{bx}{n}{<->ssub*ipaex/m/n}{}
\def\pgfgoname{\textsc{pgfgo}}
\def\karu{K\={a}ru}

\usepackage[english]{babel}

\usepackage{float}
\usetikzlibrary{positioning, fit}
\usepackage{caption}



\definecolor{pgfgo}{RGB}{220,179,92}
\colorlet{tablero2}{red!10!blue!10}
\colorlet{tablero3}{blue!10!green!10}
\definecolor{tablero4}{HTML}{9B8465}

\def\pgfgoversion{\leavevmode\hbox{0.2026.10.30}}



%LAS LINEAS CON CODIGO EN EL TEXTO SE ESTIRAN EN LUGAR DE SALIRSE DEL MARGEN.
\setlength{\emergencystretch}{3em}

\begin{document}
	\title{\Huge\bfseries pgf-go\\[1em]
		\normalsize

		\begin{center}
			\pgfgoset{profile = pgfgosmooth}
			\begin{goban}[partial = B2:E3]
				\stones[black]{C2,D2,E3}
				\stones[white] {B2,C3,D3}
				\mark[w, custom = {36}]{B2}
				\mark[custom = {a}] {E2}

				\marks[circle, line = {color = red!80!black}] {C3,D3}

			\end{goban}
		\end{center}

		\large Manual for version~\pgfgoversion}
		\author{Jonat\'an Perren Hug\\
		\texttt{jonatan.perren@gmail.com}
		}

	\maketitle



	\label{table-of-contents}
	\tableofcontents

	\clearpage

	\part{Introduction}
	\section{Introduction}

\subsection{Motivation}

The following package offers a collection of macros designed to represent Go boards and positions, fully implemented in \pgfname. This project, although still under development, incorporates innovative features aimed at simplifying and streamlining the creation of Go-related diagrams:
\begin{itemize}
	\item Profile management for boards and players.
	\item Loading multiple coordinates simultaneously using a syntax that some might find familiar.
\end{itemize}

The development of \pgfgoname stems from frustration with not finding packages that met my needs. The packages I explored allow for the representation of black-and-white boards, suitable for printing, but impractical when combined with technologies like |beamer| or in documents where a color presentation is desired.

Initially, \pgfgoname was intended to be limited to these tasks; however, I now envision it more as an interface for managing styles and states, avoiding excessive use of direct formatting. This will be reflected primarily in the structure of the \emph{keys} tree.

Regarding the name, although the macros are implemented in \pgfname and L3, the board is built on |\tikzpicture| so that users familiar with this environment can take advantage of its macros.

\subsection{What’s New}

\begin{itemize}
	\item \version:
	\begin{itemize}
			\item The \emph{keys} |inner| and |outer| have been added for the \emph{rounded corner} of the board.
			\item Three new coordinate operators have been added: |*| for all intersections, |L| for vertical lines, and |N| for horizontal lines.
	\end{itemize}
	\item 0.2025.03.07: First release.
\end{itemize}

	\part{Tutorials}\label{part:tutoriales}
	\def\karu{K\={a}ru}

\section{The Boards of \karu}

\subsection{Board and Stones}

\karu\ is writing an article to share with his students. He has noticed that diagramming \emph{software} does not offer satisfactory formats or results. Therefore, he decides to try \pgfgoname\ and starts by representing a game like the one shown in Figure~\ref{fig:first-example}.

\begin{figure}[H]
	\centering
	\pgfgouseprofile{pgfgo}
	\begin{goban}[use i = false, label = {at = {north, west}}]
		\stones[black]{B2,E2,J2,C3,D3,G3,H3,C5,D5,E5,F5,G5,E6,G6}
		\stones[white]{D8,F8,E7,H7,C6,D6,F6,H6,B5,J5,B4,E4,J4,J3,H2}
	\end{goban}
	\caption{}
	\label{fig:first-example}
\end{figure}

\karu\ must start with the board. In \pgfgoname, there is currently a single environment where elements can be placed: |goban|. This environment controls by default the size and color of the board, the position and type of labels, among other properties. It also allows partial board representations.

\begin{codeexample}[width=4.4cm]
	\begin{goban}
	\end{goban}
\end{codeexample}

\karu\ is satisfied. The environment provides a default $9\times 9$ board, which is exactly what he needs, but he must add side labels. The |goban| environment does not require mandatory parameters, although it accepts an optional one: a list of |keys| or |key=values| for configuration. \karu\ consults the section on |goban| (p.~\pageref{sec:board}) and decides to add labels to the environment. He knows that the \emph{keys} controlling the objects in the goban are divided into three groups: |grid|, for the central grid; |background|, for the border and board fill; and |label|, for coordinate labels. The latter is the one that interests him. This \emph{key}, among other options, can take the \emph{key} |at| to define the label position with values such as |above|, |west|, |below|, |right|, |all|, or |none|. Cardinal points can also be used. Since he only wants labels above and to the left of the goban, he will use the \emph{keys} |above| and |left|.

\begin{codeexample}[width=4.7cm]
	\begin{goban}[label = {at = {above, left}}]
	\end{goban}
\end{codeexample}

\karu\ feels that everything is going well, but he notices that his board includes an |I| instead of omitting it and going directly to |J|. By default, in \pgfgoname, "I" exists and is used both as a coordinate and in the side labels, although it is possible to remove it. There are two ways to do this: setting the global \emph{key} |use i| to |false| (|\pgfgoset{use i = false}|), which would prevent \karu\ from worrying about this throughout the rest of the article, or applying it directly to the board with the same \emph{key} and value.

\begin{codeexample}[width=4.8cm]
	\pgfgoset{use i = false}
	
	\begin{goban}[label = {at = {above, left}}]
	\end{goban}
\end{codeexample}

Now he is ready to place the stones. For this, \karu\ has several macros. Let’s start with the most basic one: |\stone[options]{intersection}|. This command accepts two parameters: the first, optional, controls the appearance and additional elements of the stones (such as marks); the second, mandatory, is the intersection. \karu\ does not need to convert intersections to numbers since he can specify them alphanumerically. By default, the stones are neutral, but he can choose the first player by passing the \emph{key} |black| or |b|.

\begin{codeexample}[width=4.8cm]
	\begin{goban}[label = {at = {above, left}}, use i = false]
		\stone[black] {D8}
		\stone[black] {F8}
		\stone[b] {E7}
		\stone[b] {H7}
	\end{goban}
\end{codeexample}


	\part{Reference}\label{part:referencia}
	\section{Settings}
\label{sec:configuracion}

|\pgfgoset{...}| sets options for everything that follows: the look of the board, of the stones, of the moves, etc. It is equivalent to |\keys_set:nn{pgfgo}{...}|. It is also a door to direct formatting.\footnote{Use with discretion\dots\ not like footnotes.}

\begin{command}{\pgfgoset\marg{options}}
\end{command}

\begin{key}{/pgfgo/goban = \meta{goban options}}
	Options of the board (p.~\pageref{sec:tablero}): the size, the grid, the background, the coordinates, etc.
\end{key}

\begin{key}{/pgfgo/stone = \meta{stone options}}
	Options of the stones of the three players (p.~\pageref{sec:piedras}), the same as |players = {all = {stone = {...}}}|.
\end{key}

\begin{key}{/pgfgo/move = \meta{move options}}
	Options of the moves, such as the format of their numbers. Except for |numbers| (p.~\pageref{subsec:movimientos}), which holds for every move, it is the same as |players = {all = {move = {...}}}|.
\end{key}

\begin{key}{/pgfgo/mark = \meta{mark options}}
	Options of the marks: the default shape of |\marks| (|\pgfgoset{mark = {triangle}}|) and the style of the marks of the three players (|line|, |fill|, |label|, |sequence| and |scale|): |mark = {scale = 2}| is the same as |players = {all = {mark = {scale = 2}}}|.
\end{key}

\begin{key}{/pgfgo/sgf animation = \meta{sgf animation options}}
	Options of every |\sgfanimation| of the document (p.~\pageref{subsec:sgfanimacion}).
\end{key}

\begin{key}{/pgfgo/sgf diagrams = \meta{sgf diagrams options}}
	Options of every |\sgfdiagrams| of the document (p.~\pageref{subsec:sgfdiagramas}).
\end{key}

\begin{key}{/pgfgo/sgf problem = \meta{sgf problem options}}
	Options of every |\sgfproblem| of the document (p.~\pageref{subsec:sgfproblemas}).
\end{key}

\begin{key}{/pgfgo/sgf comment = \meta{sgf comment options}}
	Options of every |\sgfcomment| of the document (p.~\pageref{subsec:sgfdatos}).
\end{key}

\begin{key}{/pgfgo/sgf info = \meta{sgf info options}}
	Options of every |\sgfinfo| of the document (p.~\pageref{subsec:sgfdatos}).
\end{key}

\begin{key}{/pgfgo/sgf annotations = \meta{texts}}
	The texts of the annotations of SGF files (p.~\pageref{subsec:sgfdatos}).
\end{key}

\begin{key}{/pgfgo/arrow = \meta{arrow options}}
	Options of every arrow of the document (p.~\pageref{subsec:flechas}).
\end{key}

\begin{key}{/pgfgo/line = \meta{line options}}
	Options of every line of the document (p.~\pageref{subsec:flechas}).
\end{key}

\begin{key}{/pgfgo/territory = \meta{territory options}}
	Options of every |\territory| of the document, such as |draw = fill| (p.~\pageref{subsec:territorio}).
\end{key}

\begin{key}{/pgfgo/score = \meta{score options}}
	Options of every |\pgfgoscore| of the document, such as |komi = 6.5| (p.~\pageref{subsec:territorio}).
\end{key}

\begin{key}{/pgfgo/players = \meta{players options}}
	Options of each player: black, white and neutral (p.~\pageref{sec:jugadores}).
\end{key}

\begin{codeexample}[]
	\pgfgoset{
		goban = {
			background = {
				fill = {color = pgfgo},
				line = false},
			grid = {line = {color = white}}},
		players = {
			black = {
				stone = {
					line = {color = magenta},
					fill = {color = blue!50!black}}},
			white = {
				stone = {line = {color = blue}},
				mark = {line = {color = orange}}
			}
		}
	}

	\begin{goban}[3]
		\stone[black] {A1}
		\stone[white, mark = {triangle}] {A2}
	\end{goban}

	\begin{goban}[4]
		\stone[black] {A2}
		\stones[white, mark = {circle}] {B2,A3}
	\end{goban}
\end{codeexample}

The |goban| environments that follow inherit these options, and each one can change them. The options of a |goban| accept, besides those of the board, those of |\pgfgoset|, which hold only for that board: |\begin{goban}[9, stone = {radius = 4pt}, mark = {line = {color = red}}]|.

\label{subsec:prender}Many \emph{keys} turn something on or off: the grid, the background, the coordinates, the border and the fill of the stones, their shadows and their shines, etc. They all follow the same rule:
\begin{itemize}
	\item the \emph{key} alone, or with |true|, turns it on: |shine|, |shine = true|;
	\item with |false| it turns it off: |shine = false|;
	\item giving options does not turn it on: |shine = {angle = 45}| changes the angle, but the shine stays as it was. To turn it on and change it at once, |shine = {true, angle = 45}|.
\end{itemize}
The exceptions are the \emph{keys} whose value is the content itself, such as |partial| (the part of the board) and |sequence| (where the numbering starts).
\subsection{The compute kernel}
\label{subsec:kernel}

\pgfgoname\ separates the computations from the drawing. The computations (reading the coordinates, knowing which part of the board is visible, numbering the moves and keeping the turn, remembering the boards, etc.) are done by a \emph{kernel}. The drawing, the \emph{keys} and the profiles are always the same. There are two kernels:
\begin{itemize}
	\item |lua|: written in Lua. It needs Lua\LaTeX, and it is the default with Lua\LaTeX.
	\item |expl3|: written in L3. It works with any engine, and it is the default with the others.
\end{itemize}

Both produce exactly the same diagrams\footnote{Or so the tests say.}. The Lua kernel also knows three things the other does not: captures (p.~\pageref{subsec:capturas}), finding the territory with |\territory{auto}| (p.~\pageref{subsec:territorio}) and reading games in SGF (p.~\pageref{subsec:sgf}). This manual is compiled with the Lua kernel.

\begin{key}{/pgfgo/package/kernel = \meta{expl3 or lua} (initially lua with Lua\LaTeX, expl3 with the other engines)}
	Package option that chooses the kernel. It is only needed to use |expl3| with Lua\LaTeX. It can only be given when \pgfgoname\ is loaded; it cannot be changed with |\pgfgoset|. If |kernel = lua| is asked for with pdf\LaTeX, the package warns and uses |expl3|. With |expl3|, |captures|, |\territory{auto}| and |\sgf| give an error.

	\begin{codeexample}[code only]
		\usepackage[kernel = expl3]{pgf-go}
	\end{codeexample}
\end{key}

\subsection{Languages}
\label{subsec:idiomas}

Some texts are written by the package itself: the captions of |\sgfdiagrams| and |\sgfproblem|, the result of the game, the annotations of the comments (|/pgfgo/sgf annotations|) and the word of the ``N at M'' notes. They are translated into English (|en|), Spanish (|es|), French (|fr|), German (|de|), Portuguese (|pt|), Italian (|it|), Esperanto (|eo|), Japanese (|ja|), Chinese (|zh|) and Korean (|ko|). Japanese, Chinese and Korean need a font with those characters.

By default, the texts are in the language set by babel or polyglossia, and they follow it if it changes inside the document (|\selectlanguage|, |otherlanguage|). Without babel or polyglossia, or in a language without a translation, they are in English. Each text can be changed with its \emph{key} (for example, |/pgfgo/sgf diagrams/title|); what is set that way holds in any language.

\begin{key}{/pgfgo/lang = \meta{language} (initially auto)}
	Language of the texts: |auto| (that of babel or polyglossia), the code of a language (|fr|) or the name babel or polyglossia give it (|french|). If the language is not known, the package warns and uses English.

	\begin{codeexample}[code only]
		\pgfgoset{lang = fr}
	\end{codeexample}
\end{key}

\begin{key}{/pgfgo/package/lang = \meta{language} (initially auto)}
	The same, as a package option: |\usepackage[lang = fr]{pgf-go}|.
\end{key}

\begin{command}{\pgfgotext\marg{text}}
	Writes one of the texts in the current language. The texts, with their value in English and in Spanish, are:

	\begin{center}
		\small
		\begin{tabular}{lll}
			|tesuji| & Tesuji & Tesuji \\
			|very good tesuji| & Very good tesuji & Muy buen tesuji \\
			|bad move| & Bad move & Mala jugada \\
			|very bad move| & Very bad move & Muy mala jugada \\
			|doubtful| & Doubtful move & Jugada dudosa \\
			|interesting| & Interesting move & Jugada interesante \\
			|good for black| & Good for Black & Bueno para Negras \\
			|very good for black| & Very good for Black & Muy bueno para Negras \\
			|good for white| & Good for White & Bueno para Blancas \\
			|very good for white| & Very good for White & Muy bueno para Blancas \\
			|even| & Even position & Posición pareja \\
			|unclear| & Unclear position & Posición incierta \\
			|diagram| & Diagram & Diagrama \\
			|variation| & Variation & Variante \\
			|black to play| & Black to play & Juegan Negras \\
			|white to play| & White to play & Juegan Blancas \\
			|at| & at & en \\
			|result| & Result & Resultado \\
			|black| & Black & Negras \\
			|white| & White & Blancas \\
			|win| & \meta{winner} wins & Ganan \meta{winner} \\
			|win by points| & \meta{winner} wins by \meta{points} points & Ganan \meta{winner} por \meta{points} puntos \\
			|win by resignation| & \meta{winner} wins by resignation & Ganan \meta{winner} por abandono \\
			|win on time| & \meta{winner} wins on time & Ganan \meta{winner} por tiempo \\
			|win by forfeit| & \meta{winner} wins by forfeit & Ganan \meta{winner} por incomparecencia \\
			|draw| & Draw (jigo) & Empate (jigo) \\
			|void| & No result & Sin resultado \\
			|unknown result| & Unknown result & Resultado desconocido \\
		\end{tabular}
	\end{center}

	\begin{codeexample}[]
		\pgfgotext{black to play}. \pgfgotext{tesuji}.

		\pgfgoset{lang = fr}
		\pgfgotext{black to play}.
	\end{codeexample}
\end{command}

\begin{command}{\pgfgodefinelanguage\marg{language}\marg{texts}}
	Adds a language, or changes texts of one that already exists. The missing texts are taken from English. In the texts of the result, |\pgfgowinner| is the name of the winner and |\pgfgopoints| the difference (p.~\pageref{subsec:sgfdiagramas}).

	\begin{codeexample}[code only]
		\pgfgodefinelanguage{es}{black = Negro, white = Blanco,
			win by points = Gana \pgfgowinner{} por \pgfgopoints{} puntos}
	\end{codeexample}
\end{command}

\begin{command}{\pgfgolanguagealias\marg{name}\marg{language}}
	Another name for a language, for example the one babel or polyglossia give it, if \pgfgoname\ does not know it.

	\begin{codeexample}[code only]
		\pgfgolanguagealias{castellano}{es}
	\end{codeexample}
\end{command}

\subsection{One board per page}
\label{subsec:paginas}

With the option |pages|, the document comes out with one |goban| per page, each one the size of its board, and nothing else: the text, the captions and the notes do not come out. Neither do the views in the text (|\pgfgoview|) and the animations. So the boards of a document can be used in another program (a word processor, an ebook, a web page) without copying their code.

At the end, the package also writes |<document>-pages.lua|, which turns each page into a file in the folder |<document>-pages|. The file is named like the |goban| (|name|; only letters, digits, |.|, |_| and |-|), or |diagram-| and the number of the page if it has no name. It is run with |texlua|:

\begin{codeexample}[code only]
texlua book-pages.lua          % SVG, with pdftocairo or dvisvgm
texlua book-pages.lua png 600  % PNG at 600 dpi (300 if not given), with pdftocairo or pdftoppm
texlua book-pages.lua pdf      % one PDF per board, with pdftocairo or pdfseparate
\end{codeexample}

\begin{key}{/pgfgo/package/pages = \meta{boolean} (default true, initially false)}
	Package option: one |goban| per page. It needs Lua\LaTeX{} and |kernel = lua|; otherwise, the package warns and the document comes out as always.

	\begin{codeexample}[code only]
		\usepackage[pages]{pgf-go}
	\end{codeexample}
\end{key}

	\section{Coordinate syntax}
\label{subsec:sintaxis-de-coordenadas}

Something that sets \pgfgoname\ apart from other packages is the possibility of entering coordinates with several syntaxes:
\begin{itemize}
	\item Multiple entries, separated by commas.
	\item Removing stones by adding the subtractive group after a <<;>>.
\end{itemize}

These forms only work in the plural commands (|\stones|, |\marks|, |\prisoners|); |\stone|, |\mark| and |\prisoner| take a single intersection. |\moves| accepts all of them except the subtractive groups, and warns if they have no fixed order (p.~\pageref{subsec:orden-coordenadas}). The writing has some restrictions. For example, to refer to an intersection, the alphabetic coordinate must be written in capitals and without spaces. So the intersection B5 must be entered as |B5|, and not as |B 5|, |B  5| or |b5|, since spaces between letter and number are not valid in this context. On boards of more than 25 columns, the columns after |Z| have two letters: |AA|, |AB|, \dots\ (|AA5|, |AA:AC 3|, |partial = AZ5:BG8|). The restrictions will be explained in each case.


When \pgfgoname\ has to read the argument of the coordinates, it may find something like this
\begin{center}
	|\stones{|{\color{red}|<item>, <item>, |\dots;} {\color{blue} |<item>, <item>, |\dots; }{\color{red} \dots}|}|
\end{center}

The first thing it will try to do is split the argument at the |;| and build two groups: the additive group and the subtractive group. The additive group adds stones with our syntax. The subtractive group removes them. These groups alternate automatically with each |;|. This is seen in more detail in the corresponding subsection. If no |;| is given, \pgfgoname\ will add all the stones to the additive group. Several stones or several groups can be added, separated by |,|. Altogether, this provides an interface to create complex groups with few \emph{tokens};

\begin{codeexample}[width=0pt]
	\pgfgoset{profile = pgfgosmooth}

	\begin{goban}[17, partial = A1:R5, label = {at = {north, west}}]
		\stones[b]{*;B,D,H,M,R2,Q,F2:4,K2rL5,O[1:2,4];B3}
	\end{goban}
\end{codeexample}


\subsection{Groups of stones}

\def\sssseccion#1{\vskip0.5em\noindent\textbf{#1}\par}

\sssseccion{All the stones of the board (\emph{Why not?})}

Adds to the corresponding group every \textbf{visible} intersection of the board. That is, on a partial board, only those of the visible region are added. This is worth making clear in case it is used together with |resume|. Its syntax is |*|.

\begin{codeexample}[]
	\pgfgoset{goban = {5, partial = B2:D4}}

	\begin{goban}[label = {at = {south, west}}]
		\stones[b]{*}
	\end{goban}
\end{codeexample}


\sssseccion{Intersection}

Intersections are given by the letter followed by the number, without spaces between the alphabetic and the numeric component: |LN|

\begin{codeexample}[]
	\pgfgoset{goban = {5, partial = B2:D4}}

	\begin{goban}[label = {at = {south, west}}]
		\stones[b]{C2, D3}
		\stone[w]{B4}
		\stone[w]{B 3} % Not recognized because it has a space.
	\end{goban}
\end{codeexample}

In the example, |B 3| is not recognized because of the space, while |B4| is drawn correctly as a white stone.

\sssseccion{Vertical or horizontal lines}
Lines are groups of vertical or horizontal stones implicitly bounded by the size of the board: a whole column or row.

\begin{itemize}
	\item For a column (a vertical line), the syntax is |L|: |D| is column D.
	\item For a row (a horizontal line), the syntax is |N|: |4| is row 4.
\end{itemize}

\begin{codeexample}[width=6cm]
	\pgfgoset{goban = {5, partial = B2:D4}}

	\begin{goban}
		\stones[b]{B2,B3,B4}
		\stones[fill = {color = red}]{D}
	\end{goban}

	\begin{goban}
		\stones[b]{B2,C2,D2}
		\stones[fill = {color = blue}]{4}
	\end{goban}

	\begin{goban}
		%Add the lines C and 3. Then remove C3.
		\stones[fill = {color = orange}]{C, 3; C3} %CAREFUL: C,3 != C3.
	\end{goban}
\end{codeexample}

\sssseccion{Segments}

Segments are groups of vertical or horizontal stones explicitly bounded. Since they go in one direction (only vertical or only horizontal), one of their coordinates stays constant and need not be repeated; only the range of the other is given.

\begin{itemize}
	\item For horizontal segments, the syntax is |L:LN| or |L:L N|.
	\item For vertical segments, the syntax is |LN:N| or |L N:N|.
\end{itemize}

\begin{codeexample}[width=6cm]
	\pgfgoset{goban = {5, partial = B2:D4}}

	\begin{goban}
		\stones[b]{B2,B3,B4}
		\stones[w]{C2,D2}
	\end{goban}

	\begin{goban}
		\stones[b]{B2:4}
		\stones[w]{C:D2}
	\end{goban}

	\begin{goban}
		% Spaces work between different elements.
		\stones[b]{B 2:4}
		\stones[w]{C:D 2}
	\end{goban}
\end{codeexample}

Here, |B2:4| stands for a vertical segment from B2 to B4, while |C:D2| is a horizontal segment from C2 to D2. Segments go from the first end to the second: |H9:2| goes down from H9 to H2 and |E:A2| goes from E2 to A2. They can also be chained with the operators we will see later.

\begin{codeexample}[width=6cm]
	\begin{goban}[label = {at = all}]
		\stones[fill = {color = red}]{A:E5}
		\stones[fill = {color = blue}]{H9:2}
		\stones[fill = {color = orange}]{E:A2}
		\stones[fill = {color = magenta}]{A:D7, A:C9} % Combined.
	\end{goban}
\end{codeexample}

\sssseccion{Rectangle: an artificial consequence}

The segment syntax creates ranges between coordinates to insert stones. I thought it useful to extend it to two ranges, as a first approach to rectangles (although it is not my favorite method, here it is). In this case, the syntax is |L:L N:N| or |L:LN:N|.

\begin{codeexample}[width=6cm]
	\begin{goban}[label = {at = all}]
		% The space is not required.
		\stones[fill = {color = red}]{A:E 7:8}
		\stones[fill = {color = blue}]{I:G 9:2}
		\stones[fill = {color = orange}]{A:B 2:4, F:D 1:3}
	\end{goban}
\end{codeexample}

For example, |A:E 7:8| creates a rectangle that covers the columns from A to E in rows 7 and 8, while |I:G 9:2| goes down from I9 to G2. It works, but the operators we will see later are clearer.


\sssseccion{Cartesian products}
	In this context, the Cartesian product is a more powerful extension of the segments, but at the same time a generalization of all the notations presented so far. The idea behind it is the following: suppose we want to enter a group of stones on column A: |A1:3, A5, A7:8|. So far there was no simpler way to write this. But they all repeat A at the beginning: what if, instead of writing it three times, we write it only once? Something like
	\begin{center}
		|A[1:3, 5, 7:8]|
	\end{center}
	In this case, the interpreter distributes the A on the left to each of the elements on the right.

	The elements of the list on the right need not be made of numbers. Something like |A:D 5:7, A1:2, A3, A:G9| can also be rewritten as
	\begin{center}
		|A[:D 5:7, 1:2, 3, :G9]|
	\end{center}
	Of course, nobody abstracts that much when entering coordinates, so this is a mere whim. The only restriction on the content in brackets is that brackets cannot be nested.

	Just as there is a syntax for letters on the left, there is one for numbers on the right, and it works the same way. The expression |A12, C12, E:G12| can be written as |[A,C,E:G]12|.

	There is a third operator which, yes, is more of a whim than anything else: the unrestricted Cartesian product. I will not dwell on it, because it requires a higher degree of abstraction and really has no human use. This operator is a combination of the previous two: the expression |[A,D:F][2,5,7:9]| is read as
	\begin{center}
		|A2, A5, A7:9, D:F2, D:F5, D:F 7:9|.
	\end{center}

	\begin{codeexample}[width=4cm]
		\begin{goban}
			\stones[b]{[A, D:F][2, 5, 7:9]}
			\marks[b]{A2, A5, A7:9, D:F2, D:F5, D:F 7:9} %<- what the interpreter reads.
		\end{goban}
	\end{codeexample}

	The Cartesian product syntax can then be used in three different ways:
\begin{itemize}
	\item |L[itm1, itm2,...]|.
	\item |[itm1, itm2,...]N|.
	\item |[itm1,itm2,...][itm3,itm4,...]|.
\end{itemize}
where brackets cannot be nested.

This can also be combined with some of the operators we will see next, to create nodes in the paths.

So far we have seen forms that cannot be combined with others in one element and that, therefore, \textbf{must} be separated by commas. Now we will see operators that join two intersections with a path and that can be chained.

\sssseccion{Horizontal--vertical and vertical--horizontal}

Given two intersections not aligned vertically or horizontally, these commands generate the intersections that connect the first with the second along a ``horizontal--vertical'' or ``vertical--horizontal'' path. This may remind you of the \verb+-|+ operator of TikZ, from which I borrowed the idea. The syntax is \verb+LN -| LN+ (for the first case) or \verb+LN |- LN+ (for the second), with or without spaces. These commands, together with the following ones, can be chained to make stairs or paths of the Manhattan metric.

\begin{codeexample}[width=6cm]
	\begin{goban}[label = {at = all}]
		\stones[fill = {color = red}]{A1 -| D3}
		\stones[fill = {color = blue}]{D4 -| A2}
		\stones[fill = {color = orange}] {A9 -| C7 -| G5} % Chained.
		\stones[fill = {color = magenta}] {F1 |- I3 -| E4} % Try to guess.
	\end{goban}
\end{codeexample}

In this example, \verb+A1 -| D3+ goes horizontally from A1 to D1 and then vertically to D3.

\sssseccion{Lines in any direction}

For any two intersections, |LN -- LN| draws the line of stones closest to the straight line that joins them: one stone per row (or per column, if the line moves sideways more than up or down), the one closest to the straight line. If both intersections are on the same row or column, it is the straight line between them.

\begin{codeexample}[width=6cm]
	\begin{goban}[label = {at = all}]
		\stones[fill = {color = red}]{A1 -- A5}
		\stones[fill = {color = blue}] {A1 -- D1}
		\stones[fill = {color = magenta}] {B9 -- D3}
		\stones[fill = {color = orange}] {I9 -| F6 -- I6} % Chained.
	\end{goban}
\end{codeexample}

When the straight line passes exactly between two intersections, the stone goes on the side of the nearer end, and in the exact middle, on the lower column (or row). So |A1 -- C5| and |C5 -- A1| give the same stones, and the line is symmetric, like the straight line.

\sssseccion{Stairs}

The stones of a slanted line touch only at the corners at each step, as on a diagonal, so in Go they do not form a group. To make them form one, \verb+LN -|- LN+ and \verb+LN |-| LN+ fill each step with the corner of \verb+-|+ (horizontal first) or of \verb+|-+ (vertical first): the result is a connected staircase, as the drawing of the operator suggests. If both intersections are on the same row or column, they give the same straight line as |--|.

\begin{codeexample}[width=6cm]
	\begin{goban}[label = {at = all}]
		\stones[fill = {color = red}]{A1 -- C7}
		\stones[fill = {color = blue}]{D1 -|- F7}
		\stones[fill = {color = orange}]{G1 |-| I7}
	\end{goban}
\end{codeexample}

In stairs the order of the ends does matter, as in \verb+-|+ and \verb+|-+.

\sssseccion{Filled and empty rectangles}

The previous syntax for rectangles, |L:L N:N|, is not very intuitive and rather a legacy of the bounded segments. It is more direct to give two diametrically opposite vertices: |LN:LN| fills the rectangle (it is the same writing as that of |partial|) and |LN r LN| draws only its outline. Although the spaces around |r| are not required, I recommend them for readability.

\begin{codeexample}[width=6cm]
	\begin{goban}[label = {at = all}]
		\stones[fill = {color = red}]{A1:D3}
		\stones[fill = {color = blue}] {A5 r D8}
		\stones[fill = {color = orange}] {I8 r G4, H6} % It's alive!
	\end{goban}
\end{codeexample}

For example, |A1:D3| fills a rectangle from A1 to D3, while |A5 r D8| draws only the border of a larger rectangle. A rectangle of a single column or a single row is a segment: |A5:A9| is |A5:9|.

\sssseccion{Chaining}

The operators (\verb+-|+, \verb+|-+, |--|, \verb+-|-+, \verb+|-|+ and |r|) can be chained: \verb+A1 -| C3 -- E3+ is \verb+A1 -| C3+ followed by \verb+C3 -- E3+, and C3 appears only once. Between two operators there can be not only an intersection: also a segment (|B5:7|, |A:B3|) or a rectangle (|A1:C3|). Each operator goes from the last point of what comes before to the first of what comes after. So \verb+A2 -| B5:7+ is \verb+A2 -| B5+ followed by |B5:7|, and \verb+B7 -| A:B3+ is \verb+B7 -| A3 -- B3+.

\begin{codeexample}[width=6cm]
	\begin{goban}[label = {at = all}]
		\stones[fill = {color = red}]{A2 -| B5:7}
		\stones[fill = {color = blue}]{I9 -| D:F6}
		\stones[fill = {color = orange}]{F1:3 |- I5}
	\end{goban}
\end{codeexample}

Rectangles can also be chained, to create more complex patterns.

\begin{codeexample}[width=6cm]
	\begin{goban}[label = {at = all}]
		\stones[fill = {color = orange}] {I8 r G6 r I4} % You guess.
		\marks[sequence]{I8,G6,I4}
		\stones[fill = {color = blue!20}] {A1:E3 r A6 |- E9 -- E8} % Wow!
		\marks[sequence]{A1,E3,A6,E9,E8}
	\end{goban}
\end{codeexample}

In this case, |I8 r G6 r I4| chains two outline rectangles, while \verb+A1:E3 r A6 |- E9 -- E8+ combines a filled rectangle, an outline, a ``vertical--horizontal'' path and a line. The forms that go alone (|*|, |L|, |N|, |L:L|, |N:N| and |L:L N:N|) cannot be chained.

\sssseccion{The order of the intersections}
\label{subsec:orden-coordenadas}

The intersections come out in the order in which they are written: segments go from the first end to the second (|N10:7| goes down from N10 to N7), the operators follow their path and brackets are distributed in place (|D4, B[3:1]| is |D4, B3:1|). That is the order numbered by |\moves| and |sequence|.

\begin{codeexample}[width=6cm]
	\begin{goban}[7, label = {at = {south, west}}]
		\stones[w]{A1 -|- F4}
		\marks[sequence]{A1 -|- F4}
	\end{goban}

	\begin{goban}[7, label = {at = {south, west}}]
		\stones[w]{G7:4 -| D2 -- A2}
		\marks[sequence = {format = Alph}]{G7:4 -| D2 -- A2}
	\end{goban}
\end{codeexample}

Some forms have no fixed path: rectangles (|LN:LN|, |L:L N:N| and |LN r LN|), whole columns and rows, and |*|. |\stones| and |\marks| accept them like the others, but |\marks[sequence]| and |\moves| warn (|unordered-coordinate|), because the order in which they go through them is not chosen by the writer.

\subsection{Additive and subtractive groups}

Suppose we want to draw the following diagram:

\begin{figure}[H]
	\centering
	\begin{goban}
		\stones[b]{A1:E6; C6 |- E3, A 1:2;}
	\end{goban}
\end{figure}

This can be done with the operators already seen, splitting the pattern into several parts:

\begin{codeexample}[width=6cm]
	\begin{goban}[label = {at = all}]
		\stones[fill = {color = orange}] {A6:B3, B2:E1, D4:E6}
	\end{goban}
\end{codeexample}

But we can also imagine that we start from a full rectangle from which the marked stones have been removed:

\begin{figure}[H]
	\centering
	\begin{goban}
		\stones[b]{A1:E6}
		\marks[cross, line = {color = white}]{C6 |- E3, A1:2}
	\end{goban}
\end{figure}

That is what additive and subtractive groups are about. An additive group adds intersections to the list, while a subtractive one removes them. By default, the commands always start by processing additive groups, which explains why they draw stones directly. In \pgfgoname, we can alternate between both kinds with |;|.

\begin{codeexample}[width=6cm]
	\begin{goban}[label = {at = all}, partial = A1:E6]
		\stones[b]{A1:E6} % Additive group.
	\end{goban}

	\begin{goban}[label = {at = all}, partial = A1:E6]
		\stones[b]{A1:E6; A 1:2} % Additive group; subtractive.
	\end{goban}

	\begin{goban}[label = {at = all}, partial = A1:E6]
		\stones[b]{A1:E6; A 1:2, C6 |- E3} % Additive group; subtractive.
	\end{goban}
\end{codeexample}

In this sequence, first a full rectangle is drawn (|A1:E6|), then the vertical segment |A 1:2| is removed, and finally the path \verb+C6 |- E3+ is taken out. The result is the diagram of the beginning.

Sometimes a pattern is ideal except for one or two stones: a subtractive group removes them without redoing everything.

\begin{codeexample}[width=6cm]
	\begin{goban}[label = {at = all}]
		\stones[b]{A1:E6; B:D 5, C2}
	\end{goban}
\end{codeexample}

Here, |B:D 5| removes a horizontal segment on row 5, and |C2| removes one stone.

Additive and subtractive groups alternate with each |;|, as many times as needed.

\begin{codeexample}[width=6cm]
	\begin{goban}[label = {at = all}, partial = A1:E6]
		\stones[b]{A1:E6} % AD
	\end{goban}

	\begin{goban}[label = {at = all}, partial = A1:E6]
		\stones[b]{A1:E6; A1:C4, E6} % AD, SU
	\end{goban}

	\begin{goban}[label = {at = all}, partial = A1:E6]
		\stones[b]{A1:E6; A1:C4, E6; A 1:2} % AD, SU, AD
	\end{goban}
\end{codeexample}

In this last example, a rectangle is created, two regions are removed (|A1:C4| and |E6|), and then a segment is added (|A 1:2|).

	\section{The Board}
\label{sec:board}

When we refer to the "board," we mean an environment or \emph{scope} in which both the goban (the game grid) and its contents—stones, coordinate labels, and other relevant information—are rendered. This environment is designed to be flexible and customizable, allowing users to adapt its appearance and functionality according to the needs of the diagram, whether to represent full games, specific sections, or even add visual annotations.

\begin{environment}{{goban}\opt{\oarg{options}}}
	
	Initializes the goban environment. By default, it draws a square 9×9 intersection board with a visible border, white background, and no labels on the sides.
	
	\begin{codeexample}[width=5cm]
		\begin{goban}
		\end{goban}
	\end{codeexample}
	
	Since the board is the only macro that utilizes the TikZ layer, each intersection can be accessed using its alphanumeric coordinate (as in \pgfname) or numeric values in the \tikzname\ coordinate system. This allows combining Go’s intuitive syntax with TikZ’s graphical capabilities, which is especially useful for adding custom elements, such as arrows or marks, directly onto the board.
	
	\begin{codeexample}[width=5.85cm]
		\begin{goban}[19, partial = {from = J1 to S9},
			label = {at = {north, east}}]
			\stones {K8,R8,Q5,R5,O4,Q4,L3,O3,M2,N2}
			\mark[black, circle] {K8}
			\stones[white] {K6,R4,R3,P3,N1,O2,Q2}
			
			\tikzset{flecha/.style = {-{latex[]}, line width = 2pt, blue}}
			\draw[flecha] (12,8)--(17,8);
		\end{goban}
	\end{codeexample}
	
	\begin{key}{/pgfgo/goban/use i=\meta{boolean} (default true, initially true)}
		
		Enables or disables the use of the letter "I" in the board’s coordinates and labels. When disabled ("false"), any reference to "I" is interpreted as "J".
		
		\begin{codeexample}[]
			\pgfgoset{goban = {13, label = {at = north}}}
			
			\begin{goban}[label = {at = west}, partial = {from = H1 to M3}]
				\stone[b]{J2}
			\end{goban}
			
			\begin{goban}[use i = false, partial = {from = H1 to N3}]
				\stone[b]{J2}
			\end{goban}
			
			\begin{goban}[use i = false, label = {at = east}, partial = {from = H1 to N3}]
				\stone[white, mark = {custom = !}] {I2} % If it’s I, it shifts to J.
			\end{goban}
		\end{codeexample}
	\end{key}
	
	\begin{key}{/pgfgo/goban/size=\meta{nlines or nlines |x| nlines} (initially 9)}
		
		Defines the number of vertical and horizontal lines on the board. The word "size" can be omitted, and the value specified directly. If a single integer (e.g., |9|) is provided, the board will be square (9×9). If an expression like |5x3| is used, a rectangular board with 5 columns and 3 rows will be created.
		
		\begin{codeexample}[width=6.5cm]
			\begin{goban}[5]
			\end{goban}
			
			\begin{goban}[5x3]
			\end{goban}
		\end{codeexample}
	\end{key}
	
	\begin{key}{/pgfgo/goban/sgf size=\meta{file}}
		Takes the size of the board from an SGF game (its |SZ| property). Only with the Lua kernel (p.~\pageref{subsec:kernel}).

		\begin{codeexample}[]
			\begin{goban}[sgf size = variations.sgf, label = {at = {north, west}}]
				\sgf{variations.sgf}
			\end{goban}
		\end{codeexample}
	\end{key}

	\begin{key}{/pgfgo/goban/scale=\meta{float} (initially 1)}
		
		Controls the global scale of the board, affecting both its overall dimensions and the thickness of certain elements (e.g., the "line width" of the lines).
		
		\begin{codeexample}[width=5.5cm]
			\begin{goban}[2, scale = 2]
				\stone[b, mark]{A1}
			\end{goban}
			
			\begin{goban}[2, scale = 1]
				\stone[b, mark]{A1}
			\end{goban}
			
			\begin{goban}[2, scale = 0.5]
				\stone[b, mark]{A1}
			\end{goban}
		\end{codeexample}
	\end{key}
	
	\subsection{Grid}
	
	The grid refers to the central lines that form the board’s intersections.
	
	\begin{key}{/pgfgo/goban/grid=\meta{options}}
		Controls the behavior and appearance of the board’s central lines. This set of options allows modifying their visibility, color, thickness, and spacing, offering a high degree of customization.
	\end{key}
	
	\begin{stylekey}{/pgfgo/goban/grid/true}

		Turns the central lines back on after a |grid = false| (for example, a global one set with |\pgfgoset|).
	\end{stylekey}

	\begin{stylekey}{/pgfgo/goban/grid/false}
		
		Disables the central lines, leaving only the border and stones visible.
		
		\begin{codeexample}[]
			\begin{goban}[3]
				\stone[b]{A1}
				\stone[w]{B2}
			\end{goban}
			
			\begin{goban}[3, grid = false]
				\stone[b]{A1}
				\stone[w]{B2}
			\end{goban}
		\end{codeexample}
		
		The first board shows the default grid.
	\end{stylekey}
	
	\begin{key}{/pgfgo/goban/grid/line=\meta{options}}
		Allows controlling the grid lines.
	\end{key}
	
	\begin{key}{/pgfgo/goban/grid/line/color =\meta{color} (initially black)}
		
		Allows changing the color of the grid lines.
		
		\begin{codeexample}[]
			\pgfgoset{goban = {5, partial = {from = B2 to D4}}}
			
			\begin{goban}[grid = {line = {color = blue}}, new remember]
				\stones[b]{C2,D3}
				\stone[w]{B4}
			\end{goban}
			
			\begin{goban}[grid = {line = {color = red}}, resume]
			\end{goban}
		\end{codeexample}
	\end{key}
	
	\begin{key}{/pgfgo/goban/grid/line/opacity =\meta{float} (initially 1)}
		
		Sets the opacity of the grid lines, from 0 (transparent) to 1 (opaque).
		
		\begin{codeexample}[]
			\pgfgoset{goban = {5, partial = {from = B2 to D4}}}
			
			\begin{goban}[new remember]
				\stones[b]{C2,D3}
				\stone[w]{B4}
			\end{goban}
			
			\begin{goban}[grid = {line = {opacity = 0.2}}, resume]
			\end{goban}
		\end{codeexample}
	\end{key}
	
	\begin{key}{/pgfgo/goban/grid/line/width =\meta{dimension} (initially 0.55pt)}
		
		Defines the thickness of the grid lines.
		
		\begin{codeexample}[]
			\pgfgoset{goban = {5, partial = {from = B2 to D4}}}
			
			\begin{goban}[new remember]
				\stones[b]{C2,D3}
				\stone[w]{B4}
			\end{goban}
			
			\begin{goban}[grid = {line = {width = 1.5pt}}, resume]
			\end{goban}
		\end{codeexample}
	\end{key}
	
	\begin{key}{/pgfgo/goban/grid/x sep =\meta{dimension} (initially 1.25em)}
		Sets the spacing between vertical lines, adjusting the width of the board’s columns.
	\end{key}
	
	\begin{key}{/pgfgo/goban/grid/y sep =\meta{dimension} (initially 1.25em)}
		Sets the spacing between horizontal lines, adjusting the height of the rows.
	\end{key}
	
	\begin{key}{/pgfgo/goban/grid/sep =\meta{dimension} (initially 1.25em)}
		Simultaneously defines the vertical and horizontal spacing of the lines, ensuring uniform spacing in both directions. This is a convenient option for maintaining balanced proportions.
	\end{key}
	
	\subsection{Border}
	
	The board’s border is the outline that frames the grid and stones, providing a clear visual boundary. By default, it is active, has a white background, and a black outline.
	
	\begin{key}{/pgfgo/goban/background=\meta{options}}
		Controls the properties of the board’s border, such as its visibility, color, thickness, and shape, allowing detailed customization.
	\end{key}
	
	\begin{stylekey}{/pgfgo/goban/background/true}

		Turns the board’s background and border back on after a |background = false|.
	\end{stylekey}

	\begin{stylekey}{/pgfgo/goban/background/false}
		
		Completely disables the board’s border, removing both the outline and fill. This is useful for minimalist diagrams or when the document’s background already provides sufficient contrast.
		
		\begin{codeexample}[]
			\pgfgoset{goban = {5, partial = {from = B2 to D4}}}
			
			\begin{goban}[new remember]
				\stones[b]{C2,D3}
				\stone[w]{B4}
			\end{goban}
			
			\begin{goban}[background = false, resume]
			\end{goban}
		\end{codeexample}
	\end{stylekey}
	
	\begin{key}{/pgfgo/goban/background/fill=\meta{options}}
		Controls options related to the board’s background fill.
	\end{key}
	
	\begin{key}{/pgfgo/goban/background/fill=\meta{boolean} (initially true)}
		
		Determines whether the board’s background fill is active. If disabled, the area inside the border will be transparent, showing only the outline (if enabled).
		
		\begin{codeexample}[]
			\pgfgoset{goban = {5, partial = {from = B2 to D4}}}
			
			\begin{goban}[background = {line = {color = red}}, new remember]
				\stones[b]{C2,D3}
				\stone[w]{B4}
			\end{goban}
			
			\begin{goban}[background = {line = false}, resume]
			\end{goban}
		\end{codeexample}
	\end{key}
	
	\begin{key}{/pgfgo/goban/background/fill/color =\meta{color} (initially white)}
		
		Defines the fill color of the board.
		
		\begin{codeexample}[]
			\pgfgoset{goban = {5, partial = {from = B2 to D4}}}
			
			\begin{goban}[background = {fill = {color = orange!40}}, new remember]
				\stones[b]{C2,D3}
				\stone[w]{B4}
			\end{goban}
			
			\begin{goban}[background = {fill = {color = olive!30}}, resume]
			\end{goban}
		\end{codeexample}
	\end{key}
	
	\begin{key}{/pgfgo/goban/background/fill/opacity =\meta{float} (initially 1)}
		
		Controls the opacity of the fill, from 0 (transparent) to 1 (opaque).
		
		\begin{codeexample}[]
			\pgfgoset{goban = {5, partial = {from = B2 to D4}}}
			
			\begin{goban}[background = {fill = {color = orange!40}}, new remember]
				\stones[b]{C2,D3}
				\stone[w]{B4}
			\end{goban}
			
			\begin{goban}[background = {fill = {color = orange!40}, fill = {opacity = 0.5}}, resume]
			\end{goban}
		\end{codeexample}
		
	\end{key}
	
	\begin{key}{/pgfgo/goban/background/line=\meta{options}}
		Controls the options for the background’s outline.
	\end{key}
	
	\begin{key}{/pgfgo/goban/background/line=\meta{boolean} (initially true)}
		Enables or disables the border line.
		
		\begin{codeexample}[]
			\pgfgoset{goban = {5, partial = {from = B2 to D4}}}
			
			\begin{goban}[background = {line = {color = red}}, new remember]
				\stones[b]{C2,D3}
				\stone[w]{B4}
			\end{goban}
			
			\begin{goban}[background = {line = false}, resume]
			\end{goban}
		\end{codeexample}
	\end{key}
	
	\begin{key}{/pgfgo/goban/background/line/color =\meta{color} (initially black)}
		
		Defines the color of the border line, allowing it to stand out from the background or grid lines.
		
		\begin{codeexample}[]
			\pgfgoset{goban = {5, partial = {from = B2 to D4}}}
			
			\begin{goban}[background = {line = {color = red}}, new remember]
				\stones[b]{C2,D3}
				\stone[w]{B4}
			\end{goban}
			
			\begin{goban}[background = {line = {color = blue}}, resume]
			\end{goban}
		\end{codeexample}
	\end{key}
	
	\begin{key}{/pgfgo/goban/background/line/opacity =\meta{float} (initially 1)}
		
		Adjusts the opacity of the border line.
		
		\begin{codeexample}[]
			\pgfgoset{goban = {5, partial = {from = B2 to D4}}}
			
			\begin{goban}[background = {line = {color = red}}, new remember]
				\stones[b]{C2,D3}
				\stone[w]{B4}
			\end{goban}
			
			\begin{goban}[background = {line = {color = red}, line = {opacity = 0.5}}, resume]
			\end{goban}
		\end{codeexample}
	\end{key}
	
	\begin{key}{/pgfgo/goban/background/line/width =\meta{dimension} (initially 0.75pt)}
		
		Controls the thickness of the border line.
		
		\begin{codeexample}[]
			\pgfgoset{goban = {5, partial = {from = B2 to D4}}}
			
			\begin{goban}[new remember]
				\stones[b]{C2,D3}
				\stone[w]{B4}
			\end{goban}
			
			\begin{goban}[background = {line = {width = 1.5pt}}, resume]
			\end{goban}
		\end{codeexample}
	\end{key}
	
	\begin{key}{/pgfgo/goban/background/rounded corners =\meta{dimension} (initially 0pt)}
		
		Adjusts the rounding of the border corners, from 0pt (straight corners) to larger values for a smoother, more stylized effect.
		
		\begin{codeexample}[]
			\pgfgoset{goban = {5, partial = {from = A1 to C3}}}
			
			\begin{goban}[new remember]
				\stones[b]{B1,C2}
				\stone[w]{A3}
			\end{goban}
			
			\begin{goban}[background = {rounded corners = 10pt}, resume]
			\end{goban}
		\end{codeexample}
	\end{key}
	
	\begin{key}{/pgfgo/goban/background/rounded corners/outer =\meta{dimension} (initially 0pt)}
		
		Adjusts only the rounding of the outer corners of the border.
		
		\begin{codeexample}[]
			\pgfgoset{goban = {5, partial = {from = A1 to C3}}}
			
			\begin{goban}[background = {rounded corners = 0pt}, new remember]
				\stones[b]{B1,C2}
				\stone[w]{A3}
			\end{goban}
			
			\begin{goban}[background = {rounded corners = {outer = 10pt}}, resume]
			\end{goban}
		\end{codeexample}
	\end{key}
	
	\begin{key}{/pgfgo/goban/background/rounded corners/inner =\meta{dimension} (initially 0pt)}
		
		Adjusts only the rounding of the inner corners of the board (only visible on partial boards).
		
		\begin{codeexample}[]
			\pgfgoset{goban = {5, partial = {from = A1 to C3}}}
			
			\begin{goban}[background = {rounded corners = 0pt}, new remember]
				\stones[b]{B1,C2}
				\stone[w]{A3}
			\end{goban}
			
			\begin{goban}[background = {rounded corners = {outer = 8pt, inner = 2pt}}, resume]
			\end{goban}
		\end{codeexample}
	\end{key}
	
	\begin{key}{/pgfgo/goban/background/sep =\meta{dimension} (initially 0.825em)}
		
		Defines the separation between the grid and the border, based by default on the stone radius multiplied by 1.5. This adjustment ensures that stones near the border do not overlap with it.
		
		\begin{codeexample}[]
			\pgfgoset{goban = {5, partial = {from = B2 to D4}}}
			
			\begin{goban}[background = {sep = 0.625em}, new remember]
				\stones[b]{C2,D3}
				\stone[w]{B4}
			\end{goban}
			
			\begin{goban}[background = {sep = 1.5em}, resume]
			\end{goban}
		\end{codeexample}
	\end{key}
	
	\begin{key}{/pgfgo/goban/background/drop shadow=\meta{options} (default true)}
		Draws a shadow under the board. Without a value, it turns it on.
		\begin{codeexample}[]
			\pgfgoset{goban = {5, partial = {from = B2 to D4}}}

			\begin{goban}[background = {drop shadow}, new remember]
				\stones[b]{C2,D3}
				\stone[w]{B4}
			\end{goban}

			\begin{goban}[background = {drop shadow = {color = blue, opacity = 0.3}}, resume]
			\end{goban}
		\end{codeexample}
	\end{key}

	\begin{key}{/pgfgo/goban/background/drop shadow=\meta{boolean} (initially false)}
		Turns the shadow on (|true|) or off (|false|).
	\end{key}

	\begin{key}{/pgfgo/goban/background/drop shadow/angle =\meta{float} (initially -45)}
		Direction of the shadow, in degrees.
		\begin{codeexample}[]
			\pgfgoset{goban = {5, partial = {from = B2 to D4}}}

			\begin{goban}[background = {drop shadow}, new remember]
				\stones[b]{C2,D3}
				\stone[w]{B4}
			\end{goban}

			\begin{goban}[background = {drop shadow = {angle = 45}}, resume]
			\end{goban}
		\end{codeexample}
	\end{key}

	\begin{key}{/pgfgo/goban/background/drop shadow/offset =\meta{dimension} (initially 1.5pt)}
		Distance between the board and its shadow.
		\begin{codeexample}[]
			\pgfgoset{goban = {5, partial = {from = B2 to D4}}}

			\begin{goban}[background = {drop shadow}, new remember]
				\stones[b]{C2,D3}
				\stone[w]{B4}
			\end{goban}

			\begin{goban}[background = {drop shadow = {offset = 4pt}}, resume]
			\end{goban}
		\end{codeexample}
	\end{key}

	\begin{key}{/pgfgo/goban/background/drop shadow/opacity =\meta{float} (initially 0.5)}
		Opacity of the shadow.
	\end{key}

	\begin{key}{/pgfgo/goban/background/drop shadow/color =\meta{color} (initially black)}
		Color of the shadow.
	\end{key}

	\subsection{Partial Board}
	
	In many cases, it’s not necessary to represent a full board, but rather just a rectangular section—whether due to space constraints, to focus attention on a specific area, or to simplify diagram creation. The |goban| environment includes the \emph{key} |partial|, which allows defining this region using the sub-\emph{key} |from|. This specifies the opposite vertices of a rectangle, regardless of the order in which they are provided, making it easy to use. For example, |from = A1 to D3| is equivalent to |from = D3 to A1|.
	
	\begin{figure}[H]
		\centering
		
		\begin{tikzpicture}
			\node[anchor=north] (X) {
				\begin{goban}[label = {at = {north, west}}]
					\coordinate (A) at (4,6);
					\coordinate (B) at (9,9);
					\coordinate (C) at (A|-B);
					\coordinate (D) at (A-|B);
					
					\tikzset{punto/.style = {circle, fill = #1, draw, inner sep=0.25em, fill opacity=0.5, anchor=center}}
					
					\node[punto = red] at (A){};
					\node[punto = red] at (B){};
					\node[punto = blue] at (C){};
					\node[punto = blue] at (D){};
					
					\draw[red, {latex}-{latex}, shorten <=5pt, shorten >= 5pt] (A)--(B);
					\draw[blue, {latex}-{latex}, shorten <=5pt, shorten >= 5pt] (C)--(D);
					
					\node[fit=(A)(B), draw, inner sep=0.6em, rounded corners, line width=1pt, fill opacity=0.1, fill=orange] {};
				\end{goban}
			};
			
			\node[anchor=north, xshift=15em] {
				\tikz{
					\node[inner sep=0] (Y) {
						\begin{goban}[label = {at = {north, west}}, partial = {from = I9 to D6}]
						\end{goban}
					};	
					\node[align = left, anchor=north, inner sep=0, yshift=-1em] at (Y.south){
						{\color{blue}|partial = {from = D9 to I6}|}\\
						{\color{blue}|partial = {from = I6 to D9}|}\\
						{\color{red} |partial = {from = D6 to I9}|}\\
						{\color{red} |partial = {from = I9 to D6}|}\\
					};
				}
			};
		\end{tikzpicture}
	\end{figure}
	
	\begin{key}{/pgfgo/goban/partial=\meta{options}}
		
		Controls options for representing partial boards, allowing the delimitation of a specific region of the goban.
	\end{key}
	
	\begin{key}{/pgfgo/goban/partial/from=\meta{intersection |to| intersection}}
		
		Specifies the opposite vertices of the rectangle defining the partial board. These intersections can be any pair of points forming a diagonal of the desired area.
		
	\end{key}
	
	\begin{codeexample}[]
		\tikzset{flecha/.style = {-{latex[]}, line width = 1pt, red}}
		
		\begin{goban}[partial = {from = A1 to D3}]
			\draw[flecha] (1,1)--(4,3); % From A1 to D3
		\end{goban}
		
		\begin{goban}[partial = {from = A9 to I6}]
			\draw[flecha] (1,9)--(9,6); % From A9 to I6
		\end{goban}
		
		\begin{goban}[partial = {from = E2 to B4}]
			\draw[flecha] (5,2)--(2,4); % From E2 to B4
		\end{goban}
	\end{codeexample}
	
	\begin{key}{/pgfgo/goban/partial/false}
		Shows the whole board again (for example, after a global |partial| set with |\pgfgoset|).
	\end{key}

	\subsection{Board Labels}
	
	Labels are the alphanumeric indicators (letters and numbers) that appear on the sides of the board to identify intersections.
	
	\begin{key}{/pgfgo/goban/label=\meta{options}}
		Controls the properties of the labels, such as their position and format.
	\end{key}
	
	\begin{key}{/pgfgo/goban/label/false}
		Disables all labels.
	\end{key}
	
	\begin{key}{/pgfgo/goban/label/at =\meta{options}}
		
		Defines on which sides of the board the labels appear. Options include combinations of cardinal positions or predefined values like "all" or "none".
		
		\begin{stylekey}{/pgfgo/goban/label/at/all}
			
			Displays labels on all sides of the board (north, south, east, and west).
		\end{stylekey}
		
		\begin{stylekey}{/pgfgo/goban/label/at/none}
			
			Equivalent to |label=false|, removing all labels.
		\end{stylekey}
		
		\begin{stylekey}{/pgfgo/goban/label/at/east}
			Places labels only on the right side (vertical numbers).
		\end{stylekey}
		\begin{stylekey}{/pgfgo/goban/label/at/right}
			Equivalent to |east|.
		\end{stylekey}
		
		\begin{stylekey}{/pgfgo/goban/label/at/north}
			Places labels only at the top (horizontal letters).
		\end{stylekey}
		\begin{stylekey}{/pgfgo/goban/label/at/above}
			Equivalent to |north|.
		\end{stylekey}
		
		\begin{stylekey}{/pgfgo/goban/label/at/south}
			Places labels only at the bottom (horizontal letters).
		\end{stylekey}
		\begin{stylekey}{/pgfgo/goban/label/at/below}
			Equivalent to |south|.
		\end{stylekey}
		
		\begin{stylekey}{/pgfgo/goban/label/at/west}
			Places labels only on the left side (vertical numbers).
		\end{stylekey}
		\begin{stylekey}{/pgfgo/goban/label/at/left}
			Equivalent to |west|.
		\end{stylekey}
		
		\begin{codeexample}[]
			\begin{goban}[12, label = {at = {all}}]
			\end{goban}
			
			\begin{goban}[partial = {from = A6 to I9}, label = {at = north}]
			\end{goban}
			
			\begin{goban}[partial = {from = B2 to E4}, label = {at = {north, east}}]
			\end{goban}
		\end{codeexample}
	\end{key}
	
	\begin{key}{/pgfgo/goban/label/font size = \meta{dimension} (initially the current font size)}

		Adjusts the font size of the labels.

		\begin{codeexample}[]
			\begin{goban}[5, label = {at = {south, west}}]
			\end{goban}

			\begin{goban}[5, label = {at = {south, west}, font size = 6pt}]
			\end{goban}
		\end{codeexample}
	\end{key}

	\begin{key}{/pgfgo/goban/label/font = \meta{font commands} (initially \texttt{\string\sffamily})}

		Sets the font of the labels.

		\begin{codeexample}[]
			\begin{goban}[5, label = {at = {south, west}, font = \bfseries}]
			\end{goban}
		\end{codeexample}
	\end{key}
	
	\begin{key}{/pgfgo/goban/label/sep = \meta{dimension} (initially 0.35cm)}
		
		Controls the distance between the labels and the board’s border.
	\end{key}
	
	\begin{key}{/pgfgo/goban/label/text color = \meta{color} (initially black)}
		
		Defines the text color of the labels.
	\end{key}
	
	\subsection{Captures}
	\label{subsec:captures}

	With the Lua kernel (p.~\pageref{subsec:kernel}), \pgfgoname\ can keep track of the stones on the board and remove the ones left without liberties, as in a game.

	\begin{key}{/pgfgo/goban/captures=\meta{remove, fade or false} (default remove, initially false)}

		Turns captures on. When a move leaves an opposing group without liberties, that group is removed from the board. If it leaves its own group without liberties (a suicide), that group is removed. Stones placed with |\stone| and |\stones| also count; neutral stones and prisoners do not. A move outside a partial board also captures. A move on an occupied point produces a warning, because it is not possible in a game. With |kernel = expl3|, this \emph{key} raises an error.

		\begin{codeexample}[]
			\pgfgoset{goban = {5, partial = {from = A1 to C3}}}

			\begin{goban}
				\moves*{B1,A1,A2}
			\end{goban}

			\begin{goban}[captures]
				\moves*{B1,A1,A2} % 3 captures 2.
			\end{goban}
		\end{codeexample}

		Since \pgfname\ cannot erase what it has already drawn, with |captures| the stones are drawn when the |goban| ends. For this reason, if the style is changed with |\pgfgoset| inside a |goban| with captures, the change does not reach the stones; the options of each command do.
	\end{key}

	\begin{stylekey}{/pgfgo/goban/captures/remove}

		Captured stones disappear. It is the same as |captures| or |captures = true|.
	\end{stylekey}

	\begin{stylekey}{/pgfgo/goban/captures/fade}

		Captured stones stay on the board, drawn as prisoners (with the style of |\prisoner| and |/pgfgo/prisoner|).

		\begin{codeexample}[]
			\pgfgoset{goban = {5, partial = {from = A1 to E3}}}

			\begin{goban}[captures = fade]
				\stones[w]{C1,D1}
				\stones[b]{B1,C2,D2}
				\moves*{E1}
			\end{goban}
		\end{codeexample}
	\end{stylekey}

	\begin{stylekey}{/pgfgo/goban/captures/false}

		Turns captures off. It is the initial value.
	\end{stylekey}

	\begin{command}{\pgfgoprisoners\marg{player}}
		Writes how many prisoners the player (|black| or |b|, |white| or |w|) has in the last |goban| with |captures|: the opponent’s stones that their moves captured.

		\begin{codeexample}[]
			\pgfgoset{goban = {5, partial = {from = A1 to C3}}}

			\begin{goban}[captures]
				\moves*{B1,A1,A2}
			\end{goban}

			Black: \pgfgoprisoners{black}. White: \pgfgoprisoners{white}.
		\end{codeexample}
	\end{command}

	\subsection{Moves on Points Already Played}
	\label{subsec:notes}

	In numbered diagrams, when a move is played on a point where there already was a move (for example, when a ko is retaken), books do not draw it: the point still shows the first stone, and below the diagram they write, for example, "3 at 1".

	\begin{key}{/pgfgo/goban/notes=\meta{boolean} (default true, initially false)}

		A numbered move on a point where there already was a numbered move of the same |goban| is not drawn: the point still shows the first one, even if it was captured, and "N at M" is noted. Only numbered moves (|\moves*|) that are visible on the board count. It works with both kernels.

		\begin{codeexample}[]
			\pgfgoset{goban = {5, partial = {from = A2 to E4}}}

			\begin{goban}[notes]
				\stones{B3,C2,C4}
				\stones[w]{D2,D4,E3}
				\moves*[w]{C3}
				\moves*{D3}
				\moves*[from = 3, w]{C3}
			\end{goban}

			\pgfgomovenotes
		\end{codeexample}
	\end{key}

	\begin{command}{\pgfgomovenotes\opt{\oarg{options}}}
		Writes the notes of the last |goban|, separated by commas. If there are no notes, it writes nothing. The options are those of |/pgfgo/notes|.
	\end{command}

	\begin{key}{/pgfgo/notes = \meta{options}}
		Options of |\pgfgomovenotes| for the whole document, for example |\pgfgoset{notes = {at = en}}|.
	\end{key}

	\begin{key}{/pgfgo/notes/at = \meta{text} (initially at)}
		Text between the two numbers.
	\end{key}

	\begin{key}{/pgfgo/notes/separator = \meta{text} (initially a comma and a space)}
		Text between the notes.
	\end{key}

	\subsection{Remembering a Board}
	The "remember" system stores what is drawn in a |goban| (stones, moves and marks, with their options) to draw it again in another one with |resume|. Moves keep their number and color, and the numbering continues from there. Only what is drawn is stored: the board options (size, |partial|, labels) are given in each |goban|.

	\begin{codeexample}[]
		\pgfgoset{goban = {5, partial = {from = B2 to D4}}}

		\begin{goban}[new remember]
			\stones[b]{C2,D3}
			\moves*[w]{B4}
		\end{goban}

		\begin{goban}[resume]
			\moves*{C3} % Continues with 2 (Black).
		\end{goban}
	\end{codeexample}

	\begin{stylekey}{/pgfgo/goban/new remember}
		Empties the memory and stores what is drawn in this |goban|.
	\end{stylekey}

	\begin{stylekey}{/pgfgo/goban/remember}
		Adds what is drawn in this |goban| to the memory, without erasing what was there.
	\end{stylekey}

	\begin{stylekey}{/pgfgo/goban/add to remember}
		The same as |remember|.
	\end{stylekey}

	\begin{stylekey}{/pgfgo/goban/forget}
		Empties the memory and stores nothing from this |goban|, except what is drawn with the option |remember| (for example, |\stone[remember]{C3}|).
	\end{stylekey}

	\begin{stylekey}{/pgfgo/goban/resume}
		Draws what is in the memory at the beginning of the |goban|.
	\end{stylekey}

\end{environment}
	\section{Stones, Moves, and Prisoners}

Stones are the fundamental elements. In \pgfgoname, commands and options are provided to draw them flexibly, allowing customization of their appearance and behavior according to the needs of the diagram.

\subsection{Stones}

\begin{command}{\stone\opt{\oarg{options}}\marg{intersection}}
	This command draws an individual stone at the specified intersection. The mandatory argument is the coordinate (e.g., "C2"), while optional options allow adjusting its style, such as color, size, or markings, following the convention of the first player in Go.
\end{command}

\begin{command}{\stones\opt{\oarg{options}}\marg{intersection \opt{, intersection, ...}}}
	Similar to |\stone|, but allows placing multiple stones using the corresponding syntax (see Coordinate Syntax).
\end{command}

It is possible to change an option for all players using |\pgfgoset{stone = {...}}|. This is equivalent to |\pgfgoset{players = {all = {stone = {...}}}}|.

\begin{key}{/pgfgo/stone/fill = \meta{option} (default true)}
	Controls the fill options of a stone.
\end{key}

\begin{key}{/pgfgo/stone/fill/color = \meta{color}}
	Defines the fill color of a stone.
	
	\begin{codeexample}[]
		\pgfgoset{goban = {5, partial = {from = B2 to D4}}}
		
		\begin{goban}
			\stones[b]{C2,D3}
			\stone[w]{B4}
		\end{goban}
		
		\begin{goban}
			\stones[fill = {color = cyan!80!black}]{C2,D3}
			\stone[fill = {color = magenta!80!black}]{B4}
		\end{goban}
	\end{codeexample}
\end{key}

\begin{key}{/pgfgo/stone/fill/opacity = \meta{float} (initially 1)}
	
	Controls the opacity of the stone’s fill, ranging from 0 (transparent) to 1 (opaque).
	
	\begin{codeexample}[]
		\pgfgoset{goban = {5, partial = {from = B2 to D4}}}
		
		\begin{goban}
			\stones[b]{C2,D3}
			\stone[w]{B4}
		\end{goban}
		
		\begin{goban}
			\stones[b, fill = {opacity = 0.8}]{C2,D3}
			\stone[w, fill = {opacity = 0.5}]{B4}
		\end{goban}
	\end{codeexample}
\end{key}

\begin{key}{/pgfgo/stone/line = \meta{option} (initially true)}
	Controls the border options of a stone.
\end{key}

\begin{key}{/pgfgo/stone/line/color = \meta{color} (initially black)}
	
	Sets the border color of the stone, which is black by default for all stones.
	
	\begin{codeexample}[]
		\pgfgoset{goban = {5, partial = {from = B2 to D4}}}
		
		\begin{goban}
			\stones[b]{C2,D3}
			\stone[w]{B4}
		\end{goban}
		
		\begin{goban}
			\stones[b, line = {color = magenta}]{C2,D3}
			\stone[w, line = {color = cyan}]{B4}
		\end{goban}
	\end{codeexample}
\end{key}

\begin{key}{/pgfgo/stone/line/opacity = \meta{float} (initially 1)}
	
	Defines the opacity of the stone’s border, allowing it to be more subtle or transparent.
	
	\begin{codeexample}[]
		\pgfgoset{goban = {5, partial = {from = B2 to D4}}}
		
		\begin{goban}
			\stones[b]{C2,D3}
			\stone[w]{B4}
		\end{goban}
		
		\begin{goban}
			\stones[b, line = {opacity = 0.5}]{C2,D3}
			\stone[w, line = {opacity = 0.8}]{B4}
		\end{goban}
	\end{codeexample}
\end{key}

\begin{key}{/pgfgo/stone/line/width = \meta{dimension} (initially 0.9pt)}
	
	Controls the thickness of the stone’s border.
	
	\begin{codeexample}[]
		\pgfgoset{goban = {5, partial = {from = B2 to D4}}}
		
		\begin{goban}
			\stones[b]{C2,D3}
			\stone[w]{B4}
		\end{goban}
		
		\begin{goban}
			\stones[b]{C2,D3}
			\stone[w, line = {width = 2pt}]{B4}
		\end{goban}
	\end{codeexample}
\end{key}

\begin{key}{/pgfgo/stone/radius = \meta{dimension} (initially 0.55em)}
	
	Sets the size of the stone by its radius.
	
	\begin{codeexample}[]
		\pgfgoset{goban = {5, partial = {from = B2 to D4}}}
		
		\begin{goban}
			\stones[b]{C2,D3}
			\stone[w]{B4}
		\end{goban}
		
		\begin{goban}
			\stones[b, radius = 0.15cm]{C2,D3}
			\stone[w, radius = 0.25cm]{B4}
		\end{goban}
	\end{codeexample}
\end{key}

\begin{key}{/pgfgo/stone/scale = \meta{float} (initially 1)}
	
	Defines a scaling factor that affects all elements of the stone (radius, border, markings), maintaining their proportions.
\end{key}

\begin{key}{/pgfgo/stone/black}
	
	Assigns the stone the characteristics of the first player (black by default), including fill and border color. This is the initial state of any stone unless another option is specified.
	
	\begin{codeexample}[]
		\pgfgoset{goban = {5, partial = {from = B2 to D4}}}
		
		\begin{goban}
			\stones[b]{C2,D3}
			\stone[w]{B4}
		\end{goban}
		
		\begin{goban}
			\stones[b]{C2,D3}
			\stone[black]{B4}
		\end{goban}
	\end{codeexample}
\end{key}

\begin{key}{/pgfgo/stone/b}
	
	Abbreviation for "black," with identical functionality.
\end{key}

\begin{key}{/pgfgo/stone/white}
	
	Assigns the stone the characteristics of the second player (white by default), with a white fill and black border.
	
	\begin{codeexample}[]
		\pgfgoset{goban = {5, partial = {from = B2 to D4}}}
		
		\begin{goban}
			\stones[b]{C2,D3}
			\stone[white]{B4}
		\end{goban}
		
		\begin{goban}
			\stones[white]{C2,D3}
			\stone[white]{B4}
		\end{goban}
	\end{codeexample}
\end{key}

\begin{key}{/pgfgo/stone/w}
	
	Abbreviation for "white," with the same functionality.
\end{key}

\begin{key}{/pgfgo/stone/mark = \meta{mark keys}}
	
	Places a mark on the stone (e.g., triangle, circle). Marks can be further customized with other keys, such as color or size.
	
	\begin{codeexample}[]
		\pgfgoset{goban = {5, partial = {from = B2 to D4}}}
		
		\begin{goban}
			\stones[b]{C2,D3}
			\stone[w]{B4}
		\end{goban}
		
		\begin{goban}
			\stones[b, mark = triangle]{C2,D3}
			\stone[w, mark = t-circle]{B4}
		\end{goban}
	\end{codeexample}
\end{key}

\begin{key}{/pgfgo/stone/forget}
	
	When the board uses a memory system (|new remember|, |resume|), this option indicates that the stone should not be remembered, allowing it to be excluded from subsequent states.
	
	\begin{codeexample}[]
		\pgfgoset{goban = {5, partial = {from = B2 to D4}}}
		
		\begin{goban}[new remember]
			\stones[b]{C2,D3}
			\stone[w, forget]{B4}
		\end{goban}
		
		\begin{goban}[resume]
		\end{goban}
	\end{codeexample}
\end{key}

\begin{key}{/pgfgo/stone/remember}
	
	In a board with active memory, ensures that the stone is remembered even if the general environment uses |forget|.
	
	\begin{codeexample}[]
		\pgfgoset{goban = {5, partial = {from = B2 to D4}}}
		
		\begin{goban}[forget]
			\stones[b]{C2,D3}
			\stone[w, remember]{B4}
		\end{goban}
		
		\begin{goban}[resume]
		\end{goban}
	\end{codeexample}
\end{key}

\begin{key}{/pgfgo/stone/prisoner=\meta{prisoner options} (default fill opacity = 0.5)}
	
	Characterizes a stone as a prisoner, typically with reduced opacity (0.5 by default) to indicate it has been captured. Other options can be added to customize its appearance.
	
	\begin{codeexample}[]
		\pgfgoset{goban = {5, partial = {from = B2 to D4}}}
		
		\begin{goban}
			\stones[b]{C2,D3}
			\stone[w]{B4}
		\end{goban}
		
		\begin{goban}
			\stones[b, prisoner]{C2,D3}
			\stone[white]{B4}
		\end{goban}
	\end{codeexample}
\end{key}

\subsection{Moves}

Move commands allow representing sequences of plays, automatically alternating between black and white. They include options for numbering the moves and customizing the associated stones.

\begin{command}{\move\opt{*}\opt{\oarg{options}}\marg{intersection}}
	
	Draws a stone as part of a move sequence, automatically alternating colors (black, white, black, etc.). The starred version (|\move*|) adds a numeric label to the stone.
\end{command}

\begin{command}{\moves\opt{*}\opt{\oarg{options}}\marg{intersection \opt{, intersection, ...}}}
	
	Similar to |\move|, but applies the sequence to multiple intersections in a single line.
\end{command}

It is possible to change options for all players using |\pgfgoset{move = {...}}|, which is equivalent to |\pgfgoset{players = {all = {move = {...}}}}|.

\begin{key}{/pgfgo/move/black}
	
	Sets the move sequence to start with the black player, useful for forcing the start of a play in specific diagrams.
	
	\begin{codeexample}[]
		\pgfgoset{goban = {5, partial = {from = B2 to D4}}}
		
		\begin{goban}
			\stones[b]{C2,D3}
			\stone[w]{B4}
		\end{goban}
		
		\begin{goban}
			\moves[b]{C2,B4}
			\move{D3}
		\end{goban}
	\end{codeexample}
\end{key}

\begin{key}{/pgfgo/move/white}
	
	Sets the sequence to start with the white player, altering the natural order of moves.
	
	\begin{codeexample}[]
		\pgfgoset{goban = {5, partial = {from = B2 to D4}}}
		
		\begin{goban}
			\stones[b]{C2,D3}
			\stone[w]{B4}
		\end{goban}
		
		\begin{goban}
			\moves[white]{C2,B4}
			\move{D3}
		\end{goban}
	\end{codeexample}
\end{key}

\begin{key}{/pgfgo/move/from = \meta{number}}
	
	Defines the starting number of a numbered move sequence, allowing continuation from a specific point instead of starting at 1.
	
	\begin{codeexample}[]
		\pgfgoset{goban = {5, partial = {from = B2 to D4}}}
		
		\begin{goban}
			\stones[b]{C2,D3}
			\stone[w]{B4}
		\end{goban}
		
		\begin{goban}
			\moves*[from = 7]{C2,B4}
			\move*{D3}
		\end{goban}
	\end{codeexample}
\end{key}

\begin{key}{/pgfgo/move/label = \meta{options}}
	
	Controls the properties of numeric labels in starred moves, such as size, color, or format.
\end{key}

%\begin{key}{/pgfgo/move/label/font size=\meta{dimension} (initially 0.5em)}
%	
%	Adjusts the font size of the labels, allowing them to adapt to the stone size or overall design.
%	
%	\begin{codeexample}[]
%		\pgfgoset{goban = {5, partial = {from = B2 to D4}}}
%		
%		\begin{goban}
%			\stones[b]{C2,D3}
%			\stone[w]{B4}
%		\end{goban}
%		
%		\begin{goban}
%			\moves*[from = 7]{C2,B4}
%			\move*[label = {font size = 0.125em}]{D3}
%		\end{goban}
%	\end{codeexample}
%\end{key}

\begin{key}{/pgfgo/move/label/format = \meta{options} (initially arabic)}
	
	Defines the numbering style of the labels.
\end{key}

\begin{key}{/pgfgo/move/label/format/arabic}
	
	Uses Arabic numerals, the default format for move labels.
	
	\begin{codeexample}[]
		\pgfgoset{goban = {5, partial = {from = B2 to D4}}}
		
		\begin{goban}
			\stones[b]{C2,D3}
			\stone[w]{B4}
		\end{goban}
		
		\begin{goban}
			\moves*[from = 7]{C2,B4}
			\move*{D3}
		\end{goban}
	\end{codeexample}
\end{key}

\begin{key}{/pgfgo/move/label/format/alph}
	
	Uses lowercase letters for the labels.
	
	\begin{codeexample}[]
		\pgfgoset{goban = {5, partial = {from = B2 to D4}}}
		
		\begin{goban}
			\stones[b]{C2,D3}
			\stone[w]{B4}
		\end{goban}
		
		\pgfgoset{move={label = {format=alph}}}
		\begin{goban}
			\moves*[from = 7]{C2,B4}
			\move*{D3}
		\end{goban}
	\end{codeexample}
\end{key}

\begin{key}{/pgfgo/move/label/format/Alph}
	
	Uses uppercase letters.
	
	\begin{codeexample}[]
		\pgfgoset{goban = {5, partial = {from = B2 to D4}}}
		
		\begin{goban}
			\stones[b]{C2,D3}
			\stone[w]{B4}
		\end{goban}
		
		\pgfgoset{move = {label = {format = Alph}}}
		\begin{goban}
			\moves*[from = 7]{C2,B4}
			\move*{D3}
		\end{goban}
	\end{codeexample}
\end{key}

\begin{key}{/pgfgo/move/label/format/roman}
	
	Uses lowercase Roman numerals.
	
	\begin{codeexample}[]
		\pgfgoset{goban = {5, partial = {from = B2 to D4}}}
		
		\begin{goban}
			\stones[b]{C2,D3}
			\stone[w]{B4}
		\end{goban}
		
		\pgfgoset{move = {label = {format = roman}}}
		\begin{goban}
			\moves*{C2,B4}
			\move*{D3}
		\end{goban}
	\end{codeexample}
\end{key}

\begin{key}{/pgfgo/move/label/format/Roman}
	
	Uses uppercase Roman numerals.
	
	\begin{codeexample}[]
		\pgfgoset{goban = {5, partial = {from = B2 to D4}}}
		
		\begin{goban}
			\stones[b]{C2,D3}
			\stone[w]{B4}
		\end{goban}
		
		\pgfgoset{move = {label = {format = Roman}}}
		\begin{goban}
			\moves*{C2,B4}
			\move*{D3}
		\end{goban}
	\end{codeexample}
\end{key}

\begin{key}{/pgfgo/move/label/color = \meta{color}}
	
	Defines the color of the labels.
	
	\begin{codeexample}[]
		\pgfgoset{goban = {5, partial = {from = B2 to D4}}}
		
		\begin{goban}
			\moves*{C2,B4}
			\move*{D3}
		\end{goban}
		
		\begin{goban}
			\moves*[label = {color = red}]{C2,B4}
			\move*[label = {color = blue}]{D3}
		\end{goban}
	\end{codeexample}
\end{key}

\begin{key}{/pgfgo/move/label/forget}
	
	Prevents move labels from being stored in the board’s memory when using |remember|, allowing stones to persist without their numbers.
	
	\begin{codeexample}[]
		\pgfgoset{goban = {5, partial = {from = B2 to D4}}}
		
		\begin{goban}[new remember]
			\moves*[label = forget]{C2,B4,D3}
		\end{goban}
		
		\begin{goban}[resume]
			\moves*{C:D4}
		\end{goban}
	\end{codeexample}
\end{key}

\begin{key}{/pgfgo/move/label/remember}
	
	Ensures that labels are stored in memory, even if the environment uses |forget|, maintaining their continuity in subsequent diagrams.
\end{key}

\begin{key}{/pgfgo/move/forget}
	
	Indicates that moves should not be remembered in the board’s memory, removing them from future states when using |resume|.
\end{key}

\begin{key}{/pgfgo/move/remember}
	
	Ensures that moves are remembered, preserving them for reuse with |resume|.
\end{key}

\subsection{Prisoners}

\pgfgoname\ provides specific commands to represent prisoners, typically with reduced opacity to distinguish them from active stones.

\begin{command}{\prisoner\opt{\oarg{options}}\marg{intersection}}
	
	Draws a prisoner stone at a specific intersection, with options to customize its style.
\end{command}

\begin{command}{\prisoners\opt{\oarg{options}}\marg{intersection \opt{, intersection, ...}}}
	
	Places multiple prisoner stones in a single line, efficient for showing captured groups.
\end{command}

\begin{key}{/pgfgo/prisoner=\meta{stone options} (default fill opacity = 0.5)}
	
	Defines the characteristics of prisoner stones, with a default fill opacity of 0.5 to indicate their status. Other |\stone| options, such as color or size, can be applied.
	
\end{key}
	\section{Marks}

Marks are graphical elements placed at the intersections of the board to highlight specific positions, indicate points of interest, or number sequences without the need to add stones. \pgfgoname\ offers versatile commands and options to customize these marks, ranging from predefined shapes to custom labels.

\begin{command}{\mark\opt{\oarg{options}}\marg{intersection}}
	
	Draws an individual mark at the specified intersection (e.g., "C3"). Options allow defining its shape, color, and other attributes.
\end{command}

\begin{command}{\marks\opt{\oarg{options}}\marg{intersection \opt{, intersection, ...}}}
	
	Similar to |\mark|, but applies the same mark to multiple intersections using the coordinate syntax.
\end{command}

\begin{key}{/pgfgo/mark = \meta{options} (default cross)}
	
	Defines the type of mark to be placed, with a cross ("cross") as the default value if no other option is specified.
	
	\begin{codeexample}[]
		\pgfgoset{goban = {5, partial = {from = B2 to D4}}}
		
		\begin{goban}
			\stones[b]{C2,D3}
			\stone[w]{B4}
		\end{goban}
		
		\begin{goban}
			\stones[b]{C2,D3}
			\stone[w]{B4}
			\mark{C3}
		\end{goban}
	\end{codeexample}
\end{key}

\begin{key}{/pgfgo/mark/circle}
	
	Places a hollow circle as a mark.
	
	\begin{codeexample}[]
		\pgfgoset{goban = {5, partial = {from = B2 to D4}}}
		
		\begin{goban}
			\stones[b]{C2,D3}
			\stone[w]{B4}
		\end{goban}
		
		\begin{goban}
			\stones[b]{C2,D3}
			\stone[w]{B4}
			\mark[circle]{C3}
		\end{goban}
	\end{codeexample}
\end{key}

\begin{key}{/pgfgo/mark/triangle}
	
	Draws a hollow triangle.
	
	\begin{codeexample}[]
		\pgfgoset{goban = {5, partial = {from = B2 to D4}}}
		
		\begin{goban}
			\stones[b]{C2,D3}
			\stone[w]{B4}
		\end{goban}
		
		\begin{goban}
			\stones[b]{C2,D3}
			\stone[w]{B4}
			\mark[triangle]{C3}
		\end{goban}
	\end{codeexample}
\end{key}

\begin{key}{/pgfgo/mark/square}
	
	Places a hollow square.
	
	\begin{codeexample}[]
		\pgfgoset{goban = {5, partial = {from = B2 to D4}}}
		
		\begin{goban}
			\stones[b]{C2,D3}
			\stone[w]{B4}
		\end{goban}
		
		\begin{goban}
			\stones[b]{C2,D3}
			\stone[w]{B4}
			\mark[square]{C3}
		\end{goban}
	\end{codeexample}
\end{key}

\begin{key}{/pgfgo/mark/cross}
	
	Draws a cross.
	
	\begin{codeexample}[]
		\pgfgoset{goban = {5, partial = {from = B2 to D4}}}
		
		\begin{goban}
			\stones[b]{C2,D3}
			\stone[w]{B4}
		\end{goban}
		
		\begin{goban}
			\stones[b]{C2,D3}
			\stone[w]{B4}
			\mark[cross]{C3}
		\end{goban}
	\end{codeexample}
\end{key}

\begin{key}{/pgfgo/mark/filled square}
	
	Places a filled square.
	
	\begin{codeexample}[]
		\pgfgoset{goban = {5, partial = {from = B2 to D4}}}
		
		\begin{goban}
			\stones[b]{C2,D3}
			\stone[w]{B4}
		\end{goban}
		
		\begin{goban}
			\stones[b]{C2,D3}
			\stone[w]{B4}
			\mark[filled square]{C3}
		\end{goban}
	\end{codeexample}
\end{key}

\begin{key}{/pgfgo/mark/filled circle}
	
	Draws a filled circle.
	
	\begin{codeexample}[]
		\pgfgoset{goban = {5, partial = {from = B2 to D4}}}
		
		\begin{goban}
			\stones[b]{C2,D3}
			\stone[w]{B4}
		\end{goban}
		
		\begin{goban}
			\stones[b]{C2,D3}
			\stone[w]{B4}
			\mark[filled circle]{C3}
		\end{goban}
	\end{codeexample}
\end{key}

\begin{key}{/pgfgo/mark/t-circle}
	
	Places a small circle with a thicker border, specifically designed for marking territories. It can also be written |tcircle|.
	
	\begin{codeexample}[]
		\pgfgoset{goban = {5, partial = {from = B2 to D4}}}
		
		\begin{goban}
			\stones[b]{C2,D3}
			\stone[w]{B4}
		\end{goban}
		
		\begin{goban}
			\stones[b]{C2,D3}
			\stone[w]{B4}
			\mark[t-circle]{C3}
		\end{goban}
	\end{codeexample}
\end{key}

\begin{key}{/pgfgo/mark/t-square}
	
	Draws a small square with a thick border, similar to "t-circle," used to mark territories. It can also be written |tsquare|.
	
	\begin{codeexample}[]
		\pgfgoset{goban = {5, partial = {from = B2 to D4}}}
		
		\begin{goban}
			\stones[b]{C2,D3}
			\stone[w]{B4}
		\end{goban}
		
		\begin{goban}
			\stones[b]{C2,D3}
			\stone[w]{B4}
			\mark[t-square]{C3}
		\end{goban}
	\end{codeexample}
\end{key}

\begin{key}{/pgfgo/mark/hatched lines}
	
	Places a pattern of crossed (hatched) lines.
	
	\begin{codeexample}[]
		\pgfgoset{goban = {5, partial = {from = B2 to D4}}}
		
		\begin{goban}
			\stones[b]{C2,D3}
			\stone[w]{B4}
		\end{goban}
		
		\begin{goban}
			\stones[b]{C2,D3}
			\stone[w]{B4}
			\mark[hatched lines]{C3}
		\end{goban}
	\end{codeexample}
\end{key}

\begin{key}{/pgfgo/mark/sequence=\meta{options} (default from = 1)}
	Allows numbering a series of marks in order.
\end{key}

\begin{key}{/pgfgo/mark/sequence/use}
	Displays the |sequence| mark on the board. It is the same as |sequence = true|.
\end{key}

\begin{key}{/pgfgo/mark/sequence/false}
	Turns the sequence off (for example, one turned on with |\pgfgoset|).
\end{key}

\begin{key}{/pgfgo/mark/sequence/format = \meta{format}}
	Turns the sequence on and chooses the number format. It is the same as |sequence, label = {format = ...}| (see |label/format|).
\end{key}

\begin{key}{/pgfgo/mark/sequence/from=\meta{number}}
	
	Sets the starting number of a numbered mark sequence, allowing continuation from a specific point instead of 1.
	
	\begin{codeexample}[]
		\pgfgoset{goban = {5, partial = {from = B2 to D4}}}
		
		\begin{goban}
			\stones[b]{C2,D3}
			\stone[w]{B4}
		\end{goban}
		
		\begin{goban}
			\stones[b]{C2,D3}
			\stone[w]{B4}
			\marks[sequence = {use, from = 3}]{C3,B3,C4}
		\end{goban}
	\end{codeexample}
\end{key}

\begin{key}{/pgfgo/mark/custom=\meta{token list}}
	
	Allows defining a custom mark using text or symbols (e.g., "!" or "?") and even images.
	
	\begin{codeexample}[]
		\pgfgoset{goban = {5, partial = {from = B2 to D4}}}
		
		\begin{goban}
			\stones[b]{C2,D3}
			\stone[w]{B4}
		\end{goban}
		
		\begin{goban}
			\stones[b]{C2,D3}
			\stone[w]{B4}
			\mark[custom = {!}]{C3}
		\end{goban}
		
		\begin{goban}
			\stones[b]{C2,D3}
			\stone[w]{B4}
			\mark[custom = {\tikz{\node[rounded corners, draw, inner sep = 1em, red]{\dots};}}]{C3}
		\end{goban}
	\end{codeexample}
\end{key}

\begin{key}{/pgfgo/mark/line = \meta{options}}
	Controls the outline of the mark.
\end{key}

\begin{key}{/pgfgo/mark/line/color = \meta{color}}
	Defines the outline color of the mark.
	
	\begin{codeexample}[]
		\pgfgoset{goban = {5, partial = {from = B2 to D4}}}
		
		\begin{goban}
			\stones[b]{C2,D3}
			\stone[w]{B4}
		\end{goban}
		
		\begin{goban}
			\stones[b]{C2,D3}
			\stone[w]{B4}
			\mark[triangle, line = {color = red}]{C3}
		\end{goban}
	\end{codeexample}
\end{key}

\begin{key}{/pgfgo/mark/line/width = \meta{dimension}}
	
	Adjusts the thickness of the mark’s outline.
	
	\begin{codeexample}[]
		\pgfgoset{goban = {5, partial = {from = B2 to D4}}}
		
		\begin{goban}
			\stones[b]{C2,D3}
			\stone[w]{B4}
			\mark[triangle]{C3}
		\end{goban}
		
		\begin{goban}
			\stones[b]{C2,D3}
			\stone[w]{B4}
			\mark[triangle, line = {width = 0.1pt}]{C3}
		\end{goban}
	\end{codeexample}
\end{key}

\begin{key}{/pgfgo/mark/fill = \meta{options}}
	Controls the fill of the filled marks (|filled square|, |filled circle| and |hatched lines|).
\end{key}

\begin{key}{/pgfgo/mark/fill/true}
	Filled marks use |fill/color|. Without |fill = true|, they are filled in gray.
\end{key}

\begin{key}{/pgfgo/mark/fill/false}
	Filled marks are filled in gray again. It is the initial value.
\end{key}

\begin{key}{/pgfgo/mark/fill/color = \meta{color}}
	
	Sets the fill color of a mark. It is only used with |fill = true|.
	
	\begin{codeexample}[]
		\pgfgoset{goban = {5, partial = {from = B2 to D4}}}
		
		\begin{goban}
			\stones[b]{C2,D3}
			\stone[w]{B4}
		\end{goban}
		
		\begin{goban}
			\stones[b]{C2,D3}
			\stone[w]{B4}
			\mark[filled circle, fill = {true, color = red}]{C3}
		\end{goban}
	\end{codeexample}
\end{key}

\begin{key}{/pgfgo/mark/label = \meta{option}}
	Controls the appearance of labels for marks that have them.
\end{key}

\begin{key}{/pgfgo/mark/label/color = \meta{color}}
	
	Controls the color of the text or symbol in custom marks (e.g., with |custom|).
	
	\begin{codeexample}[]
		\pgfgoset{goban = {5, partial = {from = B2 to D4}}}
		
		\begin{goban}
			\stones[b]{C2,D3}
			\stone[w]{B4}
		\end{goban}
		
		\begin{goban}
			\stones[b]{C2,D3}
			\stone[w]{B4}
			\mark[custom = {!}, label = {color = red}]{C3}
		\end{goban}
	\end{codeexample}
\end{key}

\begin{key}{/pgfgo/mark/label/font = \meta{font commands}}
	Font of the labels of marks (|sequence| and |custom|).

	\begin{codeexample}[]
		\pgfgoset{goban = {5, partial = {from = B2 to D4}}}

		\begin{goban}
			\stones[b]{C2,D3}
			\stone[w]{B4}
			\marks[sequence, label = {font = \itshape}]{C3,B3,C4}
		\end{goban}
	\end{codeexample}
\end{key}

\begin{key}{/pgfgo/mark/label/format=\meta{options} (initially arabic)}
	
	Controls the numbering style of enumeration labels (e.g., |sequence|), with options such as Arabic numerals, letters, or Roman numerals.
\end{key}

\begin{key}{/pgfgo/mark/label/format/arabic}
	
	Uses Arabic numerals for sequences.
	
	\begin{codeexample}[]
		\pgfgoset{goban = {5, partial = {from = B2 to D4}}}
		
		\begin{goban}
			\stones[b]{C2,D3}
			\stone[w]{B4}
		\end{goban}
		
		\begin{goban}
			\stones[b]{C2,D3}
			\stone[w]{B4}
			\marks[label = {format = arabic}, sequence]{C3,B3,C4}
		\end{goban}
	\end{codeexample}
\end{key}

\begin{key}{/pgfgo/mark/label/format/alph}
	
	Uses lowercase letters for the sequence.
	
	\begin{codeexample}[]
		\pgfgoset{goban = {5, partial = {from = B2 to D4}}}
		
		\begin{goban}
			\stones[b]{C2,D3}
			\stone[w]{B4}
		\end{goban}
		
		\begin{goban}
			\stones[b]{C2,D3}
			\stone[w]{B4}
			\marks[label = {format = alph}, sequence]{C3,B3,C4}
		\end{goban}
	\end{codeexample}
\end{key}

\begin{key}{/pgfgo/mark/label/format/Alph}
	
	Uses uppercase letters.
	
	\begin{codeexample}[]
		\pgfgoset{goban = {5, partial = {from = B2 to D4}}}
		
		\begin{goban}
			\stones[b]{C2,D3}
			\stone[w]{B4}
		\end{goban}
		
		\begin{goban}
			\stones[b]{C2,D3}
			\stone[w]{B4}
			\marks[label = {format = Alph}, sequence]{C3,B3,C4}
		\end{goban}
	\end{codeexample}
\end{key}

\begin{key}{/pgfgo/mark/label/format/roman}
	
	Uses lowercase Roman numerals.
	
	\begin{codeexample}[]
		\pgfgoset{goban = {5, partial = {from = B2 to D4}}}
		
		\begin{goban}
			\stones[b]{C2,D3}
			\stone[w]{B4}
		\end{goban}
		
		\begin{goban}
			\stones[b]{C2,D3}
			\stone[w]{B4}
			\marks[label = {format = roman}, sequence]{C3,B3,C4}
		\end{goban}
	\end{codeexample}
\end{key}

\begin{key}{/pgfgo/mark/label/format/Roman}
	
	Uses uppercase Roman numerals.
	
	\begin{codeexample}[]
		\pgfgoset{goban = {5, partial = {from = B2 to D4}}}
		
		\begin{goban}
			\stones[b]{C2,D3}
			\stone[w]{B4}
		\end{goban}
		
		\begin{goban}
			\stones[b]{C2,D3}
			\stone[w]{B4}
			\marks[label = {format = Roman}, sequence]{C3,B3,C4}
		\end{goban}
	\end{codeexample}
\end{key}

\begin{key}{/pgfgo/mark/black}
	
	Applies the properties of the first player (black) to the mark.
	
	\begin{codeexample}[]
		\pgfgoset{goban = {5, partial = {from = B2 to D4}}}
		
		\begin{goban}
			\stones[b]{C2,D3}
			\stone[w]{B4}
			\mark[triangle]{C3}
		\end{goban}
		
		\begin{goban}
			\stones[b]{C2,D3}
			\stone[w]{B4}
			\mark[black, t-square]{C3}
		\end{goban}
	\end{codeexample}
\end{key}

\begin{key}{/pgfgo/mark/white}
	
	Applies the properties of the second player (white) to the mark.
\end{key}

\begin{key}{/pgfgo/mark/neutral}
	
	Keeps the mark unassociated with any player.
\end{key}

\begin{key}{/pgfgo/mark/b}
	Short for |black|.
\end{key}

\begin{key}{/pgfgo/mark/w}
	Short for |white|.
\end{key}

\begin{key}{/pgfgo/mark/n}
	Short for |neutral|.
\end{key}

\begin{key}{/pgfgo/mark/forget}
	This mark is not stored in the board’s memory, even if the |goban| uses |remember|.
\end{key}

\begin{key}{/pgfgo/mark/remember}
	This mark is stored in the board’s memory, even if the |goban| uses |forget|.
\end{key}
	\section{The game: captures, notes and territory}
\label{sec:juego}

With the Lua kernel, a |goban| can follow the rules of the game: remove the captured stones, note the moves that fall on points already played and count the territory.

	\subsection{Captures}
	\label{subsec:capturas}

	With the Lua kernel (p.~\pageref{subsec:kernel}), \pgfgoname\ can keep track of the stones on the board and remove those left without liberties, as in a game.

	\begin{key}{/pgfgo/goban/captures=\meta{remove, fade or false} (default remove, initially false)}

		Turns captures on. When a move leaves an opposing group without liberties, that group is removed from the board. If it leaves its own group without liberties (a suicide), that group is removed. Stones put with |\stone| and |\stones| also count; neutral stones and prisoners do not. A move outside a partial board also captures. A move on an occupied point gives a warning, because it is not possible in a game. With |kernel = expl3|, this \emph{key} gives an error.

		\begin{codeexample}[]
			\pgfgoset{goban = {5, partial = A1:C3}}

			\begin{goban}
				\moves{B1,A1,A2}
			\end{goban}

			\begin{goban}[captures]
				\moves{B1,A1,A2} % 3 captures 2.
			\end{goban}
		\end{codeexample}

		Since \pgfname\ cannot erase what it has already drawn, with |captures| the stones are drawn when the |goban| ends. So, if the style is changed with |\pgfgoset| inside a |goban| with captures, the change does not reach the stones; the options of each command do.
	\end{key}

	\begin{stylekey}{/pgfgo/goban/captures/remove}

		The captured stones disappear. It is the same as |captures| or |captures = true|.
	\end{stylekey}

	\begin{stylekey}{/pgfgo/goban/captures/fade}

		The captured stones stay on the board, drawn as prisoners (with the style of |\prisoner| and of |/pgfgo/prisoner|).

		\begin{codeexample}[]
			\pgfgoset{goban = {5, partial = A1:E3}}

			\begin{goban}[captures = fade]
				\stones[w]{C1,D1}
				\stones[b]{B1,C2,D2}
				\moves{E1}
			\end{goban}
		\end{codeexample}
	\end{stylekey}

	\begin{stylekey}{/pgfgo/goban/captures/false}

		Turns captures off. It is the initial value.
	\end{stylekey}

	\begin{key}{/pgfgo/goban/captured marks=\meta{keep or remove} (initially keep)}

		Decides what happens to the marks on a stone when it is captured. It only has an effect with |captures|. Marks on empty points do not change.

		\begin{codeexample}[]
			\pgfgoset{goban = {5, partial = A1:C3, captures}}

			\begin{goban}
				\stones[w]{A1}\stones[b]{B1}
				\marks[triangle]{A1}
				\moves{A2}
			\end{goban}

			\begin{goban}[captured marks = remove]
				\stones[w]{A1}\stones[b]{B1}
				\marks[triangle]{A1}
				\moves{A2} % the triangle goes with the white stone.
			\end{goban}
		\end{codeexample}
	\end{key}

	\begin{stylekey}{/pgfgo/goban/captured marks/keep}

		The marks stay on the empty point. It is the initial value.
	\end{stylekey}

	\begin{stylekey}{/pgfgo/goban/captured marks/remove}

		The marks go with the stone, also with |captures = fade|. Sequences keep their numbering.
	\end{stylekey}

	\begin{command}{\pgfgoprisoners\marg{player}}
		Writes how many prisoners the player (|black| or |b|, |white| or |w|) has in the last |goban| with |captures|: the opposing stones captured by its moves. With another player, it gives an error and writes 0.

		\begin{codeexample}[]
			\pgfgoset{goban = {5, partial = A1:C3}}

			\begin{goban}[captures]
				\moves{B1,A1,A2}
			\end{goban}

			Black: \pgfgoprisoners{black}. White: \pgfgoprisoners{white}.
		\end{codeexample}
	\end{command}

	\subsection{Moves on points already played}
	\label{subsec:notas}

	In numbered diagrams, when a move falls on a point where there was already a move (for example, when taking back a ko), books do not draw it: the point still shows the first stone and below the diagram it is noted, for example, ``3 at 1''.

	\begin{key}{/pgfgo/goban/notes=\meta{boolean} (default true, initially false)}

		A numbered move on a point where there was already a numbered move of the same |goban| is not drawn: the point still shows the first one, even if it was captured, and ``N at M'' is noted. Only the numbered moves (|\moves| without |numbers = false|) that are seen on the board count. It works with both kernels.

		\begin{codeexample}[]
			\pgfgoset{goban = {5, partial = A2:E4}}

			\begin{goban}[notes]
				\stones[b]{B3,C2,C4}
				\stones[w]{D2,D4,E3}
				\moves[w]{C3}
				\moves{D3}
				\moves[from = 3, w]{C3}
			\end{goban}

			\pgfgomovenotes
		\end{codeexample}
	\end{key}

	\begin{command}{\pgfgomovenotes\opt{\oarg{options}}}
		Writes the notes of the last |goban|, separated by commas. If there are no notes, it writes nothing. The options are those of |/pgfgo/notes|.
	\end{command}

	\begin{key}{/pgfgo/notes = \meta{options}}
		Options of |\pgfgomovenotes| for the whole document, for example |\pgfgoset{notes = {at = en}}|.
	\end{key}

	\begin{key}{/pgfgo/notes/at = \meta{text} (initially at or its translation)}
		Text between the two numbers.
	\end{key}

	\begin{key}{/pgfgo/notes/pass = \meta{text} (initially pass or its translation)}
		A pass of an SGF file (p.~\pageref{subsec:sgf}) counts, but it is not drawn. With |notes|, if it would be numbered, ``N: pass'' is noted, in order with the other notes. The passes that end the game are not noted. This key is the text after the number.
	\end{key}

	\begin{key}{/pgfgo/notes/separator = \meta{text} (initially a comma and a space)}
		Text between the notes.

		\begin{codeexample}[]
			\begin{goban}[5, partial = A2:E4, notes]
				\stones[b]{B3,C2,C4}
				\stones[w]{D2,D4,E3}
				\moves[w]{C3}
				\moves{D3}
				\moves[from = 3, w]{C3}
				\moves{D3}
			\end{goban}

			\pgfgomovenotes[separator = {; }]
		\end{codeexample}
	\end{key}

	\subsection{Territory}
	\label{subsec:territorio}

	At the end of the game, the territory can be marked by hand with |\marks| (p.~\pageref{sec:marcas}) or with |\territory|, which also counts it.

	\begin{command}{\territory\opt{\oarg{options}}\marg{polygon}}
		Marks as territory the points of the polygon that has those vertices (in order), including its border. It does not mark the points that have a living stone: prisoners count as territory. The vertices are written as in |\stones|; a single one marks a point, and two, a segment. It works with both kernels.

		The territory is drawn when the |goban| ends. Each point keeps the last |\territory| that marked it, with the player and the options of that |\territory| (so |\territory[n]{A1}| after another |\territory| makes A1 neutral), and it is neither drawn nor counted if at the end it has a living stone.

		\begin{codeexample}[width=4.75cm]
			\begin{goban}[label = {at = {above, left}}]
		\stones[b] {A6|-B5|-D1|-E2, C4, E9-|F4, D5-|G7--D7, H6, H8, I7}
		\stones[w] {E8-|C5, B:E6, A7, B8, D9, I6|- H5, H2:3, F1|-C3, C2, D:E4, G4}
		\prisoners[b]{H1-|I2}
		\prisoner[w] {A1}

				\territory[b]{A1,A5,B5,B1}
				\territory[b]{F9,F7,G7,G6,I6,I9}
				\territory[w]{F1,F4,G4,G5,I5,I1}
				\territory[w]{A9,A6,C6,C8,D8,D9}
			\end{goban}
		\end{codeexample}
	\end{command}

	\begin{command}{\territory\opt{\oarg{options}}\marg{auto}}
		With |auto| instead of a polygon, it finds the territory by itself: each region of points without living stones belongs to the color that surrounds it. If both colors touch it (dame, seki) or neither does, it is neutral. Prisoners count as dead stones. Without a player, it marks the territory of both; with |b| or |w|, only that of that player. Only with the Lua kernel (p.~\pageref{subsec:kernel}).

		\begin{codeexample}[width=4.75cm]
			\begin{goban}[label = {at = {above, left}}]
		\stones[b] {A6|-B5|-D1|-E2, C4, E9-|F4, D5-|G7--D7, H6, H8, I7}
		\stones[w] {E8-|C5, B:E6, A7, B8, D9, I6|- H5, H2:3, F1|-C3, C2, D:E4, G4}
		\prisoners[b]{H1-|I2}
		\prisoner[w] {A1}

				\territory[dame]{auto}
			\end{goban}
		\end{codeexample}

		|auto| only looks at which color each region touches; it cannot read the position. So:
		\begin{itemize}
			\item It is for finished games. In an open position, almost the whole board touches both colors and is neutral; with stones of a single color, the whole board is theirs. If a border has a gap, the territory inside joins the dame and is neutral too.
			\item It does not know which stones are dead (sometimes the players do not either): they have to be put with |\prisoners|. A dead stone put with |\stones| counts as alive, and the region it touches is neutral.
			\item The eyes of a group in seki count as territory (under Japanese rules they do not). They can be made neutral afterwards, with a |\territory[n]| for each one.
			\item Neutral stones (without a player) occupy their point, but belong to nobody: a region that only touches neutral stones is neutral.
			\item It finds the regions with the stones there are when |\territory| appears, so it is best put at the end of the |goban|.
			\item With |partial = auto|, the board grows to show the whole territory (and the dame, with |dame|).
		\end{itemize}
	\end{command}

	\begin{key}{/pgfgo/territory/black}
		The territory is Black's. It can also be written |b|.
	\end{key}

	\begin{key}{/pgfgo/territory/white}
		The territory is White's. It can also be written |w|.
	\end{key}

	\begin{key}{/pgfgo/territory/neutral}
		Neutral points. It is what a polygon without a player marks. It can also be written |n|.
	\end{key}

	\begin{key}{/pgfgo/territory/b}
		Short for |black|.
	\end{key}

	\begin{key}{/pgfgo/territory/w}
		Short for |white|.
	\end{key}

	\begin{key}{/pgfgo/territory/n}
		Short for |neutral|.
	\end{key}

	\begin{key}{/pgfgo/territory/draw=\meta{marks, fill or region} (initially marks)}
		With |marks|, each point gets the territory mark of the player. With |fill|, the region is filled under the stones.

		\begin{codeexample}[width=4.75cm]
			\begin{goban}[label = {at = {above, left}}]
		\stones[b] {A6|-B5|-D1|-E2, C4, E9-|F4, D5-|G7--D7, H6, H8, I7}
		\stones[w] {E8-|C5, B:E6, A7, B8, D9, I6|- H5, H2:3, F1|-C3, C2, D:E4, G4}
		\prisoners[b]{H1-|I2}
		\prisoner[w] {A1}

				\territory[draw = fill]{auto}
			\end{goban}
		\end{codeexample}

		With |region|, the territory of each color is a single polygon: each square of the grid that has a corner in the territory is painted (dead stones count as territory). Its sides run along the lines of the board, its vertices are always on intersections, and it reaches the edge of the board. Two zones that touch at a single corner stay separate. All the |\territory| of a color in the same |goban| (also the territory |TB| and |TW| of |\sgf|) form a single region, which is drawn when the |goban| ends with the style of the last one. Neutral points are filled one by one, as with |fill|. The outline and its vertices can be drawn on top, with |line| and |vertices|, and |\pgfgoregion| gives its properties.

		\begin{codeexample}[width=4.75cm]
			\begin{goban}[label = {at = {above, left}}]
		\stones[b] {A6|-B5|-D1|-E2, C4, E9-|F4, D5-|G7--D7, H6, H8, I7}
		\stones[w] {E8-|C5, B:E6, A7, B8, D9, I6|- H5, H2:3, F1|-C3, C2, D:E4, G4}
		\prisoners[b]{H1-|I2}
		\prisoner[w] {A1}

				\territory[draw = region]{auto}
			\end{goban}
		\end{codeexample}
	\end{key}

	\begin{key}{/pgfgo/territory/dame=\meta{boolean} (default true, initially false)}
		With |auto|, also marks the neutral points.
	\end{key}

	\begin{key}{/pgfgo/territory/forget}
		|resume| does not bring this territory, even if |territory| is asked for (p.~\pageref{subsec:recordar}).
	\end{key}

	The other options of |\territory| change the style of the player only for that command, like |mark| and |fill| in |\territory[w, mark = {cross}]{auto}|. Without a player, they change that of the three. The style of each player is set with |/pgfgo/players|:

	\begin{key}{/pgfgo/players/black/territory=\meta{territory options}}
		Style of the territory of each player (|black|, |white| or |neutral|).
	\end{key}

	\begin{key}{/pgfgo/players/black/territory/mark=\meta{mark options} (initially t-circle)}
		Mark of each point, with |draw = marks|. It is |t-circle| for Black and |t-square| for White and the neutral points.

		\begin{codeexample}[]
			\begin{goban}[5]
				\stones[b]{C1:5}
				\stones[w]{D1:5}
				\territory[b, mark = {t-square}]{A1, A5, B5, B1}
				\territory[w]{E1, E5}
			\end{goban}
		\end{codeexample}
	\end{key}

	\begin{key}{/pgfgo/players/black/territory/fill=\meta{fill options} (default true)}
		Fill, with |draw = fill| and |draw = region|.
	\end{key}

	\begin{key}{/pgfgo/players/black/territory/fill/true}
		Draws the fill. It is the initial value.
	\end{key}

	\begin{key}{/pgfgo/players/black/territory/fill/false}
		Does not draw the fill: with |draw = region| only the outline or the vertices remain.
	\end{key}

	\begin{key}{/pgfgo/players/black/territory/fill/color=\meta{color} (initially black)}
		Color of the fill: |black| for Black, |white!80!black| for White and |gray| for the neutral points. With |none|, no color.
	\end{key}

	\begin{key}{/pgfgo/players/black/territory/fill/opacity=\meta{number} (initially 0.35)}
		Opacity of the fill: |0.35| for Black and the neutral points, |0.7| for White.

		\begin{codeexample}[]
			\begin{goban}[5]
				\stones[b]{C1:5}
				\stones[w]{D1:5}
				\territory[b, draw = fill,
					fill = {color = blue, opacity = 0.5}]{A1, A5, B5, B1}
				\territory[w, draw = fill]{E1, E5}
			\end{goban}
		\end{codeexample}
	\end{key}

	\begin{key}{/pgfgo/players/black/territory/fill/pattern=\meta{options}}
		A pattern over the fill color, as |stone/fill/pattern| (p.~\pageref{sec:piedras}). With |fill/color = none|, only the pattern is left.

		\begin{codeexample}[]
			\begin{goban}[5]
				\stones[b]{C1:5}
				\stones[w]{D1:5}
				\territory[b, draw = fill, fill = {color = none,
					pattern = {name = north east lines}}]{A1, A5, B5, B1}
				\territory[w, draw = region, fill = {color = none,
					pattern = {name = north west lines}}]{E1:5}
			\end{goban}
		\end{codeexample}
	\end{key}

	\begin{key}{/pgfgo/players/black/territory/fill/pattern/name = \meta{pattern}}
		The pattern: any of the Ti\emph{k}Z libraries |patterns| and |patterns.meta|, which \pgfgoname\ loads. One with options goes in braces: |name = {Lines[angle = 45, distance = 2pt]}|. Without |name| there is no pattern.
	\end{key}

	\begin{key}{/pgfgo/players/black/territory/fill/pattern/color = \meta{color} (initially black)}
		Color of the pattern.
	\end{key}

	\begin{key}{/pgfgo/players/black/territory/fill/pattern/false}
		No pattern.
	\end{key}

	\begin{key}{/pgfgo/players/black/territory/line=\meta{line options} (default true)}
		Outline of the region, with |draw = region|. Alone or with |true|, it draws the outline; options alone do not turn it on.

		\begin{codeexample}[width=4.75cm]
			\begin{goban}[label = {at = {above, left}}]
		\stones[b] {A6|-B5|-D1|-E2, C4, E9-|F4, D5-|G7--D7, H6, H8, I7}
		\stones[w] {E8-|C5, B:E6, A7, B8, D9, I6|- H5, H2:3, F1|-C3, C2, D:E4, G4}
		\prisoners[b]{H1-|I2}
		\prisoner[w] {A1}

				\territory[b, draw = region, line,
					vertices]{auto}
				\territory[w, draw = region,
					line = {true, color = red}]{auto}
			\end{goban}
		\end{codeexample}
	\end{key}

	\begin{key}{/pgfgo/players/black/territory/line/true}
		Draws the outline.
	\end{key}

	\begin{key}{/pgfgo/players/black/territory/line/false}
		Does not draw the outline. It is the initial value.
	\end{key}

	\begin{key}{/pgfgo/players/black/territory/line/color=\meta{color} (initially black)}
		Color of the outline: |black| for Black and |white!40!black| for White.
	\end{key}

	\begin{key}{/pgfgo/players/black/territory/line/width=\meta{dimension} (initially 1.2pt)}
		Width of the outline.
	\end{key}

	\begin{key}{/pgfgo/players/black/territory/line/opacity=\meta{number} (initially 1)}
		Opacity of the outline and of the vertices.
	\end{key}

	\begin{key}{/pgfgo/players/black/territory/line/above=\meta{boolean} (default true, initially true)}
		With |true|, the outline and the vertices are drawn on top of the stones and the marks. With |false|, under the stones, together with the fill.

		\begin{codeexample}[]
			\begin{goban}[5]
				\stones[b]{C1:5}
				\stones[w]{D1:5}
				\territory[b, draw = region, fill = false,
					line = {true, color = red, width = 2pt, opacity = 0.7}]{A1, A5, B5, B1}
			\end{goban}

			\begin{goban}[5]
				\stones[b]{C1:5}
				\stones[w]{D1:5}
				\territory[b, draw = region, fill = false,
					line = {true, color = red, width = 2pt, opacity = 0.7, above = false}]{A1, A5, B5, B1}
			\end{goban}
		\end{codeexample}
	\end{key}

	\begin{key}{/pgfgo/players/black/territory/vertices=\meta{options} (default true)}
		Points on the vertices of the outline, with |draw = region|, like the star points. They are drawn in the same layer as the outline. Alone or with |true|, it draws them; options alone do not turn them on.
	\end{key}

	\begin{key}{/pgfgo/players/black/territory/vertices/true}
		Draws the vertices.
	\end{key}

	\begin{key}{/pgfgo/players/black/territory/vertices/false}
		Does not draw the vertices. It is the initial value.
	\end{key}

	\begin{key}{/pgfgo/players/black/territory/vertices/radius=\meta{dimension} (initially 2pt)}
		Radius of the points.
	\end{key}

	\begin{key}{/pgfgo/players/black/territory/vertices/color=\meta{color}}
		Color of the points. Without this \emph{key}, they take the color of the outline.

		\begin{codeexample}[]
			\begin{goban}[5]
				\stones[b]{C1:5}
				\stones[w]{D1:5}
				\territory[b, draw = region,
					vertices = {true, radius = 3pt, color = red}]{A1, A5, B5, B1}
			\end{goban}
		\end{codeexample}
	\end{key}

	\begin{command}{\pgfgoscore\opt{\oarg{options}}\marg{player}}
		Writes the score of the player (|black| or |b|, |white| or |w|) in the last |goban|, under Japanese rules: the points of its territory, plus the dead opposing stones in it, plus the stones it captured (with |captures|), plus the |komi| for White. A point marked twice counts only once. With |neutral| or |n|, it writes the number of neutral points marked. With another player, it gives an error and writes 0.

		\begin{codeexample}[code only]
			Black: \pgfgoscore{black}. White: \pgfgoscore[komi = 6.5]{white}.
		\end{codeexample}

		With the board above, it writes ``Black: 11. White: 23.5''.
	\end{command}

	\begin{key}{/pgfgo/score/komi=\meta{number} (initially 0)}
		Komi added to White's score. For the whole document: |\pgfgoset{score = {komi = 6.5}}|.

		\begin{codeexample}[]
			\begin{goban}[5]
				\stones[b]{C1:5}
				\stones[w]{D1:5}
				\territory{auto}
			\end{goban}

			Black: \pgfgoscore{b}. White: \pgfgoscore[komi = 6.5]{w}.
		\end{codeexample}
	\end{key}

	\begin{command}{\pgfgoregion\marg{player}\marg{property}}
		Writes a property of the region (|draw = region|) of the player (|black| or |b|, |white| or |w|) in the current |goban| or in the last one. It is expandable, so it can be used in a computation or inside the coordinates of |\marks|. The properties are those of the whole outline, even if the board is partial. Without a region, the numbers are 0. The properties are:
		\begin{itemize}
			\item |polygons|: the number of outlines, outer ones and holes.
			\item |outlines|: the number of outer outlines, one for each separate zone.
			\item |holes|: the number of holes.
			\item |vertices|: the number of vertices of all the outlines.
			\item |area|: the number of painted squares.
			\item |perimeter|: the length of all the outlines, in squares.
			\item |polygon| \meta{i}: the vertices of outline \meta{i}, as coordinates separated by commas (for example, |A1,D1,D4,A4|). Empty if there are not that many outlines.
		\end{itemize}
		Each outline starts at its lowest vertex and, among those, at the leftmost one. Outer outlines go counterclockwise and holes, clockwise. The outlines are sorted by that first vertex, from bottom to top.

		\begin{codeexample}[width=4.75cm]
			\begin{goban}[label = {at = {above, left}}]
				\stones[w]{D4:F6}
				\territory[b, draw = region, line]{B2,B8,H8,H2}
				\marks[triangle]{\pgfgoregion{b}{polygon 2}}
			\end{goban}

			Outlines: \pgfgoregion{b}{polygons}
			(\pgfgoregion{b}{holes} hole).
			Squares: \pgfgoregion{b}{area}.
		\end{codeexample}
	\end{command}

	\section{Remembering boards and views}
\label{sec:memoria}

	\subsection{Remembering a board}
	\label{subsec:recordar}
	A named |goban| keeps everything it draws (stones, moves, marks, arrows, lines and territory, also what comes from |\sgf|), together with its options. Another |goban| can start with that board as it was left, with |resume|, and |\pgfgoview| shows it in the text (p.~\pageref{subsec:vistas}).

	\begin{codeexample}[]
		\pgfgoset{goban = {5, partial = B2:D4}}

		\begin{goban}[name = board1]
			\stones[b]{C2,D3}
			\moves[w]{B4}
		\end{goban}

		\begin{goban}[resume = board1]
			\moves{C3} % Goes on with 2 (Black).
		\end{goban}

		\begin{goban}[resume = {board1, numbers}]
			\moves{C3}
		\end{goban}
	\end{codeexample}

	\begin{key}{/pgfgo/goban/name=\meta{name}}
		Keeps with that name everything drawn in this |goban|, together with its options. The name is free (for example, |gob:board1|). Another |goban| with the same name replaces it.
	\end{key}

	\begin{key}{/pgfgo/goban/resume=\meta{name, options}}
		Starts the |goban| with the board that has that name, as it was at the end (with the captures resolved: if that |goban| used |captures|, this one does too). The numbering and the turn go on from there. By default it brings only the stones, without numbers, marks or territory; the options ask for the rest. The size is that of that board, unless this |goban| gives its own (|size|, a number, |sgf size| or |sgf view|; one set before with |\pgfgoset| does not count); the other options of the board (|partial|, coordinates) are given in each |goban|. If there is no |goban| with that name, an error is given.

		With |name| and |resume| in the same |goban| (in any order), the new name keeps what it brought, as it was drawn, plus its own, and the board of |resume| does not change: that is how variations are built. With the same name in both, the |goban| keeps adding to that board.

		\begin{codeexample}[]
			\pgfgoset{goban = {5, partial = B2:D4}}

			\begin{goban}[name = base]
				\stones[b]{C2,D3}
				\stones[w]{C3}
			\end{goban}

			\begin{goban}[name = variation1, resume = base]
				\moves[w]{D2}
			\end{goban}

			\begin{goban}[name = variation2, resume = base]
				\moves[w]{B3}
			\end{goban}

			\begin{goban}[resume = {variation1, numbers}]
				\moves{B4}
			\end{goban}
		\end{codeexample}
	\end{key}

	\begin{key}{/pgfgo/goban/resume/numbers=\meta{boolean} (default true, initially false)}
		Also brings the numbers of the moves.
	\end{key}

	\begin{key}{/pgfgo/goban/resume/marks=\meta{boolean} (default true, initially false)}
		Also brings the marks, the arrows and the lines.
	\end{key}

	\begin{key}{/pgfgo/goban/resume/territory=\meta{boolean} (default true, initially false)}
		Also brings the territory.

		\begin{codeexample}[]
			\pgfgoset{goban = {5, partial = B2:D4}}

			\begin{goban}[name = saved]
				\stones[b]{C2:4}
				\moves[w]{D3}
				\marks[triangle]{B3}
				\territory[b]{B2, B4}
			\end{goban}

			\begin{goban}[resume = saved]
			\end{goban}

			\begin{goban}[resume = {saved, numbers, marks, territory}]
			\end{goban}
		\end{codeexample}
	\end{key}

	So that |resume| does not bring a stone, a move, a mark, an arrow or a territory, it is drawn with the option |forget| (|\stone[forget]{C3}|, |\marks[forget, circle]{A1}|, \ldots); with |label = forget|, a move comes without its number even if |numbers| is asked for. What is forgotten stays kept: |\pgfgoview| shows it with |forgotten|.

	\subsection{Views in the text}
	\label{subsec:vistas}
	A named |goban| keeps everything it draws. Afterwards, |\pgfgoview| shows part of that board, as it was at the end (with the captures resolved), in the middle of a paragraph: a single stone or a box.

	\begin{codeexample}[]
		\begin{goban}[5, name = board1, captures]
			\stones[b]{B3, C2, C4}
			\stones[w]{D2, D4, E3}
			\moves[w]{C3, D3}
		\end{goban}

		Black at D3 \pgfgoview{board1}{D3} captures White at C3 \pgfgoview{board1}{C3}.
		The ko: \pgfgoview{board1}{B2:E4}.
	\end{codeexample}

	\begin{command}{\pgfgoview\opt{\oarg{options}}\marg{name}\marg{part}}
		Draws the intersection or the rectangle (written as in |partial|) of the |goban| with that name, as it was at the end. By default everything the |goban| drew in that part is seen (stones, numbers of the moves, marks, arrows and territory) with the background and the coordinates of that |goban|, at the size of the text and centered at the height of the lowercase letters. The options are those of |/pgfgo/view|, which remove the marks, the numbers or the territory, or add what was forgotten, and also any option of |goban|, which is applied after those of the named |goban|: that way the background, the coordinates, the grid or the star points are removed, or the style of the stones, the marks and the board is changed (|players|, |stone|, |mark|, |grid|, |background|, \ldots).

		The view uses the style that |goban| had when it was drawn (the profile, what was set with |\pgfgoset| and its options), even if it changes later. On top go the profile of the views (|view/profile|), |\pgfgoset{view = ...}| and the options of the view.

		The |goban| can come after the view (p.~\pageref{subsec:vistas-antes}) or from another document (p.~\pageref{subsec:exportar}).

		A view changes nothing of what the last |goban| keeps: |\pgfgoscore|, |\pgfgoprisoners| and |\pgfgomovenotes| still speak of it, and what each named |goban| keeps stays the same.

		\begin{codeexample}[]
			\begin{goban}[5, name = board2]
				\stones[b]{B2, C3}
				\moves[w]{C2, D3}
				\marks[b, triangle]{C3}
			\end{goban}

			\pgfgoview{board2}{B2:D3}
			\pgfgoview[marks = false, numbers = false]{board2}{B2:D3}
			\pgfgoview[label = false, grid = false]{board2}{B2:D3}
			\pgfgoview[background = false]{board2}{B2:D3}
		\end{codeexample}
	\end{command}

	\begin{key}{/pgfgo/view=\meta{options}}
		Options for every view, with |\pgfgoset|. The options of |goban| that are not of |/pgfgo/view| are added to those the views bring: |scale = 0.8, background = {line = false}| (the background without border, even if the original has one).

		\begin{codeexample}[]
			\pgfgoset{view = {background = false, scale = 1}}
			Without background and larger: \pgfgoview{board2}{C3}.
		\end{codeexample}
	\end{key}

	\begin{key}{/pgfgo/view/marks=\meta{boolean} (default true, initially true)}
		With |false|, does not draw the marks.
	\end{key}

	\begin{key}{/pgfgo/view/arrows=\meta{boolean} (default true, initially true)}
		With |false|, does not draw the arrows (|\arrow| and |AR| of |\sgf|).
	\end{key}

	\begin{key}{/pgfgo/view/lines=\meta{boolean} (default true, initially true)}
		With |false|, does not draw the lines (|\line| and |LN| of |\sgf|).

		\begin{codeexample}[]
			\begin{goban}[5, name = board-arrows]
				\stones[b]{B2}\stones[w]{D4}
				\arrow{B2}{D4}
				\line{B4}{D2}
				\marks[triangle]{C3}
			\end{goban}

			Without arrows: \pgfgoview[arrows = false]{board-arrows}{A1:E5}.
			Without lines or marks: \pgfgoview[lines = false, marks = false]{board-arrows}{A1:E5}.
		\end{codeexample}
	\end{key}

	\begin{key}{/pgfgo/view/numbers=\meta{boolean} (default true, initially true)}
		With |false|, does not draw the numbers of the moves.
	\end{key}

	\begin{key}{/pgfgo/view/territory=\meta{boolean} (default true, initially true)}
		With |false|, does not draw the territory.
	\end{key}

	\begin{key}{/pgfgo/view/forgotten=\meta{boolean} (default true, initially false)}
		Also draws what was drawn with |forget| (p.~\pageref{subsec:recordar}).

		\begin{codeexample}[]
			\begin{goban}[5, name = board3]
				\stones[b]{B2, C3}
				\stone[w, forget]{D3}
			\end{goban}

			\pgfgoview{board3}{B2:D3}
			\pgfgoview[forgotten]{board3}{B2:D3}
		\end{codeexample}
	\end{key}

	\begin{key}{/pgfgo/view/raise=\meta{dimension} (initially 0.5ex)}
		At what height above the baseline of the text the center of the view is.

		\begin{codeexample}[]
			Text \pgfgoview[raise = 0pt]{board2}{C3}
			text \pgfgoview{board2}{C3}
			text \pgfgoview[raise = 1.5ex]{board2}{C3}.
		\end{codeexample}
	\end{key}

	\begin{key}{/pgfgo/view/profile=\meta{name} (initially false)}
		A profile for the views (p.~\pageref{subsec:perfiles}), for example |simple|. It goes after the style of the named |goban| and before the other options of the view. With |false|, the view uses the style of the |goban| as it is.

		\begin{codeexample}[]
			\pgfgoview{board1}{B2:E4}
			\pgfgoview[profile = simple]{board1}{B2:E4}
			\pgfgoview[profile = wain]{board1}{B2:E4}
		\end{codeexample}
	\end{key}

	\subsection{A stone in the text}
	\label{subsec:piedra-texto}

	\begin{command}{\pgfgostone\opt{\oarg{options}}}
		Draws a single stone in the text, without a board. It uses the style of that moment and, on top, the one of the views: |view/profile|, |view/raise| and the |goban| options of |\pgfgoset{view = ...}|, such as |scale|. So it has the same size and height as the view of one intersection. The options are those of |\stone| (|b|, |w|, |n|, |prisoner|, |mark| and the style of the stone) and |number|. Without |b| or |w|, the stone is neutral.

		It changes nothing of what the last |goban| keeps.

		\begin{codeexample}[]
			The black stone \pgfgostone[b, mark = triangle] cuts and
			the white stone \pgfgostone[w, mark = {custom = A}] answers.

			\pgfgostone[b] \pgfgostone[w] \pgfgostone[n]
			\pgfgostone[b, prisoner] \pgfgostone[w, mark = cross]
		\end{codeexample}
	\end{command}

	\begin{key}{/pgfgo/pgfgostone/number=\meta{number}}
		The stone has the number, with the style of the moves of its color (|move/label|). A neutral stone has no number: it is drawn without it, with a warning.

		\begin{codeexample}[]
			Black \pgfgostone[b, number = 5] cuts,
			white \pgfgostone[w, number = 6] ataris.
		\end{codeexample}
	\end{key}

	\subsection{Views before the board}
	\label{subsec:vistas-antes}
	A view can come before the |goban| it shows, like a reference with |\ref|: the |goban| is kept in the |.aux| file and the view uses it in the next compilation. The first time the view is |??| and there is a warning; if the |goban| changes afterwards, another warning asks to compile again. If two |goban| environments have the same name, a view that comes before both shows the first one.

	\begin{codeexample}[]
		Further down, the white stone at C3 \pgfgoview{board4}{C3} is in atari.
		\begin{goban}[5, name = board4]
			\stones[b]{B3, C2, C4}
			\stones[w]{C3}
			\marks[sequence]{D3}
		\end{goban}
	\end{codeexample}

	\subsection{Using a board in another document}
	\label{subsec:exportar}
	A named |goban| can be saved in a file and used in another document, with its stones, numbers, marks, territory and style.

	\begin{key}{/pgfgo/goban/export=\meta{file}}
		When it ends, the |goban| is written to that file (if it has no extension, |.pgfgo|), for |\pgfgoimport|. It needs |name|. For security, \TeX\ only writes inside the folder of the document: the file goes with a relative name, such as |boards/ko|.

		\begin{codeexample}[code only]
			\begin{goban}[19, name = ko, export = ko]
				...
			\end{goban}
		\end{codeexample}
	\end{key}

	\begin{command}{\pgfgoimport\opt{\oarg{options}}\marg{file}}
		Brings the |goban| of an |export| file (without an extension, |.pgfgo|), as if it had been drawn in this document: it can be used with |\pgfgoview| and with |resume|. It can be used in the preamble.

		\begin{codeexample}[code only]
			\pgfgoimport{ko}
			\pgfgoimport[name = other ko]{ko}
			...
			The ko of \pgfgoview{ko}{Q16:S18}.
			\begin{goban}[resume = {other ko, marks}]
				\moves[numbers = false]{R17}
			\end{goban}
		\end{codeexample}

		A mark defined with |\pgfgodefinemark| has to be defined in this document too; otherwise, there is a warning.
	\end{command}

	\begin{key}{/pgfgo/import/name=\meta{name}}
		The name of the imported |goban|. Without |name|, the one it had.
	\end{key}

	\section{Games in SGF}
\label{subsec:sgf}

SGF is the format almost every Go program uses to save games. With the Lua kernel (p.~\pageref{subsec:kernel}), \pgfgoname\ can read them.

In every command that reads SGF files, the name can have spaces, accents, parentheses and other scripts. The special characters of \TeX\ are written with a backslash: |\sgf{a\&b.sgf}| reads |a&b.sgf|, and the same holds for |\#|, |\%|, |\$|, |\{| and |\}|. Two spaces in a row are read as one, because \TeX\ joins them before \pgfgoname\ sees the name.

Each file is read only once in the whole document, however many commands use it (if it changes during the compilation, it is read again).

\subsection{Names for the files}
\label{subsec:sgfnombres}

\begin{command}{\pgfgodefinesgf\opt{\oarg{options}}\marg{name}\marg{file}}
	Gives a name to a game: from then on, \meta{name} works in every place where an SGF file goes (|\sgf|, |\sgfinfo|, |\sgfcomment|, |\sgfdiagrams|, |\sgfproblem|, |\sgfanimation|, |sgf size| and |sgf view|). So the path of the file is written only once, and so is what says which game it is (|game| in a collection, the charset, the orientation). A name that is not defined is a file, as always; if a name matches that of a file, the name wins. The options of each command win over those of the name: |\sgfinfo[game = 1]{|\meta{name}|}{PB}| reads game 1 even if the name says another one. With |kernel = expl3|, the name is defined but cannot be used, like the SGF commands.

	\begin{codeexample}[]
		\pgfgodefinesgf{game}{game.sgf}

		\begin{goban}[sgf size = game, scale = 0.6]
			\sgf[to = 7]{game}
		\end{goban}

		Move 7: \sgfcomment{game}{7}
	\end{codeexample}
\end{command}

\begin{key}{/pgfgo/define sgf/game = \meta{number} (initially that of each command)}
	In a collection, the game used by the commands with this name, like |/pgfgo/sgf/game|.

	\begin{codeexample}[code only]
		\pgfgodefinesgf[game = 3]{third}{collection.sgf}
	\end{codeexample}
\end{key}

\begin{key}{/pgfgo/define sgf/charset = \meta{charset} (initially that of \texttt{sgf charset})}
	The charset with which the file is read if it is not UTF-8, like |/pgfgo/sgf charset|, only for this name.

	\begin{codeexample}[code only]
		\pgfgodefinesgf[charset = GBK]{qipu}{games/qipu-without-ca.sgf}
	\end{codeexample}
\end{key}

\begin{key}{/pgfgo/define sgf/origin = \meta{corner} (initially top left)}
	The corner where the point |aa| of the file goes: |top left|, |top right|, |bottom left| or |bottom right|. With |top left| it is read as SGF says. The orientation belongs only to the name: no option of the commands changes it, and the same file used by its path is read as SGF says. It holds for everything read from the file (stones, moves, marks, arrows, views) and for |partial = auto|. The coordinates written in the commands, like those of |variation = {R16, P17}|, are those of the board, not those of the file.
\end{key}

\begin{key}{/pgfgo/define sgf/transpose = \meta{boolean} (default true, initially false)}
	The first letter of each point is the row and the second the column; the size is turned too (a 13$\times$9 game becomes 9$\times$13). Together with |origin|, it gives the eight symmetries of the board.

	\begin{codeexample}[]
		\pgfgodefinesgf[origin = bottom right]{mirror}{game.sgf}
		\pgfgodefinesgf[transpose]{transposed}{game.sgf}

		\begin{goban}[sgf size = mirror, scale = 0.4]
			\sgf[to = 7]{mirror}
		\end{goban}
		\begin{goban}[sgf size = transposed, scale = 0.4]
			\sgf[to = 7]{transposed}
		\end{goban}
	\end{codeexample}
\end{key}

\subsection{Sensei's Library diagrams}
\label{subsec:sgfsensei}

Sensei's Library and forums such as L19 write diagrams as text, with lines that start with |$$|. A file with no SGF nodes and with those lines is read as a collection: each block of consecutive |$$| lines is a game (|game = 2| is the second diagram), and what is between the blocks is not read. The |[go]| and |[/go]| tags of L19 can be left in. The file works everywhere an SGF file goes, and also with |\pgfgodefinesgf|.

\begin{environment}{{pgfgodiagram}\marg{name}}
	Diagrams written in the document, which every SGF command reads as the file \meta{name}. The lines are read as they are: |\begin{pgfgodiagram}|\marg{name} ends its line and |\end{pgfgodiagram}| goes on a line of its own. It does not work inside the argument of another command or inside |pgfgolayout|. If \meta{name} is also the name of a file, the environment wins. With |kernel = expl3|, the content is skipped and there is an error.

	\begin{codeexample}[code only]
\begin{pgfgodiagram}{nadare}
$$Bc Nadare
$$  ---------------
$$ | . . . . . . .
$$ | . . . . . . .
$$ | . . . 1 6 . .
$$ | . . 5 4 . . .
$$ | . . 3 2 a . .
$$ | . . . . . . .
\end{pgfgodiagram}

\sgfdiagrams[goban = {scale = 0.5}]{nadare}
	\end{codeexample}

\begin{pgfgodiagram}{nadare}
$$Bc Nadare
$$  ---------------
$$ | . . . . . . .
$$ | . . . . . . .
$$ | . . . 1 6 . .
$$ | . . 5 4 . . .
$$ | . . 3 2 a . .
$$ | . . . . . . .
\end{pgfgodiagram}

\sgfdiagrams[goban = {scale = 0.5}]{nadare}
\end{environment}

The first line is the header: |$$|, |B| or |W| (who plays move 1; black if absent), a board size (|$$B13|; 19 if absent), |m| and the number of move 1 (|$$Wm67|) and the title. The |c| for coordinates is accepted but not used: the coordinates are those of |goban/label|. Each diagram is read like this:
\begin{itemize}
	\item The title is the name of the game (|GN|) and of its figure (|FG|): |\sgfdiagrams| puts it in the caption.
	\item |X| and |O| are setup stones (|AB| and |AW|). |B|, |#|, |Y| and |Z| are black stones with a circle, square, triangle and cross; |W|, |@|, |Q| and |P|, white ones. |C|, |S|, |T| and |M| are those marks on empty points, and lowercase letters are letters.
	\item The numbers (|1| to |9|, and |0| is 10) are moves, in order: the odd ones belong to whoever plays 1. If a number is missing (``5 at 1''), the following moves keep theirs (|MN|).
	\item The marks, the letters and the lines |{AR C14 C11}| (arrow) and |{LN 5:6 5:9}| (line) go on the last move. Their points are board points (|C14|) or diagram points (|5:6|, column and row from the top left).
	\item The edges (|---|, \verb+|+, |+|) say where the diagram is on the board. Without a size, two opposite edges give that side of the board. Without edges, the diagram goes in the middle. What is drawn is the view of the game (|VW|): |\sgfdiagrams| only shows that part.
	\item |.| and |,| are empty points (the star points come from |goban/hoshi|). |_| is not part of the board: it has no lines (|goban/blank|, p.~\pageref{subsec:parcial}). Links (\verb+[a|page]+) are not read.
\end{itemize}

A line that is not understood is skipped with a warning. A diagram bigger than the size of the header makes the board bigger, also with a warning.

\subsection{A game on a board}
\label{subsec:sgftablero}

\begin{command}{\sgf\opt{\oarg{options}}\marg{file}}

	Draws on the |goban| the main line of the first game of the file: the initial stones (|AB| and |AW|) as with |\stones|, and the moves as with |\moves|. Variations are not drawn. Passes count for the numbering, but are not drawn (with |goban/notes| they are noted, ``N: pass''). It only exists inside |goban|.
\end{command}

The stones that the nodes after the first move add (|AB|, |AW|) or remove (|AE|) are applied together with the previous move. Removing stones needs |captures|; without |captures|, |\sgf| warns once and does not remove them.

At the end, the numbering and the turn are ready to continue the game with |\moves|. |\sgf| works with |partial|, |captures|, |name| and |resume|. If the game is for a board of another size, a warning is given. With |kernel = expl3|, |\sgf| gives an error.

In the examples, the files are used by their name (p.~\pageref{subsec:sgfnombres}). |game| is |game.sgf|, a short 9$\times$9 game:
\begin{verbatim}
(;GM[1]FF[4]SZ[9]PB[Black]PW[White]C[A short game.]
;B[ee];W[ef];B[df];W[ce];B[ff];W[gd];B[eg]C[Black captures.];W[dd];B[cf];W[fd])
\end{verbatim}

\begin{codeexample}[width=5cm]
	\begin{goban}[label = {at = {north, west}}]
		\sgf{game} % Without captures, 2 is not removed.
	\end{goban}
\end{codeexample}

\begin{codeexample}[width=5cm]
	\begin{goban}[label = {at = {north, west}}, captures]
		\sgf{game}
	\end{goban}
\end{codeexample}

\begin{key}{/pgfgo/sgf/from = \meta{number} (initially 1)}

	Number of the first numbered move. The previous moves are drawn without a number.

	\begin{codeexample}[width=5cm]
		\begin{goban}[captures]
			\sgf[from = 7]{game}
		\end{goban}
	\end{codeexample}
\end{key}

\begin{key}{/pgfgo/sgf/to = \meta{number} or |last| (initially last)}

	Number of the last move played. The following ones are not drawn. With |last| (or a negative number), up to the last one.

	\begin{codeexample}[width=5cm]
		\begin{goban}[captures]
			\sgf[to = 6]{game}
			\moves{F5} % Goes on with 7 (Black), elsewhere.
		\end{goban}
	\end{codeexample}
\end{key}


|from| and |to| count the moves of the file. The number shown can be another: the property |MN| of the file changes it (the following moves go on from there), and so does |restart|. In the examples that follow, |variations| is |variations.sgf|, a 7$\times$7 game with two branchings, marks and |MN[10]| on the fourth move of the main line.

\begin{key}{/pgfgo/sgf/variation = \meta{numbers or moves} (initially the main line)}
	Chooses a variation: at each branching of the game, the given variation is taken, in order (1 is the main one). For example, |variation = 2.1| takes the second variation at the first branching and the first one at the second.

	In a file with many branchings, such as a joseki dictionary, it is easier to give the moves of the line from the first one, separated by commas or spaces: |variation = {R16, P17, Q15}|. At each branching, the variation that plays the next move is taken; a pass is written |pass|. Each intersection is a column (one or two capital letters) and a number, as in |\stone|.

	In both cases, when the list ends, the first variation is followed. If a variation (or a move) does not exist, a warning is given and the first one is followed from there.

	\begin{codeexample}[]
		\begin{goban}[sgf size = variations]
			\sgf[variation = 2]{variations}
		\end{goban}

		\begin{goban}[sgf size = variations]
			\sgf[variation = 1.2]{variations}
		\end{goban}

		\begin{goban}[sgf size = variations]
			\sgf[variation = {C5, C3}]{variations}
		\end{goban}
	\end{codeexample}
\end{key}

\begin{key}{/pgfgo/sgf/marks = \meta{boolean} (default true, initially true)}
	Draws the marks of the file (|TR|, |SQ|, |CR|, |MA|, |SL| and the labels of |LB|), the arrows (|AR|) and the lines (|LN|) of the last node shown, and dims the points of |DD| (which stay dimmed until the next |DD|; p.~\pageref{subsec:flechas}). Each mark takes the color of the stone on that point, so that it is seen, and replaces the number of the move. Labels are the exception: they do not cover a number, so on a numbered stone they are not drawn. The territory of the file (|TB| and |TW|) is drawn like that of |\territory| (p.~\pageref{subsec:territorio}) and counts for |\pgfgoscore|.

	\begin{codeexample}[]
		\begin{goban}[sgf size = variations]
			\sgf[to = 2]{variations}
		\end{goban}

		\begin{goban}[sgf size = variations]
			\sgf[to = 2, marks = false]{variations}
		\end{goban}
	\end{codeexample}
\end{key}

\begin{key}{/pgfgo/sgf/step = \meta{number} or \meta{list of numbers} (initially 100)}
	The size of the stretches in which |restart| restarts the numbering. With a list, each number is the size of a stretch, and the last one is repeated: with |step = {50, 100}|, the first stretch has 50 moves and the others 100. It is the same \emph{key} as |/pgfgo/sgf diagrams/step|, where each stretch is a diagram.
\end{key}

\begin{key}{/pgfgo/sgf/restart = \meta{restart options} (default true)}
	From the second stretch of |step| moves on, each one is numbered from 1 (or from |restart/from|), as in books: with |step = 100, restart|, move 101 is shown as 1. The first stretch keeps the numbers of the game. |\sgfanimation| and |\sgfdiagrams| have the same options.

	\begin{codeexample}[]
		\begin{goban}[sgf size = variations]
			\sgf[]{variations}
		\end{goban}

		\begin{goban}[sgf size = variations]
			\sgf[step = 3, restart]{variations}
		\end{goban}
	\end{codeexample}
\end{key}

\begin{key}{/pgfgo/sgf/restart/true}
	Restarts the numbering in each stretch.
\end{key}

\begin{key}{/pgfgo/sgf/restart/false}
	Does not restart it: the moves carry the numbers of the game (with |MN|).
\end{key}

\begin{key}{/pgfgo/sgf/restart/auto}
	Does what the file asks for with |PM|: with |PM[2]|, the numbering restarts every 100 moves, regardless of |step| (101 is shown as 1); otherwise, it does not restart. It is the initial value.
\end{key}

\begin{key}{/pgfgo/sgf/restart/from = \meta{number} (initially 1)}
	The number of the first move of each stretch, from the second on. It also turns |restart| on. With |step = 99| and |restart = {from = 0}|, move 100 is shown as 0, 101 as 1, and so on.

	\begin{codeexample}[]
		\begin{goban}[sgf size = variations]
			\sgf[step = 3, restart = {from = 0}]{variations}
		\end{goban}
	\end{codeexample}
\end{key}

\begin{key}{/pgfgo/sgf/numbers = \meta{true, false or auto} (default true, initially auto)}
	With |false|, the moves are drawn without a number. With |auto|, it does what the file asks for with |PM|: with |PM[0]|, they are not numbered. |\sgfdiagrams| has its own \emph{key} |numbers|, and |\sgfproblem| always numbers.

	\begin{codeexample}[]
		\begin{goban}[sgf size = variations]
			\sgf[numbers = false]{variations}
		\end{goban}
	\end{codeexample}
\end{key}

\begin{key}{/pgfgo/sgf/last mark = \meta{mark options} or |false| (initially false)}
	Marks the last move shown, with the options of a mark (|circle|, |triangle|, |custom = ...|, etc.). A pass is not marked, nor is a point the file already marks or a move with a number (the number already says which is the last): in practice, it marks something with |numbers = false|.

	\begin{codeexample}[]
		\begin{goban}[sgf size = variations]
			\sgf[numbers = false, last mark = circle]{variations}
		\end{goban}
	\end{codeexample}
\end{key}

\begin{key}{/pgfgo/sgf/game = \meta{number} (initially 1)}
	An SGF file can have several games (a collection). Chooses which one is read; if it does not exist, the first one is read with a warning. It also holds for |\sgfanimation|, |\sgfdiagrams|, |\sgfcomment|, |\sgfinfo| and |sgf size|, which have their own \emph{key} |game|; for the whole document, |\pgfgoset{sgf = {game = 2}}|.
\end{key}

\begin{key}{/pgfgo/sgf/variations = \meta{none, children, siblings or auto} (default auto, initially none)}
	Marks with letters the variations of the last move shown: with |children|, the moves that can follow; with |siblings|, the alternatives to that move (the letter replaces its number). |auto| does what the file asks for with |ST|. They are only marked if there is more than one.

	\begin{codeexample}[]
		\begin{goban}[sgf size = variations]
			\sgf[to = 1, marks = false, variations = children]{variations}
		\end{goban}
	\end{codeexample}
\end{key}

|\sgf| also reads |PL| (who plays next: |\moves| goes on with that color) and |CA| (the charset of the file; see |sgf charset|).

\begin{key}{/pgfgo/sgf charset = \meta{charset} (initially auto)}
	Charset of the SGF files that are not in UTF-8, for every command that reads files. With |auto|, it is the one the file gives with |CA| or, if it gives none, Latin-1 (the SGF default). |Latin-1|, |GBK| (and |GB2312|), |Big5|, |Shift_JIS|, |EUC-JP| and |EUC-KR| are read; the names ignore case and punctuation (|shift-jis| and |EUC_KR| also work). A file that is valid UTF-8 is always read as UTF-8, even if |CA| says otherwise; one in UTF-16 (with the byte order mark at the beginning) is read too.

	Many files from Asian servers do not give their charset; to read them, for example, |\pgfgoset{sgf charset = GBK}|. If a file without |CA| seems to be in one of those charsets, a warning reminds of it. With another charset a warning is given and the file is read as Latin-1: it has to be converted to UTF-8 first, for example with |iconv -f KOI8-R -t UTF-8 game.sgf > game-utf8.sgf|.

	To write the comments in Chinese, Japanese or Korean, the package of that language is also needed: |ctex|, |luatexja| or |luatexko|. Examples 14, 18 and 19 of |doc/examples| show a commented game in each one. Without it, LaTeX only notes the missing characters in the |.log|; so, if the font with which a comment or an |\sgfinfo| is written lacks some Chinese, Japanese or Korean character, a warning is given, once per font. With those packages (or |xeCJK|) this is not checked: they write those characters with another font.
\end{key}

\begin{key}{/pgfgo/sgf = \meta{sgf options}}
	Options of every |\sgf| of the document, for example |\pgfgoset{sgf = {variations}}|.
\end{key}

\subsection{Information and comments}
\label{subsec:sgfdatos}

\begin{command}{\sgfinfo\opt{\oarg{options}}\marg{file}\marg{property}}
	Writes a property of the root of the game, for example |PB| and |PW| (the players), |RE| (the result), |KM| (the komi) or |DT| (the date), or with |move| that of a move. If it is not there, it writes nothing. It can be used outside |goban|.

	\begin{codeexample}[]
		\sgfinfo{variations}{PB} (Black) against
		\sgfinfo{variations}{PW} (White),
		komi \sgfinfo{variations}{KM}.
		Result: \sgfinfo{variations}{RE}.
	\end{codeexample}

	The value is written as it is in the file (with the escaped characters already resolved and in the charset of |CA|): |RE| stays as |B+R|, not in words. The properties with information about the game, according to the SGF format, are:

	\begin{center}
		\small
		\begin{tabular}{lp{6.5cm}l}
			\emph{Property} & \emph{What it is} & \emph{Example} \\[2pt]
			|PB|, |PW| & The name of Black and that of White. & |Lee Sedol| \\
			|BR|, |WR| & The rank of Black and that of White. & |9p| \\
			|BT|, |WT| & The team of Black and that of White. & \\
			|RE| & The result: |B+| or |W+| and the points, |R| (resignation), |T| (time) or |F| (forfeit); |0| or |Draw|, a draw; |Void|, no result; |?|, unknown. & |W+3.5| \\
			|KM| & The komi. & |6.5| \\
			|HA| & The handicap: how many stones. & |2| \\
			|DT| & The date, |YYYY-MM-DD| (or several, separated by commas). & |2016-03-09| \\
			|EV| & The tournament or event. & |Meijin| \\
			|RO| & The round. & |3 (final)| \\
			|PC| & The place. & |Tokyo| \\
			|GN| & The name of the game. & \\
			|GC| & A comment about the game: the background, the players. & \\
			|ON| & The opening (fuseki). & |chinese| \\
			|RU| & The rules. & |Japanese| \\
			|TM| & The time of each player, in seconds. & |7200| \\
			|OT| & The overtime (byo-yomi). & |5x30 byo-yomi| \\
			|SO| & The source: a book, a magazine. & \\
			|AN| & Who commented on the game. & \\
			|US| & Who wrote the file. & \\
			|CP| & The copyright. & \\
			|AP| & The program that wrote the file and its version. & |CGoban:3| \\
			|SZ| & The size of the board (|19|, or |19:13| if it is not square). & |19| \\
			|CA| & The charset of the file. & |UTF-8| \\
			|FF|, |GM| & The version of the format and the game (|1|: Go). & |4| \\
		\end{tabular}
	\end{center}

	Any other property of the root (or of the move, with |move|) works. If it has several values, they are written separated by commas: |\sgfinfo{game}{AB}| writes the black stones put before the start in SGF coordinates, |dd, pp|. The long names of the old versions of the format (|PlayerBlack|) are read as their abbreviation: they are asked for with |PB|.
\end{command}

\begin{command}{\sgfcomment\opt{\oarg{options}}\marg{file}\marg{move}}
	Writes the comment of a move of the file (|0|: the one before the first move): the name of the node (|N|), in bold; the annotations of the move (|TE|, |BM|, |DO|, |IT|) and of the position (|GB|, |GW|, |DM|, |UC|); and the comment (|C|). The annotations are written with the texts of |/pgfgo/sgf annotations|. The values (|V|) and the hotspots (|HO|), which are for the game tree, are not written. It can be used outside |goban|. |\sgfanimation| and |\sgfdiagrams| write the comments the same way.

	\begin{codeexample}[]
		Move 7: \sgfcomment{game}{7}
	\end{codeexample}
\end{command}

\begin{key}{/pgfgo/sgf info/game = \meta{number} (initially 1)}
	Like |/pgfgo/sgf/game|, for |\sgfinfo|.
\end{key}

\begin{key}{/pgfgo/sgf info/move = \meta{number} (initially 0)}
	The property of that move: that of its node and that of the nodes without a move that follow it. With |0|, that of the root (and of the nodes before the first move). For example, the time White has left at move 120 of a KGS or OGS game (|WL|, in seconds): |\sgfinfo[move = 120]{game}{WL}|.
\end{key}

\begin{key}{/pgfgo/sgf info/variation = \meta{numbers} (initially the main line)}
	With |move|, the line of the move, like |/pgfgo/sgf/variation|.
\end{key}

\begin{key}{/pgfgo/sgf comment/game = \meta{number} (initially 1)}
	Like |/pgfgo/sgf/game|.
\end{key}

\begin{key}{/pgfgo/sgf comment/variation = \meta{numbers} (initially the main line)}
	Like |/pgfgo/sgf/variation|.

	\begin{codeexample}[]
		\sgfcomment[variation = 2]{problem}{1}
	\end{codeexample}
\end{key}

\begin{key}{/pgfgo/sgf comment/annotations = \meta{boolean} (default true, initially true)}
	Writes the name and the annotations. With |false|, only the comment.

	\begin{codeexample}[]
		With: \sgfcomment{game}{7}
		\quad Without: \sgfcomment[annotations = false]{game}{7}
	\end{codeexample}
\end{key}

\begin{key}{/pgfgo/sgf comment/chat = \meta{boolean} (default true, initially true)}
	With |false|, removes the chat of the servers from the comment: on OGS, everything after the line ``|-- chat --|'' (what comes before is the comment of the reviewer); on KGS, the lines ``|name [rank]: ...|'' and, if there is any, also the lines ``|name: ...|'' of the same comment (those of the players without a rank). If nothing is left, the move has no comment.
\end{key}

The texts of the annotations are in the language of the document (p.~\pageref{subsec:idiomas}), and they can be changed for the whole document with |\pgfgoset{sgf annotations = {...}}|, for example |bad move = Mistake|. The emphasized annotations of |DM| and |UC| use the same text.

\begin{key}{/pgfgo/sgf annotations/tesuji = \meta{text} (initially Tesuji or its translation)}
	Text of |TE|.

	\begin{codeexample}[]
		\pgfgoset{sgf annotations = {tesuji = Good move}}
		\sgfcomment{game}{7}
	\end{codeexample}
\end{key}

\begin{key}{/pgfgo/sgf annotations/very good tesuji = \meta{text} (initially Very good tesuji or its translation)}
	Text of |TE[2]|.
\end{key}

\begin{key}{/pgfgo/sgf annotations/bad move = \meta{text} (initially Bad move or its translation)}
	Text of |BM|.
\end{key}

\begin{key}{/pgfgo/sgf annotations/very bad move = \meta{text} (initially Very bad move or its translation)}
	Text of |BM[2]|.
\end{key}

\begin{key}{/pgfgo/sgf annotations/doubtful = \meta{text} (initially Doubtful move or its translation)}
	Text of |DO|.
\end{key}

\begin{key}{/pgfgo/sgf annotations/interesting = \meta{text} (initially Interesting move or its translation)}
	Text of |IT|.
\end{key}

\begin{key}{/pgfgo/sgf annotations/good for black = \meta{text} (initially Good for Black or its translation)}
	Text of |GB|.
\end{key}

\begin{key}{/pgfgo/sgf annotations/very good for black = \meta{text} (initially Very good for Black or its translation)}
	Text of |GB[2]|.
\end{key}

\begin{key}{/pgfgo/sgf annotations/good for white = \meta{text} (initially Good for White or its translation)}
	Text of |GW|.
\end{key}

\begin{key}{/pgfgo/sgf annotations/very good for white = \meta{text} (initially Very good for White or its translation)}
	Text of |GW[2]|.
\end{key}

\begin{key}{/pgfgo/sgf annotations/even = \meta{text} (initially Even position or its translation)}
	Text of |DM|.
\end{key}

\begin{key}{/pgfgo/sgf annotations/unclear = \meta{text} (initially Unclear position or its translation)}
	Text of |UC|.
\end{key}

\subsection{Animations}
\label{subsec:sgfanimacion}

\begin{command}{\sgfanimation\opt{\oarg{options}}\marg{file}}

	Steps through the game with the |animate| package: it draws one |goban| per move, and the buttons of |animate| step forward, back or play it. Frame $k$ shows the position after move $k$; frame 0, the position before the first move. It is used outside |goban|. The boards have the size of the game (|SZ|), unless |goban| gives another one. It needs |\usepackage{animate}| and two compilations.
\end{command}

The animations of |animate| work in Acrobat Reader and in some other viewers (Foxit, PDF-XChange). In the rest only the first frame is seen. These options can be set for the whole document with |\pgfgoset{sgf animation = {...}}|.

\begin{codeexample}[width=5cm]
	\sgfanimation[comments, goban = {captures}]{game}
\end{codeexample}

\begin{key}{/pgfgo/sgf animation/from = \meta{number} (initially 0)}
	First frame.
\end{key}

\begin{key}{/pgfgo/sgf animation/to = \meta{number} or |last| (initially last)}
	Last frame. With |last| (or |-1|), that of the last move.
\end{key}

\begin{key}{/pgfgo/sgf animation/fps = \meta{number} (initially 1)}
	Frames per second when playing.
\end{key}

\begin{key}{/pgfgo/sgf animation/goban = \meta{goban options}}
	Options of each board, as in |\begin{goban}[...]|. As in |\sgfdiagrams|, the view of the file is applied after them: to fix the visible part there is |partial|.
\end{key}

\begin{key}{/pgfgo/sgf animation/animate = \meta{animate options} (initially controls)}
	Options of the |animateinline| environment of |animate|, for example |controls, loop|.
\end{key}

\begin{key}{/pgfgo/sgf animation/numbers = \meta{boolean} (default true, initially false)}
	Numbers every move, like |\sgf|. Without |numbers|, the moves have no number and the last one is marked.

	\begin{codeexample}[width=5cm]
		\sgfanimation[from = 6, to = 8, numbers, goban = {captures}]{game}
	\end{codeexample}
\end{key}

\begin{key}{/pgfgo/sgf animation/last mark = \meta{mark options or false} (initially circle)}
	Mark of the last move, with the color of whoever played it. With |false|, the last move is not marked. It is not used with |numbers|, nor when the file already marks that point.

	\begin{codeexample}[width=5cm]
		\sgfanimation[from = 5, to = 7, last mark = triangle, fps = 2,
			animate = {controls, loop}, goban = {captures}]{game}
	\end{codeexample}
\end{key}

\begin{key}{/pgfgo/sgf animation/comments = \meta{comments options} (default true)}
	Shows the comment (|C|) of the move below the board. That of frame 0 is that of the nodes before the first move. It is turned on with |comments| or |comments = true| and off with |false| (the initial value); |height| and |font| do not turn it on: |comments = {true, height = 6em}|.
\end{key}

\begin{stylekey}{/pgfgo/sgf animation/comments/true}
	Shows the comments.
\end{stylekey}

\begin{stylekey}{/pgfgo/sgf animation/comments/false}
	Does not show the comments. It is the initial value.
\end{stylekey}

\begin{key}{/pgfgo/sgf animation/comments/height = \meta{dimension} (initially 4em)}
	Height of the space for the comments. It is the same in every frame, so that they do not change size.
\end{key}

\begin{key}{/pgfgo/sgf animation/comments/font = \meta{font commands} (initially \texttt{\string\small})}
	Font of the comments.

	\begin{codeexample}[width=5cm]
		\sgfanimation[from = 7, to = 9, goban = {captures},
			comments = {true, height = 2.5em, font = \itshape}]{game}
	\end{codeexample}
\end{key}

\begin{key}{/pgfgo/sgf animation/variation = \meta{numbers} (initially the main line)}
	Like |/pgfgo/sgf/variation|: the animation goes through the line that leads to that variation.

	\begin{codeexample}[width=5cm]
		\sgfanimation[variation = 2, numbers]{variations}
	\end{codeexample}
\end{key}

\begin{key}{/pgfgo/sgf animation/marks = \meta{boolean} (default true, initially true)}
	Like |/pgfgo/sgf/marks|: each frame shows the marks of its move.
\end{key}

\begin{key}{/pgfgo/sgf animation/step = \meta{number} or \meta{list of numbers} (initially 100)}
	Like |/pgfgo/sgf/step|.
\end{key}

\begin{key}{/pgfgo/sgf animation/restart = \meta{restart options} (default true)}
	Like |/pgfgo/sgf/restart|.
\end{key}

\begin{key}{/pgfgo/sgf animation/restart/true}
	Like |/pgfgo/sgf/restart/true|.
\end{key}

\begin{key}{/pgfgo/sgf animation/restart/false}
	Like |/pgfgo/sgf/restart/false|.
\end{key}

\begin{key}{/pgfgo/sgf animation/restart/auto}
	Like |/pgfgo/sgf/restart/auto|. It is the initial value.
\end{key}

\begin{key}{/pgfgo/sgf animation/restart/from = \meta{number} (initially 1)}
	Like |/pgfgo/sgf/restart/from|.
\end{key}

\begin{key}{/pgfgo/sgf animation/annotations = \meta{boolean} (default true, initially true)}
	With |comments|, also writes the name and the annotations of each move, like |\sgfcomment|.
\end{key}

\begin{key}{/pgfgo/sgf animation/chat = \meta{boolean} (default true, initially true)}
	With |comments|, like |/pgfgo/sgf comment/chat|: with |false|, without the chat of KGS and OGS.
\end{key}

\begin{key}{/pgfgo/sgf animation/game = \meta{number} (initially 1)}
	Like |/pgfgo/sgf/game|.
\end{key}

\begin{key}{/pgfgo/sgf animation/view = \meta{boolean} (default true, initially true)}
	If the file has a view (|VW|), the boards show only that part, like |sgf view|. Every frame has the same board: if the view changes during the game, the board covers the views of every frame (and if in any of them the whole board is seen, the whole board is seen).
\end{key}

\begin{key}{/pgfgo/sgf animation/partial = \meta{auto, false or corners} (initially false)}
	With |auto|, where the file has no view, the boards show only the part used by the line that is animated (that of |variation|), as in |\sgfdiagrams|. With two corners (|partial = M12:T19|), every board shows that part, even if the file has a view. Every frame has the same board.
\end{key}

\begin{key}{/pgfgo/sgf animation/layers = \meta{boolean} (default true, initially true)}
	Builds the animation in layers: the board is drawn only once, and so is each stone, on a transparency that |animate| shows while the stone is still on the board (with the |timeline| option of |animate|). The only thing drawn in each frame is what changes: the mark of the last move, the marks of the file and the comment. So the PDF grows with the number of moves and not with its square: a game of 200 moves takes less than 1~MB, instead of more than 10~MB, and it also compiles faster. With |layers = false|, each frame is a whole |goban|. The |timeline| is kept in an auxiliary file, |<jobname>-pgfgo-<n>.tln|.
\end{key}

\subsection{Diagrams as in books}
\label{subsec:sgfdiagramas}

\begin{command}{\sgfdiagrams\opt{\oarg{options}}\marg{file}}
	Draws the game as in books: one diagram every |step| moves (1--50, 51--100, etc.). Each diagram starts from the real position, with the previous moves without a number and their captures; the stones captured inside the diagram are not removed (they are removed in the next one). Below each one go its caption and its ``N at M'' notes (p.~\pageref{subsec:notas}). The caption gives the numbers of the first and the last move, the ones on the stones if the file changes the numbering (|MN|). A game without moves is a single diagram, of its position, with no numbers in the caption. It is used outside |goban|. The options can be set for the whole document with |\pgfgoset{sgf diagrams = {...}}|.

	\begin{codeexample}[code only]
		\sgfdiagrams[step = 5, title = Diagram]{game}
	\end{codeexample}
\end{command}

\sgfdiagrams[step = 5, title = Diagram]{game}

\begin{key}{/pgfgo/sgf diagrams/step = \meta{number} or \meta{list of numbers} (initially 50)}
	Number of moves of each diagram. With a list, the first diagram has as many moves as the first number says, the second as many as the second says, and so on; the last number is repeated: with |step = {5, 10, 3}|, the diagrams have 5, 10, 3, 3, 3... moves.
\end{key}

\begin{key}{/pgfgo/sgf diagrams/from = \meta{number} (initially 1)}
	First move. The diagrams are the same as without |from|, with the same number: only those with moves from |from| on are drawn, and the first one starts at |from|. With |step = 50| and |from = 120|, the first one is ``Diagram 3 (120--150)''.
\end{key}

\begin{key}{/pgfgo/sgf diagrams/to = \meta{number} or |last| (initially last)}
	Last move. With |last| (or |-1|), the last one of the game.
\end{key}

\begin{key}{/pgfgo/sgf diagrams/numbers = \meta{boolean} (default true, initially true)}
	With |false|, the moves of each diagram are drawn without a number (the ``N at M'' notes still use the numbers of the game).
\end{key}

\begin{key}{/pgfgo/sgf diagrams/last mark = \meta{mark options} or |false| (initially false)}
	Like |/pgfgo/sgf/last mark|, in each diagram.

	\begin{codeexample}[code only]
		\sgfdiagrams[step = 5, to = 10, numbers = false, last mark = circle]{game}
	\end{codeexample}

	\sgfdiagrams[step = 5, to = 10, numbers = false, last mark = circle]{game}
\end{key}

\begin{key}{/pgfgo/sgf diagrams/restart = \meta{restart options} (default true)}
	Like |/pgfgo/sgf/restart|, with one diagram per stretch: from the second diagram on, each one is numbered from 1 (or from |restart/from|), as some books do. The first diagram keeps the numbers of the game.

	\begin{codeexample}[code only]
		\sgfdiagrams[step = {2, 3}, restart = {from = 0}]{game}
	\end{codeexample}
\end{key}

\sgfdiagrams[step = {2, 3}, to = 8, restart = {from = 0}, title = Diagram]{game}

\begin{key}{/pgfgo/sgf diagrams/restart/true}
	Restarts the numbering in each diagram.
\end{key}

\begin{key}{/pgfgo/sgf diagrams/restart/false}
	Does not restart it: the moves carry the numbers of the game.
\end{key}

\begin{key}{/pgfgo/sgf diagrams/restart/auto}
	Like |/pgfgo/sgf/restart/auto|: with |PM[2]|, every 100 moves, regardless of the diagrams. It is the initial value.
\end{key}

\begin{key}{/pgfgo/sgf diagrams/restart/from = \meta{number} (initially 1)}
	The number of the first move of each diagram, from the second on. It also turns |restart| on. With |step = 99| and |restart = {from = 0}|, the first diagram goes from 1 to 99, and in the second one move 100 is shown as 0, 101 as 1, and so on.
\end{key}

\begin{key}{/pgfgo/sgf diagrams/marks = \meta{boolean} (default true, initially true)}
	Like |/pgfgo/sgf/marks|: each diagram shows the marks, arrows and lines of its last move. With |marks = false| they are not drawn.
\end{key}

\begin{key}{/pgfgo/sgf diagrams/notes = \meta{boolean} (default true, initially true)}
	Writes its notes below each diagram: ``N at M'' and the passes, ``N: pass''.
\end{key}

\begin{key}{/pgfgo/sgf diagrams/title = \meta{text} (initially Diagram or its translation)}
	Word of the caption: ``Diagram 2 (51--100)''.
\end{key}

\begin{key}{/pgfgo/sgf diagrams/caption = \meta{caption options}}
	Options of the caption and of the notes.
\end{key}

\begin{key}{/pgfgo/sgf diagrams/caption/font = \meta{font commands} (initially \texttt{\string\small})}
	Font of the caption and of the notes: |caption = {font = \itshape}|.

	\begin{codeexample}[code only]
		\sgfdiagrams[from = 6, title = Figure, caption = {font = \itshape},
			goban = {label = {at = {north, west}}}]{game}
	\end{codeexample}

	\sgfdiagrams[from = 6, title = Figure, caption = {font = \itshape},
		goban = {label = {at = {north, west}}}]{game}
\end{key}

\begin{key}{/pgfgo/sgf diagrams/caption/min width = \meta{dimension} (initially 8em)}
	With |layout = {columns = auto}|, a caption wider than its board is broken into lines at the width of the board. If the board is narrower than this measure (in the font of the caption), it is broken at this measure; it is never broken narrower than its longest word, nor widened beyond its width on one line. In the grid, the caption is broken at the width of the cell.
\end{key}

\begin{key}{/pgfgo/sgf diagrams/goban = \meta{goban options}}
	Options of each board. By default, the boards have the size of the game and use |captures| and |notes|. The view of the file (|view|) is applied after these options: if the file has a view, |goban = {partial = ...}| does not change the visible part; for that there is |partial|.
\end{key}

\begin{key}{/pgfgo/sgf diagrams/name = \meta{name}}
	Each diagram is a named |goban| (|/pgfgo/goban/name|): this name followed by the number of its caption. With |name = var|, the diagrams are called |var1|, |var2|... and the variations |var2a|, |var2b|... They can be used with |resume| and in the views in the text (p.~\pageref{subsec:vistas}). Each one keeps the whole position up to its last move, with the captures resolved: a stone of the numbered stretch captured in the same diagram stays in the drawing, but not in the memory.

	\begin{codeexample}[]
		\sgfdiagrams[step = 2, to = 3, name = var, goban = {scale = 0.5}]{variations}

		Diagram 2 in the text: \pgfgoview{var2}{B2:F6}.
	\end{codeexample}
\end{key}

\begin{key}{/pgfgo/sgf diagrams/layout = \meta{layout options}}
	How the diagrams are laid out: options of |/pgfgo/layout| (p.~\pageref{subsec:disposicion}) only for |\sgfdiagrams|. For example, |layout = {columns = 3}|.
\end{key}

\begin{key}{/pgfgo/sgf diagrams/variation = \meta{numbers} (initially the main line)}
	Like |/pgfgo/sgf/variation|.
\end{key}

\begin{key}{/pgfgo/sgf diagrams/variations = \meta{variations options} (default commented)}
	Also draws the variations of the game, each one in its own diagram, after the diagram where its first move starts: ``Diagram 3a'', ``Diagram 3b''... The moves of each variation are numbered from 1, and the caption says which move the first one replaces (|variation title|). The previous ones are seen without a number, as they were. The options say which ones are drawn, and they add up: |variations = {commented, 5}|. Besides the ones below, anything else is the path of a variation, as in |variation| (|2|, |2.3.1|, or its moves in braces and separated by spaces: |{R16 P17 Q15}|): that one is always drawn, at any depth. The |1|'s at the end do not change the line: |2.3.1| is the same as |2.3|. If the file does not have it, a warning is given. Its comments and notes are written like those of the other diagrams. With |variation|, the variations are those that leave that line (the main one too).

	\begin{codeexample}[]
		\sgfdiagrams[step = 2, to = 3, variations = all, goban = {scale = 0.5}]{variations}
	\end{codeexample}
\end{key}

\begin{stylekey}{/pgfgo/sgf diagrams/variations/none}
	None. It is the initial value, and clears the previous ones.
\end{stylekey}

\begin{stylekey}{/pgfgo/sgf diagrams/variations/commented}
	Those that have some comment (|C|) in their own moves. It is what plain |variations| does.
\end{stylekey}

\begin{stylekey}{/pgfgo/sgf diagrams/variations/all}
	All of them.
\end{stylekey}

\begin{key}{/pgfgo/sgf diagrams/variations/depth = \meta{number} (initially 1)}
	To what depth |commented| and |all| draw: |1|, only those that leave the line that is drawn (the main one, or that of |variation|); |2|, also the variations of those, and so on. Those of a variation go after it, and their caption says which of its moves they replace.
\end{key}

\begin{key}{/pgfgo/sgf diagrams/variations/from = \meta{number} (initially 1)}
	From which move |commented| and |all| draw: only the variations that start at that move or later.

	\begin{codeexample}[code only]
		\sgfdiagrams[variation = {Q16, R14}, variations = {commented, from = 3}]{josekis}
	\end{codeexample}
\end{key}

\begin{key}{/pgfgo/sgf diagrams/variation title = \meta{text} (initially ``variation at'' and the number, or its translation)}
	What goes in parentheses in the caption of a variation: ``Diagram 3a (variation at 26)''. |\pgfgovariationmove| is the number of the move it replaces.
\end{key}

\begin{key}{/pgfgo/sgf diagrams/comments = \meta{comments options} (default true)}
	Writes below each diagram the comments of its moves, like |\sgfcomment|, each one with the number of the move. It is turned on with |comments| or |comments = true| and off with |false| (the initial value); the other options do not turn it on: |comments = {true, align = justify}|. So |\pgfgoset{sgf diagrams = {comments = {font = \itshape}}}| sets the style of the whole document without turning the comments on.
\end{key}

\begin{stylekey}{/pgfgo/sgf diagrams/comments/true}
	Writes the comments.
\end{stylekey}
\begin{stylekey}{/pgfgo/sgf diagrams/comments/false}
	Does not write the comments. It is the initial value.
\end{stylekey}

\begin{key}{/pgfgo/sgf diagrams/comments/font = \meta{font commands} (initially that of the caption)}
	Font of the comments. Without this option, that of |caption/font|.
\end{key}

\begin{key}{/pgfgo/sgf diagrams/comments/align = \meta{left, right, center or justify} (initially left)}
	Alignment of the comments in the width of the diagram. With |justify|, a narrow diagram gets large spaces between the words.
\end{key}

\begin{key}{/pgfgo/sgf diagrams/comments/root = \meta{boolean} (default true, initially false)}
	Also writes the comment of the root (that of the nodes before the first move), without a number, before those of the diagram that starts at move 1.

	\begin{codeexample}[code only]
		\sgfdiagrams[step = 10, comments = {true, root, font = \itshape}]{game}
	\end{codeexample}
\end{key}

\begin{key}{/pgfgo/sgf diagrams/comments/number = \meta{text} (initially \texttt{\string\textbf\{\#1\}:\string~})}
	What goes before each comment; |#1| is the number of the move, the one seen on the board (with |restart|, the restarted one). For example, |number = {Move~#1.~}|.
\end{key}

\begin{key}{/pgfgo/sgf diagrams/comments/separator = \meta{text} (initially one paragraph per comment)}
	What goes between the comments. With the initial value, each one goes in its own paragraph; with |separator = {\quad}|, they run on.

	\begin{codeexample}[code only]
		\sgfdiagrams[step = 10, comments = {true, number = {\textbf{#1}.~}, separator = {~}}]{game}
	\end{codeexample}
\end{key}

\begin{key}{/pgfgo/sgf diagrams/comments/position = \meta{below, right, left or after} (initially below)}
	Where the comments go: below the caption, beside the diagram, as in books of commented games, or with |after| all together after the diagrams, across the page, each group with the caption of its diagram as a heading (they can break across pages). Beside, they start at the height of the top border of the board; the caption, the notes and the result go on below the board.

	\begin{codeexample}[code only]
		\sgfdiagrams[step = 10, comments = {true, position = right}]{game}
	\end{codeexample}
\end{key}

\begin{key}{/pgfgo/sgf diagrams/comments/width = \meta{dimension} or auto (initially auto)}
	Width of the comments beside. With |auto|, what is left: outside the grid, the rest of the line (each diagram goes on its own line); in the grid (|layout = {columns = ...}|), the rest of the cell, and with |fit| the board takes half of the cell and the comments the other half. With a dimension, in the grid the board takes what is left of the cell.
\end{key}

\begin{key}{/pgfgo/sgf diagrams/comments/sep = \meta{dimension} (initially 1em)}
	Space between the diagram and the comments beside it.

	\begin{codeexample}[code only]
		\sgfdiagrams[step = 10, layout = {columns = 2, fit},
			comments = {true, position = left, sep = 2em}]{game}
	\end{codeexample}
\end{key}

\begin{key}{/pgfgo/sgf diagrams/comments/height = \meta{dimension} or none (initially none)}
	Maximum height of the comments of a diagram, below or beside. What does not fit goes after the diagrams, as with |after|, with the caption of the diagram as a heading.

	\begin{codeexample}[code only]
		\sgfdiagrams[step = 30, layout = {columns = 2},
			comments = {true, height = 6\baselineskip}]{game}
	\end{codeexample}
\end{key}

\begin{key}{/pgfgo/sgf diagrams/annotations = \meta{boolean} (default true, initially true)}
	With |comments|, also writes the name and the annotations of each move.

	\begin{codeexample}[code only]
		\sgfdiagrams[step = 10, comments]{game}
		\sgfdiagrams[step = 10, comments, annotations = false]{game}
	\end{codeexample}

	\sgfdiagrams[step = 10, comments]{game}
	\sgfdiagrams[step = 10, comments, annotations = false]{game}
\end{key}

\begin{key}{/pgfgo/sgf diagrams/chat = \meta{boolean} (default true, initially true)}
	With |comments|, like |/pgfgo/sgf comment/chat|: with |false|, without the chat of KGS and OGS. A move that only had chat left does not appear.
\end{key}

\begin{key}{/pgfgo/sgf diagrams/figures = \meta{boolean} (default true, initially true)}
	If the file marks figures (|FG|), each figure is a diagram, with its name as the caption, instead of one every |step| moves.
\end{key}

\begin{key}{/pgfgo/sgf diagrams/view = \meta{boolean} (default true, initially true)}
	If the file has a view (|VW|), each diagram shows only that part, like |sgf view|. If the view changes during the game, each diagram covers the views of its moves.
\end{key}

\begin{key}{/pgfgo/sgf diagrams/partial = \meta{auto, false or corners} (initially false)}
	With |auto|, where the file has no view, each diagram shows only the part of the board used by the lines that are drawn (that of |variation| and those of |variations|), like |/pgfgo/goban/partial = auto|. It is the same in every diagram, also in those of the variations. With two corners, as in |/pgfgo/goban/partial| (|K10:T19|, |{K10, T19}| or |{K10 to T19}|), every diagram shows that part, even if the file has a view.

	\begin{codeexample}[code only]
		\sgfdiagrams[partial = auto, variations]{josekis}
		\sgfdiagrams[partial = K10:T19, variations]{josekis}
		\sgfdiagrams[partial = auto, variation = {Q16, R14},
			variations = {commented, from = 3}, marks = false]{josekis}
	\end{codeexample}
\end{key}

\begin{key}{/pgfgo/sgf diagrams/game = \meta{number} (initially 1)}
	Like |/pgfgo/sgf/game|.
\end{key}

\begin{key}{/pgfgo/sgf diagrams/result = \meta{result options} (default true)}
	Writes below the last diagram, if it reaches the end of the game, the result (|RE|) and, if the file marks the territory (|TB|, |TW|), the score of each player: ``Result: Black wins by 0.5 points (Black 6, White 5.5)''. The score is counted as with |\pgfgoscore| (p.~\pageref{subsec:territorio}), with the komi of the file (|KM|). It is turned off with |result = false|; |title| and |format| do not turn it on.
\end{key}

\begin{stylekey}{/pgfgo/sgf diagrams/result/true}
	Writes the result. It is the initial value.
\end{stylekey}

\begin{stylekey}{/pgfgo/sgf diagrams/result/false}
	Does not write the result.
\end{stylekey}

\begin{key}{/pgfgo/sgf diagrams/result/title = \meta{text} (initially Result or its translation)}
	Word before the result.
\end{key}

\begin{key}{/pgfgo/sgf diagrams/result/format = \meta{text or raw} (initially text)}
	How the result is written. |text|: in words, in the language of the document (p.~\pageref{subsec:idiomas}), ``Black wins by resignation''; |raw|: as it is in the file, ``B+R''. With |text|, |0| or |Draw| (a draw), |Void| (no result), |?| (unknown) and |B+| or |W+| followed by nothing, by the difference, by |R| or |Resign|, by |T| or |Time|, or by |F| or |Forfeit| are understood. Any other result is written as it is.
\end{key}

\begin{key}{/pgfgo/sgf diagrams/black name = \meta{text} (initially Black or its translation)}
	Name of Black in the score and in the result.
\end{key}

\begin{key}{/pgfgo/sgf diagrams/white name = \meta{text} (initially White or its translation)}
	Name of White in the score and in the result.

	\begin{codeexample}[code only]
		\sgfdiagrams[result = {title = Final, format = raw}]{variations}
		\sgfdiagrams[black name = Ann, white name = Bob]{variations}
	\end{codeexample}

	\sgfdiagrams[result = {title = Final, format = raw}]{variations}
	\sgfdiagrams[black name = Ann, white name = Bob]{variations}
\end{key}

\begin{command}{\pgfgowinner}
	In the texts of the result (|\pgfgodefinelanguage|), the name of the winner: |black name| or |white name|.
\end{command}

\begin{command}{\pgfgovariationmove}
	In the caption of a variation (|variation title|) and in its text in each language (|variation at|), the number of the move it replaces.
\end{command}

\begin{command}{\pgfgopoints}
	In the texts of the result, the difference in points, as it is in the file.
\end{command}

\subsection{Problems}
\label{subsec:sgfproblemas}

\begin{command}{\sgfproblem\opt{\oarg{options}}\marg{file}}
	Draws a problem (tsumego). First, the statement: the initial position with its marks, who plays (|PL|, or the color of the first move) and the comment at the beginning. Then, one diagram for each line of the tree of the file (each variation of the solution, right or wrong), numbered from 1, with its ``N at M'' notes and the comments and annotations of its moves, which usually say whether the variation is right. The diagrams are drawn like those of |\sgfdiagrams|. If the file has no view (|VW|), only the part of the board the problem uses is shown. It is used outside |goban|. The options can be set for the whole document with |\pgfgoset{sgf problem = {...}}|. In the examples, |problem| is |problem.sgf|, a life-and-death problem: the right move and two wrong ones, with their comments.

	\begin{codeexample}[code only]
		\sgfproblem[title = Problem 1, variation title = Line,
			black to play = Black plays]{problem}
	\end{codeexample}
\end{command}

\sgfproblem[title = Problem 1, variation title = Line,
	black to play = Black plays]{problem}

\begin{key}{/pgfgo/sgf problem/statement = \meta{true, false or alone} (default true, initially true)}
	Draws the statement. With |alone|, the statement goes alone on its row, centered, and the variations below.
\end{key}

\begin{key}{/pgfgo/sgf problem/solutions = \meta{boolean} (default true, initially true)}
	Draws the variations. With |solutions = false|, only the statement remains; the solutions go elsewhere, with |statement = false|.

	\begin{codeexample}[code only]
		\sgfproblem[solutions = false, title = Problem 2]{problem}
	\end{codeexample}

	\sgfproblem[solutions = false, title = Problem 2]{problem}
\end{key}

\begin{key}{/pgfgo/sgf problem/title = \meta{text} (initially empty)}
	Beginning of the caption of the statement: ``Problem 1: Black to play''. With several games (|game = 1:20|), it is followed by the number of each one; empty, it is |Problem| or its translation: ``Problem 7: Black to play''. The solutions also say which problem they belong to: ``Problem 7: Variation 1''.
\end{key}

\begin{key}{/pgfgo/sgf problem/to play = \meta{auto, black, white or none} (initially auto)}
	Who plays, for the caption of the statement. With |auto|, what |PL| says or, if it is not there, the color of the first move; if there is neither, it is not said. Many collections do not say it in the file: |to play = black|. With |none|, the caption does not say it.
\end{key}

\begin{key}{/pgfgo/sgf problem/verdict = \meta{boolean} (default true, initially false)}
	The caption of each solution says whether it is right or wrong (``Variation 2 (wrong)''), with the texts |correct| and |wrong| of the language. It is decided with the annotations of the line: the last |GB| or |GW| (the position is good for Black or for White), according to who plays; if there is none, the first move with |TE| or |BM| (the odd moves belong to whoever plays: a |BM| of theirs is a mistake, one of the opponent is not); otherwise, the word |RIGHT| in a comment, as in the files of goproblems.com. If there is none of that, it says nothing.

	\begin{codeexample}[]
		\sgfproblem[verdict, statement = false, comments = false,
			goban = {scale = 0.5}]{problem}
	\end{codeexample}
\end{key}

\begin{key}{/pgfgo/sgf problem/variation title = \meta{text} (initially Variation or its translation)}
	Word of the caption of each variation: ``Variation 2''.
\end{key}

\begin{key}{/pgfgo/sgf problem/black to play = \meta{text} (initially Black to play or its translation)}
	Text of the statement when Black plays.
\end{key}

\begin{key}{/pgfgo/sgf problem/white to play = \meta{text} (initially White to play or its translation)}
	Text of the statement when White plays.
\end{key}

\begin{key}{/pgfgo/sgf problem/comments = \meta{comments options} (default true)}
	Writes the comment at the beginning below the statement and those of each variation below its diagram, like |/pgfgo/sgf diagrams/comments|. Here they are on at the start; they are turned off with |comments = false|. The options are |true|, |false|, |font|, |align|, |number|, |separator|, |position|, |width|, |sep| and |height|, as in |\sgfdiagrams|; |font| and |align| also hold for the comment of the statement.
\end{key}

\begin{stylekey}{/pgfgo/sgf problem/comments/true}
	Writes the comments. It is the initial value.
\end{stylekey}

\begin{stylekey}{/pgfgo/sgf problem/comments/false}
	Does not write the comments.
\end{stylekey}

\begin{key}{/pgfgo/sgf problem/comments/font = \meta{font commands} (initially that of the caption)}
	Like |/pgfgo/sgf diagrams/comments/font|.
\end{key}

\begin{key}{/pgfgo/sgf problem/comments/align = \meta{left, right, center or justify} (initially left)}
	Like |/pgfgo/sgf diagrams/comments/align|.
\end{key}

\begin{key}{/pgfgo/sgf problem/comments/number = \meta{text} (initially \texttt{\string\textbf\{\#1\}:\string~})}
	Like |/pgfgo/sgf diagrams/comments/number|.
\end{key}

\begin{key}{/pgfgo/sgf problem/comments/separator = \meta{text} (initially one paragraph per comment)}
	Like |/pgfgo/sgf diagrams/comments/separator|.
\end{key}

\begin{key}{/pgfgo/sgf problem/comments/position = \meta{below, right, left or after} (initially below)}
	Like |/pgfgo/sgf diagrams/comments/position|.
\end{key}

\begin{key}{/pgfgo/sgf problem/comments/width = \meta{dimension} or auto (initially auto)}
	Like |/pgfgo/sgf diagrams/comments/width|.
\end{key}

\begin{key}{/pgfgo/sgf problem/comments/sep = \meta{dimension} (initially 1em)}
	Like |/pgfgo/sgf diagrams/comments/sep|.
\end{key}

\begin{key}{/pgfgo/sgf problem/comments/height = \meta{dimension} or none (initially none)}
	Like |/pgfgo/sgf diagrams/comments/height|.
\end{key}

\begin{key}{/pgfgo/sgf problem/annotations = \meta{boolean} (default true, initially true)}
	With |comments|, also writes the name and the annotations of each move (|TE|, |BM|, etc.).
\end{key}

\begin{key}{/pgfgo/sgf problem/chat = \meta{boolean} (default true, initially true)}
	With |comments|, like |/pgfgo/sgf comment/chat|: with |false|, without the chat of KGS and OGS.
\end{key}

\begin{key}{/pgfgo/sgf problem/notes = \meta{boolean} (default true, initially true)}
	Writes the ``N at M'' notes of each variation.
\end{key}

\begin{key}{/pgfgo/sgf problem/caption = \meta{caption options}}
	Options of the captions, the notes and the comments.
\end{key}

\begin{key}{/pgfgo/sgf problem/caption/font = \meta{font commands} (initially \texttt{\string\small})}
	Font of the captions, the notes and the comments.

	\begin{codeexample}[code only]
		\sgfproblem[statement = alone, comments = false, notes = false,
			caption = {font = \itshape}]{problem}
	\end{codeexample}

	\sgfproblem[statement = alone, comments = false, notes = false,
		caption = {font = \itshape}]{problem}
\end{key}

\begin{key}{/pgfgo/sgf problem/caption/min width = \meta{dimension} (initially 8em)}
	As in |\sgfdiagrams|: with |columns = auto|, the minimum width at which a caption wider than its board is broken (``Problem 3: White to play'' below a small corner).
\end{key}

\begin{key}{/pgfgo/sgf problem/goban = \meta{goban options}}
	Options of each board, like |/pgfgo/sgf diagrams/goban|.
\end{key}

\begin{key}{/pgfgo/sgf problem/name = \meta{name}}
	Like |/pgfgo/sgf diagrams/name|: the statement is called |name|, and each variation |name| followed by its number. With several games (|game|), the game is added: the statement of problem 3 is |name3|, and its variations |name3-1|, |name3-2|...
\end{key}

\begin{key}{/pgfgo/sgf problem/layout = \meta{layout options}}
	How the statement and the variations are laid out: options of |/pgfgo/layout| (p.~\pageref{subsec:disposicion}) only for |\sgfproblem|.
\end{key}

\begin{key}{/pgfgo/sgf problem/view = \meta{boolean} (default true, initially true)}
	If the file has a view (|VW|) at the beginning, every diagram shows that part.
\end{key}

\begin{key}{/pgfgo/sgf problem/partial = \meta{auto, false or corners} (initially auto)}
	With |auto|, if the view of the file is not used, it shows only the part of the board where there are stones, moves, marks, arrows or lines in some line of the problem, like |/pgfgo/goban/partial = auto|. It is the same in every diagram. The lines added around are changed with |goban = {partial margin = ...}|. With |false|, the whole board is seen. With two corners (|partial = A14:G19|), every diagram shows that part, even if the file has a view.

	\begin{codeexample}[code only]
		\sgfproblem[solutions = false, goban = {partial margin = 3}]{problem}
	\end{codeexample}

	\sgfproblem[solutions = false, goban = {partial margin = 3}]{problem}
\end{key}

\begin{key}{/pgfgo/sgf problem/game = \meta{numbers} (initially 1)}
	Like |/pgfgo/sgf/game|: in a collection, the problem that is drawn. Also several, with numbers and ranges like those of the coordinates: |game = 1:20| or |game = {1:4, 9}|. Then all the statements go first, numbered (|title|), and then the solutions of those that have them. The games that the file does not have are left out, with a warning; if none is left, the first one is drawn. That is how a problem sheet is built:

	\begin{codeexample}[code only]
		\sgfproblem[game = 1:24, to play = black, comments = false,
			layout = {columns = 6, fit}]{collection.sgf}
	\end{codeexample}
\end{key}

\begin{key}{/pgfgo/sgf problem/links = \meta{boolean} (default true, initially false)}
	With |hyperref|, the caption of the statement is a link to its first solution, and that of each solution, a link to the statement, also if they are in another |\sgfproblem| (with |solutions = false| and |statement = false|). Two compilations are needed; a link whose target is not there is not set. Without |hyperref|, the package warns.

	\begin{codeexample}[]
		\sgfproblem[links, comments = false, notes = false,
			goban = {scale = 0.5}]{problem}
	\end{codeexample}
\end{key}

	\section{Layout of the diagrams}
\label{subsec:disposicion}

Several boards can go in a grid: the |goban| environments and the diagrams of |\sgfdiagrams| and |\sgfproblem|, inside the |pgfgolayout| environment, or those of a single |\sgfdiagrams| or |\sgfproblem|, with its \emph{key} |layout|. Without a grid, the diagrams of those commands go in the paragraph, side by side, and move to the next line when they do not fit. With |columns|, they go in a grid of columns of the same width: each row is a separate line (the page can break between rows), the boards of a row are aligned at the top and an incomplete row leaves the last columns empty. With |fit|, the boards are also scaled to fill the width of their column.

\begin{environment}{{pgfgolayout}\opt{\oarg{layout options}}}
	Lays out in a grid the |goban| environments it contains, and the diagrams of the |\sgfdiagrams| and |\sgfproblem| it contains, in the order in which they appear. The options are those of |/pgfgo/layout|; the |layout| options of the commands inside are not used. There must be no text inside, only the boards and the commands. With |fit|, each |goban| is drawn several times to measure it: its content must have no effects outside the board (for example, counters that increase).

	\begin{codeexample}[code only]
		\begin{pgfgolayout}[columns = 4, fit]
			\begin{goban}[5]\stones[b]{C3}\end{goban}
			\begin{goban}[5]\stones[w]{C3}\end{goban}
			\sgfdiagrams[step = 5]{game}
		\end{pgfgolayout}
	\end{codeexample}
\end{environment}

\begin{pgfgolayout}[columns = 4, fit]
	\begin{goban}[5]\stones[b]{C3}\end{goban}
	\begin{goban}[5]\stones[w]{C3}\end{goban}
	\sgfdiagrams[step = 5]{game}
\end{pgfgolayout}

\begin{codeexample}[code only]
	\sgfdiagrams[step = 2, to = 6, layout = {columns = 3, fit}]{game}
\end{codeexample}

\sgfdiagrams[step = 2, to = 6, layout = {columns = 3, fit}]{game}

\begin{stylekey}{/pgfgo/layout = \meta{layout options}}
	Options of the layout for the whole document (for example, |columns = 3, fit|). For a single command, |/pgfgo/sgf diagrams/layout| and |/pgfgo/sgf problem/layout|; also for the whole document, |\pgfgoset{sgf diagrams = {layout = {...}}}|.
\end{stylekey}

\begin{key}{/pgfgo/layout/columns = \meta{number or auto} (initially auto)}
	Number of columns of the grid. With |auto|, the diagrams go in the paragraph, without a grid.
\end{key}

\begin{key}{/pgfgo/layout/width = \meta{dimension} or |natural| (initially \texttt{\string\linewidth})}
	Width of the grid. If it is less than that of the line, the grid is centered. Each column measures what is left after subtracting the |column sep| from |width|, divided among the columns; a narrower board sits in its column according to |halign|, so between two boards there is more space than |column sep|.

	With |natural|, the grid fits the diagrams: the columns have the width of the widest diagram (with its caption, which is broken as without a grid, see |caption/min width|) and between two boards of the same width there is exactly |column sep|. The grid measures that times the number of columns (or of diagrams, if there are fewer) and is centered on the line; if it is wider than the line, it goes into the margin. The diagrams are all drawn together at the end (or at the end of each group of |per page| or |per frame|). |fit| does not work with |natural|: the package warns and the boards keep their scale.

	\begin{codeexample}[code only]
		\begin{pgfgolayout}[columns = 2, width = natural, column sep = 0em, row sep = 0em]
			\begin{goban}[partial = A9:E5, scale = 0.5]\stone[b]{C7}\end{goban}
			\begin{goban}[partial = E9:I5, scale = 0.5]\stone[w]{G7}\end{goban}
			\begin{goban}[partial = A1:E5, scale = 0.5]\stone[b]{C3}\end{goban}
			\begin{goban}[partial = E1:I5, scale = 0.5]\stone[w]{G3}\end{goban}
		\end{pgfgolayout}
	\end{codeexample}

	\begin{pgfgolayout}[columns = 2, width = natural, column sep = 0em, row sep = 0em]
		\begin{goban}[partial = A9:E5, scale = 0.5]\stone[b]{C7}\end{goban}
		\begin{goban}[partial = E9:I5, scale = 0.5]\stone[w]{G7}\end{goban}
		\begin{goban}[partial = A1:E5, scale = 0.5]\stone[b]{C3}\end{goban}
		\begin{goban}[partial = E1:I5, scale = 0.5]\stone[w]{G3}\end{goban}
	\end{pgfgolayout}
\end{key}

\begin{key}{/pgfgo/layout/column sep = \meta{dimension} (initially 1em)}
	Space between columns.
\end{key}

\begin{key}{/pgfgo/layout/row sep = \meta{skip} (initially \texttt{\string\medskipamount})}
	Space between rows.
\end{key}

\begin{key}{/pgfgo/layout/fit = \meta{boolean} (default true, initially false)}
	Scales each board so that it fills the width of its column; it replaces the |scale| of the board. The lines do not change width. In |\sgfdiagrams| and |\sgfproblem|, every board uses the scale that fills the column with the first one, so the stones have the same size in all of them. If with that scale a board does not fit in the column (a problem wider than the first one, for example), only that board is reduced until it fills it. Only with |columns|, and not with |width = natural|.
\end{key}

\begin{key}{/pgfgo/layout/per page = \meta{number} or |false| (initially 0)}
	Every \meta{number} diagrams, starts a new page (with |\newpage|): with |per page = 1|, one diagram per page. In the grid, a half row is drawn before the break; \meta{number} is best a multiple of |columns|. With |0| or |false|, there are no breaks. In |\sgfproblem| with |statement = alone|, the statement counts as a diagram.

	\begin{codeexample}[code only]
		\sgfdiagrams[step = 50, layout = {columns = 2, fit, per page = 4}]{lee}
	\end{codeexample}
\end{key}

\begin{key}{/pgfgo/layout/per frame = \meta{number} or |false| (default 1, initially 0)}
	Inside a |beamer| |frame|, every \meta{number} diagrams move to the next slide: with |per frame| alone, one diagram per slide, each with its caption. The diagrams use |\only<+>|, so they add up to the other overlays of the |frame|: after a |\pause|, the first diagram goes on the next slide. Outside |beamer|, the package warns and the \emph{key} is not used. In |beamer|, if there is also |per page|, |per frame| wins.

	\begin{codeexample}[code only]
		\begin{frame}{The game}
			\sgfdiagrams[step = 50, goban = {scale = 0.6}, layout = {per frame}]{lee}
		\end{frame}
	\end{codeexample}
\end{key}

\begin{key}{/pgfgo/layout/halign = \meta{center, left or right} (initially center)}
	Where each diagram goes inside its column, with its caption and its notes. The comments are always aligned to the left.

	\begin{codeexample}[code only]
		\sgfdiagrams[step = 2, to = 6, layout = {columns = 2, width = 0.6\linewidth,
			column sep = 3em, row sep = 2em, halign = left}]{game}
	\end{codeexample}

	\sgfdiagrams[step = 2, to = 6, layout = {columns = 2, width = 0.6\linewidth,
		column sep = 3em, row sep = 2em, halign = left}]{game}
\end{key}

\begin{key}{/pgfgo/layout/valign = \meta{top or bottom} (initially top)}
	How the diagrams of a row of the grid are aligned when their boards have different heights. With |top|, at the top. With |bottom|, the boards end at the same height, and their captions and notes all start on the same line, below. With |columns = auto|, the diagrams go on the baseline of the paragraph, as with |bottom|.

	\begin{codeexample}[]
		\begin{pgfgolayout}[columns = 3, valign = bottom]
			\begin{goban}[5, scale = 0.6]\stones[b]{C3}\end{goban}
			\begin{goban}[5x3, scale = 0.6]\stones[w]{C2}\end{goban}
			\begin{goban}[9, scale = 0.6]\stones[b]{E5}\end{goban}
		\end{pgfgolayout}
	\end{codeexample}
\end{key}

\begin{key}{/pgfgo/layout/align captions = \meta{boolean} (default true, initially false)}
	With |valign = top|, the captions, notes and comments of the diagrams of a row start on the same line, below the tallest board; the boards stay aligned at the top. With |valign = bottom| they are already on the same line. The |goban| environments of |pgfgolayout| have no caption: they do not change.

	\begin{codeexample}[code only]
		\sgfproblem[game = 1:3, solutions = false, partial = auto,
			layout = {columns = 3, align captions}]{problems}
	\end{codeexample}
\end{key}

	\section{Players and profiles}
\label{sec:jugadores}

Each stone, move and mark belongs to one of three players: the first (|black|), the second (|white|) or the neutral one (|neutral|). Their look is not in each command but in the player: the \emph{keys} tree of each one gathers the options of its stones, its moves, its marks and its territory, which in the previous sections appear under |stone|, |move| and |mark|. Here are the players and, afterwards, the profiles.

\subsection{Players}

The options of the board are changed with |\pgfgoset{goban = {...}}|. Those of the players, with these \emph{keys}:

\begin{key}{/pgfgo/players = \meta{players options}}
	Options of the players: |neutral|, |black| and |white|. Players are not to be confused with profiles: a profile is loaded into the players, and then these \emph{keys} change them.
\end{key}

The options of the players can be set all together with the \emph{key} |all|, or separately. The changes made here do not affect the definition of the profiles. That is, if we load the profile |pgfgo| and change the black player, this does not affect the definition of |pgfgo-black|, but the options of the first player itself.

\begin{key}{/pgfgo/players/all = \meta{all player options}}
	Applies the options to the three players.
\end{key}

\begin{key}{/pgfgo/players/neutral = \meta{neutral player options}}
	Options of the neutral player (the stones without a player).
\end{key}

\begin{key}{/pgfgo/players/black = \meta{first player options}}
	Options of the first player.
\end{key}

\begin{key}{/pgfgo/players/white = \meta{second player options}}
	Options of the second player.

	\begin{codeexample}[width=4cm]
		\pgfgoset{players = {
			all = {stone = {line = {color = red}}},
			neutral = {stone = {fill = {color = yellow}}},
			black = {stone = {fill = {color = blue!60!black}}}}}

		\begin{goban}[3]
			\stones[b]{A1, B2}
			\stones[w]{C3}
			\stone{A3}
		\end{goban}
	\end{codeexample}
\end{key}

To avoid writing the whole path, there are four commands:

\begin{command}{\pgfgosetplayer\marg{player}\marg{player options}}
	It is the same as |\pgfgoset{players = {<player> = {<player options>}}}|.
\end{command}

\begin{command}{\pgfgosetallplayers\marg{player options}}
	Applies the options to the three players. It is the same as |\pgfgoset{players = {all = {...}}}|.

	\begin{codeexample}[width=4cm]
		\pgfgosetallplayers{stone = {radius = 0.3cm}}
		\pgfgosetplayer{white}{stone = {fill = {color = orange}}}

		\begin{goban}[3]
			\stones[b]{A1, B2}
			\stones[w]{C3}
			\stone{A3}
		\end{goban}
	\end{codeexample}
\end{command}

\begin{command}{\pgfgocopyplayer\marg{target player}\marg{source player}}
	Copies the current options of |source player| into |target player|.

	\begin{codeexample}[width=3cm]
		\pgfgocopyplayer{neutral}{white}

		\begin{goban}[partial = A1:C3]
			\stone{B2} % Neutral, but with the look of White.
			\stones[b]{A1,C3}
		\end{goban}
	\end{codeexample}
\end{command}

\begin{command}{\pgfgocopyplayers\marg{neutral source}\marg{black source}\marg{white source}}
	Copies, in this order, |neutral source| into the neutral player, |black source| into the black one and |white source| into the white one (with |\pgfgocopyplayer|). Since the copies are made one after the other, it cannot swap players.

	\begin{codeexample}[width=4cm]
		\pgfgocopyplayers{neutral}{white}{black}
		% It does not swap: both are white.

		\begin{goban}[3]
			\stones[b]{A1, B2}
			\stones[w]{C3}
		\end{goban}
	\end{codeexample}
\end{command}

\label{subsec:jugadores}Each player has the same \emph{keys} tree. Each branch is documented with the object it applies to, without the player in the path: for the white player, |/pgfgo/stone/radius| is |/pgfgo/players/white/stone/radius|. The three branches below, although they mention |white|, hold for the three players.

\begin{key}{/pgfgo/players/white/stone = \meta{stone options}}
	Options of its stones and of its prisoners (p.~\pageref{sec:piedras}).
\end{key}

\begin{key}{/pgfgo/players/white/move = \meta{move options}}
	Options of its moves (p.~\pageref{subsec:movimientos}).
\end{key}

\begin{key}{/pgfgo/players/white/mark = \meta{mark options}}
	The style of its marks (p.~\pageref{sec:marcas}). The shape (|triangle|, |custom|, \dots) is not part of the style: it goes in the options of each mark (as in |\marks[w, custom = {?}]{A3}|) or in |/pgfgo/mark|. Inside |\pgfgoset{players = ...}| it gives an error.
\end{key}

The fourth branch, |territory|, is the style of its territory (p.~\pageref{subsec:territorio}).


\subsection{Profiles}
\label{subsec:perfiles}

Profiles are equivalent to \tikzname\ styles. In short, they save a configuration of board and stones with a name and call it whenever wanted.

\begin{command}{\pgfgodefinegoban\opt{\oarg{options}}\marg{goban name}\marg{goban options}}
	Gives a name to a group of |goban| options. It has two options.

	\begin{key}{/pgfgo/profiles/goban/from = \meta{goban name}}
		Defines the goban from a copy of another one.
		\begin{codeexample}[code only]
			\pgfgodefinegoban[from = gobanA]{gobanB}{...}
		\end{codeexample}

		A goban can copy itself to add options, although it is not the best alternative.
		\begin{codeexample}[code only]
			\pgfgodefinegoban[from = gobanB]{gobanB}{...}
		\end{codeexample}
	\end{key}

	\begin{key}{/pgfgo/profiles/goban/append}
		Adds options to a goban already defined.
		\begin{codeexample}[code only]
			\pgfgodefinegoban[append]{gobanB}{...}
		\end{codeexample}
		It is equivalent to:
		\begin{codeexample}[code only]
			\pgfgodefinegoban[from = gobanB]{gobanB}{...}
		\end{codeexample}
	\end{key}
\end{command}

\begin{command}{\pgfgodefineplayer\opt{\oarg{options}}\marg{player name}\marg{player options}}
	The same, for a player. Its options are those of |\pgfgodefinegoban|.

	\begin{key}{/pgfgo/profiles/player/from = \meta{player name}}
		Defines the player from a copy of another one.
	\end{key}

	\begin{key}{/pgfgo/profiles/player/append}
		Extends the options of a player already defined. It is the same as |from = |\meta{player name} with its own name.
	\end{key}
\end{command}

\begin{key}{/pgfgo/profile = \meta{name}}
	Applies the profile |name|: the goban |name| and the players |name-neutral|, |name-black| and |name-white|. It is used in |\pgfgoset|, in the options of a |goban| and in those of a view (p.~\pageref{subsec:vistas}).

	Before applying it, the board and the players go back to the initial style (that of the profile |default|): so a profile does not inherit what it does not define, and a shadow or an image set before is lost. What is not style stays as it was: the size, |partial|, |scale|, where the coordinates go and which they are, |hoshi/at|, the captures and the notes. What comes after the profile is added to it.

	\begin{codeexample}[width=3cm]
		\pgfgoset{stone = {shadow}}
		\begin{goban}[3, profile = default]
			\stones[b]{A1, B2}
			\stones[w]{C3}
		\end{goban}
	\end{codeexample}

	The goban |name| and the players |name-neutral|, |name-black| and |name-white| must exist. For example, |profile = pgfgo| works because there is a goban called |pgfgo|, and players called |pgfgo-neutral|, |pgfgo-black| and |pgfgo-white|.
\end{key}

\begin{key}{/pgfgo/goban profile = \meta{goban name}}
	Applies only the goban |goban name|, after taking the board back to the initial style.
\end{key}

\begin{key}{/pgfgo/players profile = \meta{name}}
	Applies only the players |name-neutral|, |name-black| and |name-white|, each one after going back to the initial style.

	\begin{codeexample}[]
		\begin{goban}[3, goban profile = pgfgo]
			\stones[b]{A1, B2}
			\stones[w]{C3}
		\end{goban}

		\begin{goban}[3, players profile = pgfgo]
			\stones[b]{A1, B2}
			\stones[w]{C3}
		\end{goban}
	\end{codeexample}
\end{key}

\begin{key}{/pgfgo/players/black/profile = \meta{player name}}
	Applies the player |player name| to one of the players (|neutral|, |black| or |white|), after taking it back to the initial style. With |all|, the same one to the three.
	\begin{codeexample}[code only]
		\pgfgoset{goban profile = pgfgo, players = {black = {profile = wain-black}}}
	\end{codeexample}
\end{key}

\begin{codeexample}[width=3cm]
	\pgfgodefinegoban{example}{background = {fill = {color = blue!05}}}
	\pgfgodefineplayer[from = default-neutral]{example-neutral}
	{stone = {
			fill = {color = gray},
			line = {color = magenta!20!black}}
	}

	\pgfgodefineplayer[from = example-neutral]{example-white}
	{stone = {fill = {color = yellow}}}

	\pgfgodefineplayer[from = example-neutral]{example-black}
	{stone = {fill = {color = blue}}}

	\pgfgoset{profile = example}

	\begin{goban}[partial = A1:C3]
		\moves[numbers = false]{A1:3}
		\stones[b]{B1,B3,C2}
		\stones[w]{C1,C3}
		\stone{B2}
	\end{goban}
\end{codeexample}

\begin{command}{\pgfgodefineprofile\opt{\oarg{options}}\marg{name}\marg{profile options}}
	Defines at once a whole profile for |profile|: the goban |name| and the players |name-neutral|, |name-black| and |name-white|. It is the same as the example above, in a single command.

	\begin{codeexample}[width=3cm]
		\pgfgodefineprofile[from = default]{example2}{
			goban = {background = {fill = {color = blue!05}}},
			all = {stone = {line = {color = magenta!20!black}}},
			white = {stone = {fill = {color = yellow}}},
			black = {stone = {fill = {color = blue}}}}

		\pgfgoset{profile = example2}

		\begin{goban}[partial = A1:C3]
			\moves[numbers = false]{A1:3}
			\stones[b]{B1,B3,C2}
			\stones[w]{C1,C3}
			\stone{B2}
		\end{goban}
	\end{codeexample}

	\begin{key}{/pgfgo/profiles/profile/goban = \meta{goban options}}
		The options of the goban.
	\end{key}

	\begin{key}{/pgfgo/profiles/profile/all = \meta{player options}}
		Options of the three players. They go before those of each one.
	\end{key}

	\begin{key}{/pgfgo/profiles/profile/neutral = \meta{player options}}
		Options of the neutral player.
	\end{key}

	\begin{key}{/pgfgo/profiles/profile/black = \meta{player options}}
		Options of the black player.
	\end{key}

	\begin{key}{/pgfgo/profiles/profile/white = \meta{player options}}
		Options of the white player.
	\end{key}

	\begin{key}{/pgfgo/profiles/profile/from = \meta{profile name}}
		Option of the command: each part starts with that of the profile \meta{profile name} (the goban \meta{profile name}, the player \meta{profile name}|-neutral|, etc.).
	\end{key}
\end{command}

	\section{Catalog}
\label{sec:catalog}

\subsection{Profiles}
\subsubsection{default}
\begin{codeexample}[width=4cm]
	\pgfgouseprofile{default}
	
	\begin{goban}[partial = {from = A1 to G3}]
		\moves*[from = 1]{A:F 3}
		\stones[w]{A1:2}
		\stones[b]{B1:2}
		\stones[n]{C1:2}
		\stone[w, mark = circle]{D1}
		\mark[w, circle]{D2}
		\stone[b, mark = circle]{E1}
		\mark[b, circle]{E2}
		\stone[n, mark = circle]{F1}
		\mark[n, circle]{F2}
		\mark[w, tsquare]{G1}
		\mark[b, tcircle]{G2}
		\mark[n, tsquare]{G3}
	\end{goban}
\end{codeexample}

\subsubsection{pgfgo}
\begin{codeexample}[width=4cm]
	\pgfgouseprofile{pgfgo}
	
	\begin{goban}[partial = {from = A1 to G3}]
		\moves*[from = 1]{A:F 3}
		\stones[w]{A1:2}
		\stones[b]{B1:2}
		\stones[n]{C1:2}
		\stone[w, mark = circle]{D1}
		\mark[w, circle]{D2}
		\stone[b, mark = circle]{E1}
		\mark[b, circle]{E2}
		\stone[n, mark = circle]{F1}
		\mark[n, circle]{F2}
		\mark[w, tsquare]{G1}
		\mark[b, tcircle]{G2}
		\mark[n, tsquare]{G3}
	\end{goban}
\end{codeexample}

\subsubsection{pgfgosmooth}
\begin{codeexample}[width=4cm]
	\pgfgouseprofile{pgfgosmooth}
	
	\begin{goban}[partial = {from = A1 to G3}]
		\moves*[from = 1]{A:F 3}
		\stones[w]{A1:2}
		\stones[b]{B1:2}
		\stones[n]{C1:2}
		\stone[w, mark = circle]{D1}
		\mark[w, circle]{D2}
		\stone[b, mark = circle]{E1}
		\mark[b, circle]{E2}
		\stone[n, mark = circle]{F1}
		\mark[n, circle]{F2}
		\mark[w, tsquare]{G1}
		\mark[b, tcircle]{G2}
		\mark[n, tsquare]{G3}
	\end{goban}
\end{codeexample}

\subsubsection{wain}
\begin{codeexample}[width=4cm]
	\pgfgouseprofile{wain}
	
	\begin{goban}[partial = {from = A1 to G3}]
		\moves*[from = 1]{A:F 3}
		\stones[w]{A1:2}
		\stones[b]{B1:2}
		\stones[n]{C1:2}
		\stone[w, mark = circle]{D1}
		\mark[w, circle]{D2}
		\stone[b, mark = circle]{E1}
		\mark[b, circle]{E2}
		\stone[n, mark = circle]{F1}
		\mark[n, circle]{F2}
		\mark[w, tsquare]{G1}
		\mark[b, tcircle]{G2}
		\mark[n, tsquare]{G3}
	\end{goban}
\end{codeexample}



	\appendix
	\section{What's new}
\label{sec:novedades}


\begin{itemize}
	\item Next version:
		\begin{itemize}
			\item \tikzname\ patterns in the fills of the stones, the marks, the territory and the background (|fill/pattern|, p.~\pageref{sec:piedras}).
			\item Coordinates on the border in lower case (|label/letter case = lower|, p.~\pageref{subsec:etiquetas}).
			\item The profile |draft|: black and white, with thin lines, no border on the background, and lines where there is a fill.
			\item |goban/arrow| and |goban/line|: the style of the arrows and lines is part of the profile of the board (p.~\pageref{subsec:flechas}). Every profile defines it: |pgfgo| and |pgfgosmooth| draw them in |red!80!black|.
			\item |scale line width|: |scale| also multiplies the width of the lines.
			\item |layout/valign = bottom|: the boards of a row of the grid end at the same height (p.~\pageref{subsec:disposicion}).
			\item |partial| of |\sgfdiagrams|, |\sgfproblem| and |\sgfanimation| also takes two corners (|partial = K10:T19|): the same part in every board (p.~\pageref{subsec:sgfdiagramas}).
			\item |\pgfgodefinesgf| takes |origin| and |transpose|: the corner where the point |aa| of the file goes and whether the first letter is the row, the eight symmetries of the board (p.~\pageref{subsec:sgfnombres}).
			\item |name| in |\sgfdiagrams| and |\sgfproblem|: each diagram is a named board (|var1|, |var2a|...), for |resume| and the views in the text (p.~\pageref{subsec:sgfdiagramas}).
			\item Columns of two letters after |Z|, for boards of more than 25 columns: |AA5|, |partial = AZ5:BG8| (p.~\pageref{subsec:sintaxis-de-coordenadas}).
			\item |links| in |\sgfproblem|: with |hyperref|, the statement and its solutions link to each other (p.~\pageref{subsec:sgfproblemas}).
			\item The option |pages|: one board per page, and |<document>-pages.lua| turns them into SVG, PNG or PDF (p.~\pageref{subsec:paginas}).
			\item Sensei's Library and L19 diagrams (|$$| lines), in a file or in the environment |pgfgodiagram|: every SGF command reads them (p.~\pageref{subsec:sgfsensei}).
			\item |goban/blank|: intersections that are not part of the board, without lines or star points; they are the |_| of Sensei's Library diagrams (p.~\pageref{subsec:parcial}).
			\item The sizes of the board are in |pt|: they do not change with the size of the font of the document. One character of |mark/sequence| and one of |mark/custom| have the same size.
			\item |grid/edge|: the edge of the grid, thicker or in another color.
			\item |layout/align| is renamed |layout/halign|, and |layout/align captions| puts the captions of a row on the same line with the boards at the top (p.~\pageref{subsec:disposicion}).
			\item |notes/font|: the font of the notes.
			\item |\sgfproblem| leaves out the games that the file does not have, with a single warning.
			\item In Spanish, the colors are singular and the result also agrees with a name: ``Gana Negro'', ``Juega Blanco'', ``Gana Ana por abandono''. In Portuguese, ``Vitória de Ana''.
		\end{itemize}
	\item \pgfgoversion:
		\begin{itemize}
			\item Kernel and rules of the game:
				\begin{itemize}
					\item The computations (coordinates, numbering of moves, memory, etc.) have been separated from the drawing. They are done by a \emph{kernel}, chosen with the package option |kernel| (p.~\pageref{subsec:kernel}). Besides the kernel written in L3, a kernel written in Lua has been added, which is the default with Lua\LaTeX.
					\item With the Lua kernel, captures have been added (|captures|, p.~\pageref{subsec:capturas}). With |captured marks = remove|, the marks on a captured stone go with it, and |\pgfgoprisoners| counts the prisoners of each player.
					\item ``N at M'' notes have been added for numbered moves on points already played (|notes| and |\pgfgomovenotes|, p.~\pageref{subsec:notas}). The numbered passes of an SGF are noted as ``N: pass'' (|notes/pass|).
					\item |\territory|, which marks the territory (with a polygon, or by itself with |auto| and the Lua kernel), and |\pgfgoscore|, which counts the points (p.~\pageref{subsec:territorio}), have been added. With |draw = region|, the territory of each color is a single polygon, and |\pgfgoregion| gives its measures.
				\end{itemize}
			\item Named boards and views:
				\begin{itemize}
					\item A |goban| with |name| keeps everything it draws, and another one starts with that board with |resume| (p.~\pageref{subsec:recordar}). |resume| brings the stones; |numbers|, |marks| and |territory| bring the rest. With |name| and |resume| in the same |goban|, variations are built.
					\item |\pgfgoview| draws a point or part of a named board in the middle of the text (p.~\pageref{subsec:vistas}), with the options of |view| (|marks|, |arrows|, |lines|, |numbers|, |territory|, |forgotten|, |raise|, |profile|).
					\item A view can come before its |goban|, like |\ref| (p.~\pageref{subsec:vistas-antes}), and a board can be used in another document with |export| and |\pgfgoimport| (p.~\pageref{subsec:exportar}).
				\end{itemize}
			\item SGF:
				\begin{itemize}
					\item With the Lua kernel, |\sgf| reads games in SGF (p.~\pageref{subsec:sgftablero}): variations (|variation|), marks of the file (|marks|), the numbering of |MN| and |restart|, and the size of the board taken from the file (|sgf size|).
					\item Collections of games (|game|), the view (|VW|, |sgf view|), figures (|FG|) in |\sgfdiagrams|, variations as letters (|variations|), |PL| and files in Latin-1.
					\item The numbering asked for by |PM| (|numbers| and |restart = auto|), and the view (|VW|) when it changes during the game (|sgf view move|; |\sgfdiagrams| and |\sgfanimation| use it by themselves). |step| accepts a list, and |restart = {from = |\meta{n}|}| numbers each stretch from \meta{n}.
					\item |variation = {R16, P17, pass}|: the variation can also be given with its moves, from the first one.
					\item Files in GBK, Big5, Shift\_JIS, EUC-JP and EUC-KR are read (|CA| or |sgf charset|). If the font lacks the Chinese, Japanese or Korean characters of a text of the file, the package warns.
					\item Names for the files, with |\pgfgodefinesgf| (p.~\pageref{subsec:sgfnombres}): the path, the game and the charset are written only once.
					\item The information of the game (|\sgfinfo|, also of a move with |move|). The names of the nodes and the annotations (|TE|, |BM|, |GB|, etc.) are written with the comments. |\sgfcomment| writes the comment of a move, and |chat = false| removes the chat of the servers (p.~\pageref{subsec:sgfdatos}).
					\item |\sgfdiagrams| draws a game as a series of diagrams, as in books (p.~\pageref{subsec:sgfdiagramas}). It writes the comments below each diagram, beside it or after all of them (|comments|), and the result, in words, and the score below the last one (|result|). With |variations|, it also draws the variations, each in its own diagram.
					\item |last mark| marks the last move, and |to = last| goes to the end.
					\item |\sgfproblem| has been added, which draws a problem (tsumego): the statement and one diagram for each variation of the solution, with its comments (p.~\pageref{subsec:sgfproblemas}). With |game = 1:20|, a problem sheet from a collection; with |verdict|, the caption says whether the variation is right.
					\item With the |animate| package, |\sgfanimation| steps through a game with buttons (p.~\pageref{subsec:sgfanimacion}). It builds the animation in layers (|layers|): each stone is drawn only once, and the PDF of a long game takes much less space.
				\end{itemize}
			\item Drawing:
				\begin{itemize}
					\item Arrows and lines have been added (|\arrow| and |\line|, p.~\pageref{subsec:flechas}), with |dash|, |shorten|, |bend| and |tips|, and the mark |dim|. |\sgf| also draws |AR|, |LN|, |SL| and |DD|.
					\item The board draws the star points (|hoshi|, p.~\pageref{subsec:hoshi}).
					\item The coordinates of the border can be letters, numbers or kanji on each axis, counted from either side (|columns|, |rows|, |style = japanese|, etc.).
					\item |partial = auto| shows only the part of the board used by what is drawn (p.~\pageref{subsec:parcial}); also in |\sgfdiagrams|, |\sgfanimation| and |\sgfproblem|.
					\item Images (PNG, JPG or PDF) on the stones (|stone/image|, p.~\pageref{subsec:piedrasimagen}), on the background and the grid of the board, on the star points and on the marks. With several files, each point always uses the same one, and |rotate| turns each stone or tile. Each image is kept only once in the PDF (p.~\pageref{subsubsec:imagenespdf}).
					\item Shadows and shines of your own (|shadow/image| and |shine/image|), and |shadow/scale| and |shadow/color|.
					\item Marks of your own, with PGF code or with an image (|\pgfgodefinemark|, p.~\pageref{subsec:marcaspropias}), and |mark/scale|.
					\item Numbers in kanji for the moves and the marks (|format = kanji|).
					\item Opacity of the stroke and of the fill of the marks, of the numbers of the moves, of the images of the stones and of the star points.
					\item The \tikzname\ coordinate system |go|: |(go cs:C3)| (p.~\pageref{subsec:tikz}).
				\end{itemize}
			\item Styles and profiles:
				\begin{itemize}
					\item Profiles are \emph{keys} (p.~\pageref{subsec:perfiles}): |profile|, |goban profile|, |players profile| and |players/|\meta{player}|/profile|. Before they are applied, the board and the players go back to the initial style. |\pgfgodefineprofile| defines a whole profile at once.
					\item Three new profiles: |simple|, and |shell| and |yunzi|, with images (p.~\pageref{subsec:perfilshell}): a wooden board with slate and shell stones, or matte ones. The images come with the package, next to its |.sty| files, and are used by their name.
					\item The same rule for the \emph{keys} that turn something on or off (p.~\pageref{subsec:prender}).
					\item The options of a |goban| accept those of |\pgfgoset|, only for that board, and |\pgfgoset{mark = {...}}| also sets the style of the marks of the three players.
				\end{itemize}
			\item Layout: |layout| and the |pgfgolayout| environment lay out the diagrams in a grid of columns, with |fit|, |width = natural|, |per page| and |per frame| (p.~\pageref{subsec:disposicion}).
			\item Languages: the texts of the package are translated into ten languages and follow |babel| or |polyglossia|; |lang|, |\pgfgotext|, |\pgfgodefinelanguage| and |\pgfgolanguagealias| (p.~\pageref{subsec:idiomas}).
			\item Coordinates:
				\begin{itemize}
					\item Lines in any direction and stairs: |A1 -- E6| (in any direction), \verb+A1 -|- E6+ and \verb+A1 |-| E6+.
					\item Rectangles are written |A1:C3|, as in |partial|, and stretches and rectangles are chained with the operators.
					\item |\moves| accepts every form of |\stones| and plays each one in its order (p.~\pageref{subsec:orden-coordenadas}); |\moves| and |\marks[sequence]| warn if they have no fixed order.
				\end{itemize}
			\item Behavior changes (the full list, with the old names of the \emph{keys}, is in |CHANGELOG.md|):
				\begin{itemize}
					\item |\moves| numbers its moves; |numbers = false| leaves them without a number. |\move|, |\move*| and |\moves*| have been removed.
					\item |\stones| without a player draws neutral stones, like |\stone|.
					\item The only way to remember a board is |name| and |resume|; |goban/remember|, |goban/forget| and |resume| without a name have been removed.
					\item |\pgfgouseprofile| and the other |\pgfgouse...| have been removed; now |\pgfgoset{profile = wain}| is used.
					\item The operator |R| has been removed (|A1 R C3| is |A1:C3|), and |:--:|, \verb+:-|:+ and \verb+:|-:+ are now |--|, \verb+-|-+ and \verb+|-|+.
					\item The territory is drawn when the |goban| ends, with the last |\territory| that marked each point.
					\item A move number is never covered by a label of the file (|LB|) or by |last mark|.
				\end{itemize}
		\end{itemize}
	\item 0.2025.10.13: Bug fixes. Each |goban| starts at move 1 with Black, |\mark| and |\marks| only exist inside |goban|, and what falls outside a partial board is no longer drawn. The full list is in the file |CHANGELOG.md|.
	\item 0.2025.03.10:
		\begin{itemize}
			\item The \emph{keys} |inner| and |outer| have been added for the \emph{rounded corner} of the board.
			\item Three new coordinate operators have been added: |*| for every intersection, |L| for the whole column of coordinate L and |N| for the whole row of coordinate N.
		\end{itemize}
	\item 0.2025.03.07: First release.
\end{itemize}

\end{document}
